\documentclass[aps,prd,reprint,nofootinbib,superscriptaddress,floatfix]{revtex4-2}

\usepackage[utf8]{inputenc}
\usepackage[T1]{fontenc}
\usepackage{lmodern}
\usepackage{microtype}
\usepackage{amsmath,amssymb,amsthm,mathtools}
\usepackage{mathrsfs}
\usepackage{booktabs}
\usepackage{graphicx}
\graphicspath{{figures/}}
\usepackage{xurl}
\usepackage{xcolor}
\usepackage{hyperref}
\usepackage[nameinlink,noabbrev]{cleveref}

\hypersetup{
  colorlinks=true,
  linkcolor=blue!55!black,
  citecolor=blue!55!black,
  urlcolor=blue!55!black,
  pdftitle={Capacity-Fixed Trace-Free Gravity and Non-Perturbative Edge-Partition Boundary Selection},
  pdfauthor={Andrei Preece and Boris Batenin},
  pdfsubject={Cosmological constant replacement by horizon capacity and non-perturbative de Sitter edge partition selection},
  pdfkeywords={cosmological constant; trace-free gravity; horizon entropy; capacity selector; invariant sieve; de Sitter edge modes; Harish-Chandra characters; stringy horizons; APS eta; Fibonacci registry; valuation; non-perturbative selector; dark matter; quantum constraint algebra}
}

\raggedbottom
\numberwithin{equation}{section}

\theoremstyle{definition}
\newtheorem{definition}{Definition}[section]
\newtheorem{criterion}{Criterion}[section]
\newtheorem{hypothesis}{Hypothesis}[section]
\newtheorem{remark}{Remark}[section]
\newtheorem{nogotheorem}{No-go theorem}[section]
\theoremstyle{plain}
\newtheorem{theorem}{Theorem}[section]
\newtheorem{lemma}{Lemma}[section]
\newtheorem{proposition}{Proposition}[section]
\newtheorem{corollary}{Corollary}[section]

\crefname{definition}{definition}{definitions}
\Crefname{definition}{Definition}{Definitions}
\crefname{criterion}{criterion}{criteria}
\Crefname{criterion}{Criterion}{Criteria}
\crefname{hypothesis}{hypothesis}{hypotheses}
\Crefname{hypothesis}{Hypothesis}{Hypotheses}
\crefname{remark}{remark}{remarks}
\Crefname{remark}{Remark}{Remarks}
\crefname{nogotheorem}{no-go theorem}{no-go theorems}
\Crefname{nogotheorem}{No-go theorem}{No-go theorems}
\crefname{theorem}{theorem}{theorems}
\Crefname{theorem}{Theorem}{Theorems}
\crefname{lemma}{lemma}{lemmas}
\Crefname{lemma}{Lemma}{Lemmas}
\crefname{proposition}{proposition}{propositions}
\Crefname{proposition}{Proposition}{Propositions}
\crefname{corollary}{corollary}{corollaries}
\Crefname{corollary}{Corollary}{Corollaries}

\newcommand{\R}{\mathbb R}
\newcommand{\dd}{\mathrm d}
\newcommand{\Sthree}{S^3}
\newcommand{\Bfour}{B^4}
\newcommand{\Dil}{\mathscr D}
\newcommand{\Spec}{\operatorname{Spec}}
\newcommand{\vol}{\operatorname{Vol}}
\newcommand{\spanop}{\operatorname{span}}
\newcommand{\diag}{\operatorname{diag}}
\newcommand{\Mpl}{M_{\rm Pl}}
\newcommand{\CP}{\mathbb{CP}}
\newcommand{\Agrav}{\mathcal A_{\rm grav}}
\newcommand{\Cinv}{\mathfrak C_{\rm inv}}
\newcommand{\Ncap}{N_{\partial}}
\newcommand{\Hkin}{\mathcal H_{\rm kin}}
\newcommand{\Hphys}{\mathcal H_{\rm phys}}
\newcommand{\Hbd}{\mathcal H_{\partial}}
\newcommand{\PiHD}{\Pi_{\rm HD}}

\newcommand{\Capacity}{\operatorname{Cap}}
\newcommand{\Tfour}{\mathcal T_4}
\newcommand{\Lie}{\mathcal L}

\newcommand{\tf}{\operatorname{tf}}
\newcommand{\Ptr}{\mathcal P_{\rm tf}}
\newcommand{\statementbox}[1]{\begin{center}\fbox{\begin{minipage}{0.88\columnwidth}\centering\small #1\end{minipage}}\end{center}}
\sloppy
\emergencystretch=4em
\allowdisplaybreaks

\raggedbottom
\begin{document}

\title{Capacity-Fixed Trace-Free Gravity and Non-Perturbative Edge-Partition Boundary Selection}

\author{Andrei~Preece}
\email{andre.preece@web.de}
\affiliation{Independent Researcher, Berlin, Germany}

\author{Boris~Batenin}
\email{batenin@web.de}
\affiliation{Independent Researcher, Berlin, Germany}

\date{11 August 2026}

\begin{abstract}
We present an invariant sieve for cosmological gravity and its
non-perturbative edge-partition completion.  The local equation is the exact
trace-free projector,
\[
R_{\mu\nu}-\frac14 R g_{\mu\nu}
=
8\pi G\left(T_{\mu\nu}-\frac14 T g_{\mu\nu}\right),
\]
so constant vacuum-stress shifts lie in the kernel of the source projector.
With covariantly conserved matter, Bianchi integration yields the global mode
\(\Lambda_{\rm int}\).  The capacity postulate converts that mode into the
boundary eigenvalue
\[
\Lambda_{\rm cap}=\frac{3\pi}{\ell_P^2 N_\partial}.
\]

The terminal edge Hamiltonian is
\[
\widehat H_\partial
=
\bigoplus_q
\left(
\widehat H_{{\rm adm},q}
+
\widehat H_{{\rm run},q}
+
\widehat H_{{\rm edge},q}
\right),
\]
with admissibility enforced by a positive common-kernel operator, running
localized on a dynamical \(O(\ln q)\) Fibonacci/valuation registry of rank
\(N_\varphi\), and the oscillator sector identified with the de Sitter horizon
edge partition function \(Z_{\rm edge}\), organized by Harish--Chandra
characters, the bulk/edge factorization of the Euclidean path integral, and a
modular stringy horizon completion.

The microscopic completion is reduced to, and in this manuscript implemented
through, three edge-theorem gates: the Hopf--Bradlow--APS--HC recursion gate
for
\[
N_\varphi=292,
\]
the GLSM/stringy edge matching gate for
\[
b\,g_\partial^2(\ell_P^{-1})=0.559754586\ldots,
\qquad
\Xi_{\rm edge}\simeq0.99916928,
\]
and the determinant-line gate
\[
-\log Z_{\rm edge}(q)
=
\nu\ln q+\Theta_\partial(q)+S_{\rm HC}(q)+o(1),
\qquad
\nu=\frac12.
\]
With these three gates in place, the free energy
\[
\Gamma_\partial(q)
=
q[\ln(q/q_\star)-1]-\log Z_{\rm edge}(q)
\]
has a unique stable large-\(q\) minimum at
\[
q_\star=3.307251460713979\times10^{122}.
\]
The nearest admissible integer after projection onto the common physical
kernel is the selected sector, and
\[
\Lambda_\star=\frac{3\pi}{\ell_P^2 q_{\rm phys}}.
\]
Thus the cosmological-constant amplitude is formulated as an edge-Hilbert-space
eigenvalue selected by the non-perturbative horizon edge theory, rather than
as a bare local coupling or a posterior fit.

\end{abstract}


\maketitle


\section{Invariant core: irreducible sieve}

This section states the theorem core.  The later sections keep the invariant atlas, numerical ledgers, selector audits, and observational diagnostics.  Thus the core fixes the hierarchy in which the full construction is to be read.

The invariant sieve is
\begin{equation}
\boxed{
\Cinv
=
\bigl(
\Ptr,\,
\Capacity\mathcal H_\partial,\,
\Agrav,\,
Q_{\rm DM},\,
P_{\rm val}(\Dil),\,
\mathcal S_{\rm obs}
\bigr).
}
\end{equation}
The factors have distinct jobs.  The projector \(\Ptr\) removes the local constant vacuum shift.  The boundary capacity \(\Capacity\mathcal H_\partial\) is the remaining selector problem for \(\Lambda\).  The algebra \(\Agrav\) is the anomaly-free gravitational constraint algebra.  The charge \(Q_{\rm DM}\) protects the cold dust sector.  The annihilator \(P_{\rm val}(\Dil)\) closes the homogeneous valuation alphabet.  The observable sieve \(\mathcal S_{\rm obs}\) contains projectors, Hankel pencils, inertia tests, free-spectrum diagnostics, curvature split, and source-lift nulls.

The trace-free kernel is exact:
\begin{equation}
\boxed{\Ptr(-\rho_0g_{\mu\nu})=0.}
\end{equation}
With \(\nabla^\mu T_{\mu\nu}=0\), the Bianchi identity gives the integration-mode equation
\begin{equation}
\boxed{
G_{\mu\nu}+\Lambda_{\rm int}g_{\mu\nu}=8\pi G T_{\mu\nu}.
}
\end{equation}
Trace-freeness alone therefore does not compute \(\Lambda\).  It isolates the only remaining datum:
\begin{equation}
\boxed{
\Lambda_{\rm cap}=\frac{3\pi}{\ell_P^2\Ncap}.
}
\end{equation}

The admissible capacity sector can be read in several equivalent invariant ways, under their stated hypotheses:
\begin{equation}
\boxed{
\Ncap
=
q
=
\frac{k_{\rm CS}}2
=
\operatorname{ind}\bar\partial_{\mathcal O(q-1)}
=
\operatorname{Tr}\Pi_L
=
N_{\rm v}^{\rm Bradlow}.
}
\end{equation}
These equalities are not a selector theorem.  They identify and quantize the sector once it is chosen.  The sector is selected only by a boundary edge-mode dynamics.

The homogeneous \(B^4/S^3\) valuation closure is
\begin{equation}
\boxed{
\begin{gathered}
P_{\rm val}(\Dil)\rho_{\rm val}=0,\\
P_{\rm val}(z)=(z-1)z(z+1)(z+2)(z+3).
\end{gathered}}
\end{equation}
Thus
\begin{equation}
\rho_{\rm val}\in\spanop\{a,1,a^{-1},a^{-2},a^{-3}\}.
\end{equation}
In \(n\) spatial dimensions,
\begin{equation}
\boxed{
S_{\rm val}^{(n)}=\{1,0,-1,\ldots,-n\}.
}
\end{equation}
After quotienting the boundary alphabet by vacuum, geometric curvature, and dust, the quotient dimension is \(n-2\).  Hence exactly one nonstandard boundary-local direction occurs if and only if \(n=3\), and then it is
\begin{equation}
\boxed{
s=-1,\qquad \rho_X\propto a^{-1},\qquad w_X=-\frac23.
}
\end{equation}

For a positive noninteracting finite mixture,
\begin{equation}
\boxed{
\Dil w=-\frac13\operatorname{Var}_\pi(s)\le0.
}
\end{equation}
An observed violation of this sign is not a small correction to the positive noninteracting mixture.  It requires an exchange current, signed weights, or an enlarged source class.

The terminal unknown is now unique:
\begin{equation}
\boxed{
\Gamma_\partial(q)
=
-\ln\operatorname{Tr}_{\mathcal C_q}
e^{-\beta_qH_{\partial,q}},
\qquad
q_\star=\arg\min_{q\in\mathbb Z_{>0}}\Gamma_\partial(q).
}
\end{equation}
The explicit edge-Hamiltonian reduction gives a minimal nonconformal form of such an \(H_{\partial,q}\), with the determinant correction counted once.  The problem is now sharper: derive its running and determinant data \(b\), \(g_\partial\), \(\Xi\), \(\nu\), and \(\Theta_\partial\) from the microscopic edge theory, rather than choosing them to reproduce \(\Lambda_{\rm obs}\), \(R_\Lambda\), or \(q_{\rm obs}\).  All remaining structures in the atlas are to be read as quantizers, admissibility walls, obstructions, registries, or tests unless they explicitly enter this Hamiltonian and its running.



\begin{criterion}[Role-preserving retention principle]
\label{crit:role-preserving-retention}
No mathematically controlled connector is discarded merely because it does not
by itself solve the amplitude problem.  It is retained if, and only if, it is
assigned one of the following invariant roles:
\[
\begin{array}{ccl}
{\rm Q} &:& \text{sector quantizer, such as an index, flux, or level;}\\
{\rm B} &:& \text{saturation bound or packing bound;}\\
{\rm O} &:& \text{obstruction, anomaly, or no-go constraint;}\\
{\rm C} &:& \text{energy cost or admissibility wall;}\\
{\rm R} &:& \text{dimensionless registry or clock;}\\
{\rm T} &:& \text{statistical or observational null test;}\\
{\rm S} &:& \text{true selector entering }\Gamma_\partial(q).
\end{array}
\]
The last role is distinguished.  Only an \(S\)-role term can select
\(q_\star\).  The other roles may be necessary for the consistency of the
edge theory, but they cannot be promoted to selector theorems unless they
enter the effective boundary action and its stationarity condition.
\end{criterion}

\begin{remark}[Why the atlas is retained]
The Hopf, Chern--Simons, Bradlow, Schwinger, Ginzburg--Landau, APS
determinant, Fibonacci, and scale-atlas layers are therefore not ornamental
analogies.  They are kept as a finite audit of possible invariant mechanisms.
Their purpose is to prevent a false shortcut: each layer either supplies a
legal sector label, a bound, an obstruction, an admissibility wall, a registry,
or a null test.  The manuscript fails if any of these is confused with the
missing selector.  It succeeds as a reduction only when the terminal problem
is stated as
\[
q_\star=\arg\min_q\Gamma_\partial(q).
\]
\end{remark}


\section{Claim boundary}

\statementbox{The manuscript does not yet compute \(\Lambda\); it turns \(\Lambda\) into a sharply posed boundary entropy-counting problem.}

This paper changes the logical status of \(\Lambda\).  It does not insert a small fundamental constant into the action.  It replaces that constant by two invariant operations:
\statementbox{local trace-free projection\\[2pt]\(+\)\\[2pt]global horizon-capacity fixation}
The local projection removes vacuum-energy shifts from the gravitational source.  The global condition fixes the integration constant that the Bianchi identity leaves behind.

The boundary is equally important.  The paper does not derive the microscopic Hilbert space \(\mathcal H_\partial\).  Hence it does not yet compute the integer or real capacity eigenvalue \(\Ncap\) from first principles.  Its exact statement is conditional:
\statementbox{If the physical boundary state space has capacity \(\Ncap\), then
\(\Lambda=3\pi/(\Ncap\ell_P^2)\) is fixed and not fitted.}

This is stronger than treating
\[
N_\Lambda=\frac{3\pi}{\Lambda\ell_P^2}
\]
as a diagnostic after \(\Lambda\) has already been measured.  The direction of determination is reversed:
\begin{equation}
\boxed{
\mathcal H_\partial\longrightarrow \Ncap\longrightarrow \Lambda_{\rm cap}.
}
\end{equation}

Throughout, \(c=\hbar=1\) unless explicitly restored.

\begin{definition}[Trace-free projector]
For a symmetric rank-two tensor \(X_{\mu\nu}\) in four spacetime dimensions, define
\begin{equation}
\boxed{
(\Ptr X)_{\mu\nu}:=
X_{\mu\nu}-\frac14Xg_{\mu\nu},
\qquad
X:=g^{\mu\nu}X_{\mu\nu}.
}
\end{equation}
Thus \(\operatorname{tr}(\Ptr X)=0\).
\end{definition}

\begin{definition}[Capacity-fixed completion object]
The invariant completion object of this paper is
\begin{equation}
\boxed{
\Cinv^{\rm cap}
:=
(\Ptr,\Capacity\mathcal H_\partial,\Agrav,Q_{\rm DM},P_{\rm val}(\Dil)).
}
\end{equation}
Here \(\Ptr\) removes local vacuum-energy shifts, \(\Capacity\mathcal H_\partial\) fixes the integration mode, \(\Agrav\) is the anomaly-free quantum constraint algebra, \(Q_{\rm DM}\) is an exact dark charge or equivalent stability invariant, and \(P_{\rm val}(\Dil)\) is the finite dilation annihilator of the homogeneous valuation sector.
\end{definition}



\begin{remark}[Status of the capacity postulate]
The capacity postulate does not by itself compute \(\Ncap\).  It replaces a freely inserted cosmological constant by a sharply defined microscopic eigenvalue problem.  A full solution of the cosmological-constant amplitude problem requires an independent derivation of \(\Ncap\) from the boundary state space.  The present work supplies the invariant mechanism and the falsification criteria.
\end{remark}

\section{Layer separation and invariant atlas}\label{sec:atlas}

The construction uses three different finite alphabets.  They must not be identified.

\begin{definition}[Three alphabets]
The round valuation density alphabet is
\begin{equation}
\boxed{
S_{\rm val}^{\rho}:=\{1,0,-1,-2,-3\}.
}
\end{equation}
The boundary-local valuation density alphabet is
\begin{equation}
\boxed{
S_{\partial}^{\rho}:=\{0,-1,-2,-3\}.
}
\end{equation}
If ordinary radiation is added to the Friedmann ledger, the expansion alphabet used for \(F=H^2/H_0^2\) is
\begin{equation}
\boxed{
S_{\rm cos}^{F}:=\{0,-1,-2,-3,-4\}.
}
\end{equation}
The radiation exponent \(s=-4\) is not produced by the four-dimensional Euclidean valuation density on the round \(B^4/S^3\) family.  It is an external standard source added to the cosmological ledger.
\end{definition}

\begin{criterion}[Exact invariant atlas]\label{crit:atlas}
A proposed completion passes the invariant sieve only if the following sectors are simultaneously well defined.
\begin{center}
\footnotesize
\setlength{\tabcolsep}{2pt}
\begin{tabular}{@{}p{0.23\columnwidth}p{0.31\columnwidth}p{0.37\columnwidth}@{}}
\toprule
Sector & invariant object & sharp test or equation\\
\midrule
local gravity & \(\Ptr\) & \(\Ptr(-\rho_0g_{\mu\nu})=0\)\\
\(\Lambda\)-mode & \(\Ncap=\Capacity\mathcal H_\partial\) & \(\Lambda_{\rm cap}=3\pi/(\ell_P^2\Ncap)\)\\
quantum covariance & \(\Agrav\) & anomaly-free constraint algebra\\
dark matter & \(Q_{\rm DM}\) & \(\nabla_\mu J^\mu_{\rm DM}=0\)\\
valuation density & \(P_{\rm val}(\Dil)\) & \(P_{\rm val}(\Dil)\rho_{\rm val}=0\)\\
boundary quotient & \([a^{-1}]\) & \(w_X=-2/3\)\\
curvature split & \(\mathcal C_{-2}[F]-\Omega_K^{\rm geom}\) & \(\Omega_{2,\rm val}\)\\
spectroscopy & Hankel pencil \((H_1,H_0)\) & free recovery of \(s_j\)\\
positivity & inertia of \(H_0^{\rm phys}\) & \(\nu_-=0\) after split\\
spin--2 lift & tensor causal cone & \(c_T=c\)\\
equation of state & cumulant closure & nonlinear Bell identity\\
\bottomrule
\end{tabular}
\end{center}
The atlas is a sieve, not a fit.  A candidate that fails one line has not solved the corresponding part of the problem.
\end{criterion}

\begin{remark}[Corrections relative to a loose formulation]
There are four common confusions that this manuscript excludes.  First, \(N_\Lambda=3\pi/(\Lambda\ell_P^2)\) is of order \(10^{122}\) for the observed universe, not of order \(3\pi\).  Second, \(P_{\rm val}(\Dil)\) annihilates the density profile, not the equation-of-state ratio \(w(a)\).  Third, a dark charge is a conserved current or protected discrete/topological charge, not a bare linear combination of \(\Omega_\Lambda\) and \(\Omega_X\).  Fourth, the homogeneous \(s=-2\) expansion branch is not the spin--2 graviton propagation mode; tensor propagation belongs to the perturbative lift.
\end{remark}

\section{Why a bare cosmological constant is the wrong invariant}

A bare cosmological constant is a dimensional term.  It can be made small or large by units; only the combination
\begin{equation}
\boxed{
\Lambda\ell_P^2
}
\end{equation}
or, equivalently,
\begin{equation}
\boxed{
N_\Lambda:=\frac{3\pi}{\Lambda\ell_P^2}
}
\end{equation}
is invariant.  For the observed late-time universe, \(N_\Lambda\) is of order \(10^{122}\).  The factor \(3\pi\) is only the numerator determined by the de Sitter area formula; it is not the value of the capacity.

\begin{nogotheorem}[No topology-only dimensional amplitude]\label{nogo:topology-amplitude}
Pure topology or other dimensionless data cannot naturally determine a nonzero dimensional constant such as \(\Lambda\), \(H_0\), a particle mass, or a vacuum energy density while remaining covariant under global length rescalings.
\end{nogotheorem}

\begin{proof}
Let \(X\) have length dimension \(L^k\), \(k\ne0\).  Under \(L\mapsto\lambda L\), \(X\mapsto\lambda^kX\).  Pure topology and dimensionless data are unchanged.  If \(X\) is determined by such data alone, it must both transform and remain fixed.  For all \(\lambda>0\), this forces \(X=0\).
\end{proof}

\begin{criterion}[Solved versus renamed cosmological constant]
A proposal solves the cosmological-constant amplitude problem only if it supplies a mechanism determining \(N_\Lambda\), or a distribution plus measure selecting it.  Merely rewriting
\[
N_\Lambda=\frac{3\pi}{\Lambda\ell_P^2}
\]
after \(\Lambda\) has been observed renames the problem but does not solve it.
\end{criterion}

\section{Trace-free local dynamics}

The local gravitational equation is postulated to be
\begin{equation}
\boxed{
R_{\mu\nu}-\frac14Rg_{\mu\nu}
=
8\pi G\left(T_{\mu\nu}-\frac14Tg_{\mu\nu}\right).
}
\label{eq:tracefree}
\end{equation}
Equivalently,
\begin{equation}
\boxed{
\Ptr R_{\mu\nu}=8\pi G\,\Ptr T_{\mu\nu}.
}
\end{equation}

\begin{theorem}[Vacuum-energy kernel]\label{thm:vacuum-kernel}
Let
\begin{equation}
T_{\mu\nu}\longmapsto T_{\mu\nu}'=T_{\mu\nu}-\rho_0g_{\mu\nu}
\end{equation}
with constant \(\rho_0\).  Then
\begin{equation}
\Ptr T'_{\mu\nu}=\Ptr T_{\mu\nu}.
\end{equation}
Thus a local constant vacuum-energy shift does not source the trace-free equation \eqref{eq:tracefree}.
\end{theorem}

\begin{proof}
In four dimensions \(T'=T-4\rho_0\).  Therefore
\begin{align}
T'_{\mu\nu}-\frac14T'g_{\mu\nu}
&=
T_{\mu\nu}-\rho_0g_{\mu\nu}
-\frac14(T-4\rho_0)g_{\mu\nu}\\
&=
T_{\mu\nu}-\frac14Tg_{\mu\nu}.
\end{align}
\end{proof}

\begin{theorem}[Bianchi integration mode]\label{thm:bianchi-integration}
Assume \(\nabla^\mu T_{\mu\nu}=0\).  Then \eqref{eq:tracefree} implies
\begin{equation}
\nabla_\nu(R+8\pi G T)=0.
\end{equation}
Hence on every connected spacetime region
\begin{equation}
\boxed{
R+8\pi G T=4\Lambda_{\rm int}
}
\label{eq:traceconstant}
\end{equation}
for a constant \(\Lambda_{\rm int}\), and \eqref{eq:tracefree} is equivalent to
\begin{equation}
\boxed{
G_{\mu\nu}+\Lambda_{\rm int}g_{\mu\nu}
=
8\pi G\,T_{\mu\nu}.
}
\label{eq:einstein-int}
\end{equation}
\end{theorem}

\begin{proof}
Taking the divergence of \eqref{eq:tracefree} and using the contracted Bianchi identity gives
\[
\nabla^\mu R_{\mu\nu}-\frac14\nabla_\nu R
=
\frac12\nabla_\nu R-\frac14\nabla_\nu R
=
\frac14\nabla_\nu R.
\]
The right-hand side has divergence
\[
8\pi G\left(\nabla^\mu T_{\mu\nu}-\frac14\nabla_\nu T\right)
=
-2\pi G\nabla_\nu T.
\]
Thus \(\nabla_\nu(R+8\pi GT)=0\), proving \eqref{eq:traceconstant}.  Substitution of
\[
R=4\Lambda_{\rm int}-8\pi GT
\]
back into \eqref{eq:tracefree} gives
\[
G_{\mu\nu}+\Lambda_{\rm int}g_{\mu\nu}=8\pi GT_{\mu\nu}.
\]
\end{proof}

\begin{remark}[Vacuum-shift equivalence class]
The trace-free equation is invariant under \(T_{\mu\nu}\mapsto T_{\mu\nu}-\rho_0g_{\mu\nu}\), but the reconstructed Einstein form is invariant only if either \(T_{\mu\nu}\) is understood modulo this kernel or the integration constant shifts as
\[
\Lambda_{\rm int}\mapsto\Lambda_{\rm int}+8\pi G\rho_0 .
\]
The capacity postulate fixes the representative by assigning the integration mode from boundary entropy capacity, rather than by choosing a local vacuum counterterm.
\end{remark}

\begin{remark}[Relation to unimodular gravity]
The trace-free equation is the local equation familiar from unimodular formulations \cite{AndersonFinkelstein1971,Unruh1989,HenneauxTeitelboim1989,Weinberg1989}.  The new element here is not the trace-free equation by itself.  It is the invariant prescription that fixes the resulting integration mode by a horizon-capacity eigenvalue.
\end{remark}

\section{Horizon-capacity fixation}

Let the late-time cosmological sector have a future de Sitter horizon.  Its radius and area are
\begin{equation}
R_{\rm dS}=\sqrt{\frac3{\Lambda}},
\qquad
A_{\rm dS}=4\pi R_{\rm dS}^2=\frac{12\pi}{\Lambda}.
\end{equation}
The Bekenstein--Hawking--Gibbons entropy normalization gives \cite{Bekenstein1973,Hawking1975,GibbonsHawking1977}
\begin{equation}
S_{\rm dS}=\frac{A_{\rm dS}}{4\ell_P^2}
=
\frac{3\pi}{\Lambda\ell_P^2}.
\end{equation}

\begin{definition}[Boundary entropy capacity]\label{def:entropy-capacity}
The symbol
\[
\Ncap=\Capacity(\mathcal H_\partial;\Agrav,Q_{\rm DM},P_{\rm val})
\]
denotes a logarithmic state-counting capacity, normalized so that in the semiclassical de Sitter sector it equals the horizon entropy rather than the Hilbert-space dimension itself.  Thus a finite microcanonical realization has
\[
\dim\mathcal H_\partial\sim e^{\Ncap},
\]
not \(\dim\mathcal H_\partial=\Ncap\).
\end{definition}

\begin{hypothesis}[Boundary capacity postulate]\label{hyp:capacity}
The physical cosmological boundary sector has a finite entropy-capacity eigenvalue
\begin{equation}
\boxed{
\Ncap
:=
\Capacity(\mathcal H_\partial;\Agrav,Q_{\rm DM},P_{\rm val})
}
\end{equation}
defined after imposing:
\begin{enumerate}
\item the anomaly-free quantum constraint algebra \(\Agrav\);
\item the dark-charge superselection rule \(Q_{\rm DM}\);
\item the finite homogeneous valuation closure \(P_{\rm val}(\Dil)\).
\end{enumerate}
In the semiclassical de Sitter limit this capacity equals the horizon entropy:
\begin{equation}
\boxed{
\Ncap=S_{\rm dS}.
}
\end{equation}
\end{hypothesis}

\begin{theorem}[Capacity-fixed cosmological mode]\label{thm:capacity-mode}
Under \cref{hyp:capacity}, once the entropy-capacity eigenvalue is supplied or computed, the integration constant in \eqref{eq:einstein-int} is fixed by
\begin{equation}
\boxed{
\Lambda_{\rm int}=\Lambda_{\rm cap}
=
\frac{3\pi}{\Ncap\ell_P^2}
=
\frac{3\pi}
{\ell_P^2\Capacity(\mathcal H_\partial;\Agrav,Q_{\rm DM},P_{\rm val})}.
}
\end{equation}
\end{theorem}

\begin{proof}
By \cref{hyp:capacity},
\[
\Ncap=S_{\rm dS}=\frac{3\pi}{\Lambda_{\rm cap}\ell_P^2}.
\]
Solving for \(\Lambda_{\rm cap}\) gives the result.  For a given positive entropy capacity \(\Ncap>0\), the value is single-valued.
\end{proof}

\begin{corollary}[Replacement of the cosmological constant]\label{cor:lambda-replacement}
The gravitational field equation becomes
\begin{equation}
\boxed{
G_{\mu\nu}
+
\frac{3\pi}{\Ncap\ell_P^2}g_{\mu\nu}
=
8\pi G\,T_{\mu\nu}.
}
\label{eq:capacity-einstein}
\end{equation}
The term proportional to \(g_{\mu\nu}\) is not a free bare cosmological constant.  It is an on-shell integration mode fixed once the boundary entropy capacity is supplied or computed.
\end{corollary}

\begin{theorem}[Vacuum-shift immunity of the capacity mode]\label{thm:vacuum-shift-immunity}
A shift \(T_{\mu\nu}\mapsto T_{\mu\nu}-\rho_0g_{\mu\nu}\) cannot change \(\Lambda_{\rm cap}\) unless it changes the physical boundary capacity \(\Ncap\).
\end{theorem}

\begin{proof}
By \cref{thm:vacuum-kernel}, the local trace-free equation is invariant under the shift.  By \cref{thm:capacity-mode}, the integration mode is a function only of \(\Ncap\) and \(\ell_P\).  Therefore the shift changes \(\Lambda_{\rm cap}\) only if it changes \(\Ncap\), which is a global state-counting datum rather than a local zero-point counterterm.
\end{proof}

\begin{criterion}[Exact cosmological-constant null]\label{crit:lambda-null}
The invariant residual for the replacement mechanism is
\begin{equation}
\boxed{
\mathcal R_\Lambda
:=
\Lambda_{\rm obs}
-
\frac{3\pi}
{\ell_P^2\Capacity(\mathcal H_\partial;\Agrav,Q_{\rm DM},P_{\rm val})}.
}
\end{equation}
A completed microscopic theory must give \(\mathcal R_\Lambda=0\) within observational and state-counting uncertainty.  Failure to compute the capacity is not a numerical failure of \eqref{eq:capacity-einstein}; it marks the remaining microscopic boundary state-counting task.
\end{criterion}


\section{Dynamical dark-energy update: capacity flow rather than a varying bare constant}\label{sec:dynamic-capacity}

Recent DESI DR2 analyses do not establish that a literal cosmological constant is increasing.  A cosmological constant is, by definition, constant.  The empirical statement is sharper and weaker: combinations of DESI BAO with CMB and supernova data show a preference, still below the discovery threshold and sensitive to dataset consistency, for a time-evolving dark-energy equation of state.  The DESI DR2 key paper reports that the preferred \(w_0w_a\) solution lies in the quadrant
\[
w_0>-1,\qquad w_a<0,
\]
with a \(3.1\sigma\) preference over \(\Lambda\)CDM for DESI BAO plus CMB and \(2.8\)--\(4.2\sigma\) when supernova samples are also included \cite{DESIDR2BAO2025}.  The DESI extended dark-energy paper reports that two-parameter \(w(z)\) descriptions capture the observed trends and that the evidence for dynamical dark energy is especially stable at low redshift \cite{DESIDR2Extended2025}.  Critical analyses emphasize that dataset tensions and systematics still matter, so this is not yet a final discovery of a unique dark-energy law \cite{WangMota2025}.

The invariant response is not to promote \(\Lambda\) back to a local variable.  The trace-free sector continues to remove local vacuum-energy shifts.  What may vary is the effective dark-energy density
\[
\Lambda_{\rm eff}(a):=\frac{8\pi G}{c^4}\rho_{\rm DE}(a),
\]
or equivalently the boundary capacity
\[
\boxed{
N_\partial(a):=\frac{3\pi}{\ell_P^2\Lambda_{\rm eff}(a)}.
}
\]
Thus a time-dependent dark-energy density is a capacity-flow measurement.

\begin{definition}[Capacity-flow exponent]
Let
\[
\Dil:=a\frac{d}{da}.
\]
For any effective dark-energy sector with \(\rho_{\rm DE}>0\), define
\[
s_{\rm eff}(a):=\Dil\ln\rho_{\rm DE}(a)
=\Dil\ln\Lambda_{\rm eff}(a),
\]
and
\[
\nu_\partial(a):=\Dil\ln N_\partial(a).
\]
Then
\[
\boxed{
\nu_\partial(a)=-s_{\rm eff}(a).
}
\]
Using the continuity equation
\[
\Dil\ln\rho_{\rm DE}=-3[1+w_{\rm eff}(a)],
\]
one obtains the invariant diagnostic
\[
\boxed{
s_{\rm eff}(a)=-3[1+w_{\rm eff}(a)],
\qquad
\nu_\partial(a)=3[1+w_{\rm eff}(a)].
}
\]
\end{definition}

\begin{remark}[No varying Bianchi constant]
The quantity \(\Lambda_{\rm eff}(a)\) in this section is an effective dark-energy density written in cosmological-constant units.  It is not the Bianchi integration constant \(\Lambda_{\rm int}\) of \cref{thm:bianchi-integration}.  A time-dependent \(\Lambda_{\rm eff}(a)\) is compatible with the preceding trace-free construction only as part of a conserved enlarged stress tensor or a covariant boundary-exchange law.
\end{remark}

For the CPL form
\[
w(a)=w_0+w_a(1-a),
\]
the exact size-dependent factor is
\[
\boxed{
\frac{\Lambda_{\rm eff}(a)}{\Lambda_0}
=
a^{-3(1+w_0+w_a)}
\exp\!\left[3w_a(a-1)\right],
}
\]
and the corresponding capacity factor is
\[
\boxed{
\frac{N_\partial(a)}{N_{\partial0}}
=
a^{3(1+w_0+w_a)}
\exp\!\left[-3w_a(a-1)\right].
}
\]
The instantaneous crossing of the phantom divide, if \(w_a\ne0\), occurs at
\[
\boxed{
a_\times=1+\frac{1+w_0}{w_a},
\qquad
z_\times=a_\times^{-1}-1.
}
\]
For the DESI-favored quadrant \(w_0>-1,w_a<0\), this crossing lies in the past when \(0<a_\times<1\).  In that case \(w<-1\) before the crossing and \(w>-1\) after it.  Equivalently, \(\Lambda_{\rm eff}\) grows with scale while \(w<-1\), reaches an extremum at \(w=-1\), and decreases with scale while \(w>-1\).

\begin{theorem}[Static valuation interpretation of an increasing effective \(\Lambda\)]
In a single valuation branch,
\[
\rho_s(a)\propto a^s,
\qquad
w_s=-1-\frac{s}{3},
\qquad
\frac{N_s(a)}{N_{s0}}=a^{-s}.
\]
Thus an increasing effective dark-energy density,
\[
\Dil\ln\Lambda_{\rm eff}>0,
\]
requires
\[
s_{\rm eff}>0
\qquad\Longleftrightarrow\qquad
w_{\rm eff}<-1.
\]
Within the static positive \(B^4/S^3\) valuation spectrum, the only positive root is
\begin{equation}
\boxed{\begin{gathered}
s=1,\qquad w=-\frac43,\\
\Lambda_{\rm eff}\propto a,\qquad N_\partial\propto a^{-1}.
\end{gathered}}
\end{equation}
The boundary-local quotient branch \(s=-1\) instead gives
\begin{equation}
\boxed{\begin{gathered}
s=-1,\qquad w=-\frac23,\\
\Lambda_{\rm eff}\propto a^{-1},\qquad N_\partial\propto a.
\end{gathered}}
\end{equation}
Therefore a genuine increasing-\(\Lambda_{\rm eff}\) epoch cannot be represented by the positive boundary-local \(a^{-1}\) branch.  It either selects the full bulk valuation root \(s=1\), or it requires a controlled interacting/sign-changing capacity-flow extension.
\end{theorem}

\begin{proof}
For a power branch \(\rho_s\propto a^s\), the continuity equation gives \(s=-3(1+w_s)\), hence \(w_s=-1-s/3\).  Since \(N_\partial\propto\Lambda_{\rm eff}^{-1}\), one has \(N_s/N_{s0}=a^{-s}\).  The density increases with scale exactly when \(s>0\).  In the finite valuation spectrum \(\{1,0,-1,-2,-3\}\), the only positive exponent is \(s=1\).  The boundary-local spectrum \(\{0,-1,-2,-3\}\) contains no positive exponent.
\end{proof}

\begin{corollary}[Positive-mixture monotonicity test]
For a positive finite valuation mixture,
\[
\Dil w=-\frac13{\rm Var}_\pi(s)\le0.
\]
Hence a reconstructed epoch with \(\Dil w>0\) cannot be a positive noninteracting valuation mixture on a fixed support.  If confirmed, the DESI-type trend from phantom-like behavior in the past to quintessence-like behavior today selects one of three possibilities:
\statementbox{bulk valuation activation, interacting branch exchange, or sign-indefinite effective reconstruction}
It does not invalidate the valuation annihilator.  It identifies the boundary-positive subclass as too small for that epoch.
\end{corollary}

\begin{theorem}[Interacting two-branch capacity flow and observable exchange current]
\label{thm:interacting-current}
Let
\[
\begin{gathered}
\rho(a)=C_-(a)a^{s_-}+C_+(a)a^{s_+},\\
C_\pm(a)>0,\qquad s_-<s_+
\end{gathered}
\]
with branch fractions \(\pi_\pm\) and
\[
\kappa_1:=\Dil\ln\rho .
\]
Define the exchange current
\[
J_+:=\Dil\pi_+-(s_+-\kappa_1)\pi_+,
\qquad
J_-:=-J_+ .
\]
Then
\[
\sum_sJ_s=0,
\qquad
\Dil\kappa_1=\kappa_2+(s_+-s_-)J_+ .
\]
Since
\[
\kappa_1=-3(1+w),
\]
and, for a two-point support,
\[
\kappa_2=(s_+-\kappa_1)(\kappa_1-s_-),
\]
the current is recovered pointwise from the observable pair \((w,\Dil w)\):
\begin{equation}
\boxed{
J_+(a)=
\frac{
-3\Dil w-(s_+-\kappa_1)(\kappa_1-s_-)
}{
s_+-s_-
}.
}
\label{eq:observable-current}
\end{equation}
Hence \(J_+\equiv0\) if and only if the noninteracting replicator closure holds.  Moreover, \(\Dil w>0\) on an interval if and only if
\[
J_+<-\frac{\kappa_2}{s_+-s_-}<0
\]
there; capacity then flows from the upper branch to the lower branch while both branch weights remain positive.
\end{theorem}

\begin{proof}
From \(\pi_++\pi_-=1\) and
\[
\kappa_1=s_-+(s_+-s_-)\pi_+,
\]
one has
\[
\Dil\kappa_1=(s_+-s_-)\Dil\pi_+ .
\]
Substituting
\[
\Dil\pi_+=(s_+-\kappa_1)\pi_++J_+
\]
and using
\[
(s_+-s_-)(s_+-\kappa_1)\pi_+=\kappa_2
\]
for a two-point distribution gives the stated identity.  The variance formula is
\[
\kappa_2=\pi_+\pi_-(s_+-s_-)^2
=(s_+-\kappa_1)(\kappa_1-s_-).
\]
\end{proof}

\begin{corollary}[Phantom quadrant under positivity]
A reconstructed epoch with \(w<-1\) in the past, \(w>-1\) today, and \(\Dil w>0\) cannot be a positive noninteracting mixture on any fixed support.  It is realized by a positive interacting two-branch mixture if and only if
\[
\kappa_1\in(s_-,s_+),
\qquad
J_+<-\frac{\kappa_2}{s_+-s_-}.
\]
For the valuation pair \((s_-,s_+)=(-1,1)\), this gives an exact amplitude-free diagnostic of the DESI-type quadrant.
\end{corollary}

\begin{criterion}[Interacting capacity-flow null]
The observable residual
\[
\boxed{\mathcal R_{\rm int}(a):=J_+(a)}
\]
vanishes identically for the noninteracting valuation closure.  A statistically nonzero reconstruction of \(\mathcal R_{\rm int}\) with the sign pattern of \cref{thm:interacting-current} is the invariant signature of interacting capacity exchange, distinct from both signed negative-weight shadow fits and bulk-valuation-root activation.
\end{criterion}

\begin{criterion}[Dynamical-capacity null]
The invariant observational null replacing a varying bare cosmological constant is
\[
\boxed{
\mathcal R_{\rm cap}(a):=
\Dil\ln N_\partial(a)-3[1+w_{\rm eff}(a)]=0.
}
\]
For a static valuation branch this reduces to
\[
\mathcal R_{\rm cap}=0,\qquad
\nu_\partial=-s\in\{-1,0,1,2,3\}.
\]
For a free-spectrum reconstruction one shall first infer \(s_{\rm eff}\) or the discrete pencil roots from the data, and only afterwards compare them with
\[
S_{\rm val}=\{1,0,-1,-2,-3\}.
\]
This is the same consistency-versus-identification boundary already required by the invariant test architecture.
\end{criterion}

Using the Planck-length value \( \ell_P=1.616255\times10^{-35}\,{\rm m}\) \cite{NISTPlanckLength2022}, and a representative \(\Lambda\)CDM normalization \(H_0=67.4\,{\rm km\,s^{-1}\,Mpc^{-1}}\), \(\Omega_\Lambda\simeq0.685\), one obtains
\[
\Lambda_0\simeq1.091\times10^{-52}\,{\rm m}^{-2},
\]
and therefore
\[
\boxed{
N_{\partial0}
=
\frac{3\pi}{\Lambda_0\ell_P^2}
\simeq3.31\times10^{122}.
}
\]
For the static valuation roots, the size-dependent capacity factors are
\[
\begin{array}{c|c|c|c}
s & w_s & \Lambda_s(a)/\Lambda_{s0} & N_s(a)/N_{s0}\\
\hline
1 & -4/3 & a & a^{-1}\\
0 & -1 & 1 & 1\\
-1 & -2/3 & a^{-1} & a\\
-2 & -1/3 & a^{-2} & a^2\\
-3 & 0 & a^{-3} & a^3
\end{array}
\]
This table is the exact size-dependent factorization of the capacity mechanism.




\section{Invariant frontier and all-dimensional closure ledger}
\label{sec:frontier}

This section records what is closed at theorem level and what remains a
microscopic frontier.  This distinction is part of the result.  A finite
annihilator, a projector, a sign certificate, or a null test may be exact
without determining a dimensional amplitude.

\subsection{All-dimensional valuation skeleton}

Let the spatial dimension be \(n\ge1\), with round section \(S^n_A\) and
auxiliary filling \(B^{n+1}_A\).  Since \(\vol(S^n_A)\propto A^n\), the
round restriction of a continuous rigid-motion invariant additive valuation
on convex bodies in \(\mathbb R^{n+1}\) has filling-energy powers
\[
E_{\rm val}(A)\in
\spanop\{A^{n+1},A^n,A^{n-1},\ldots,A,1\}.
\]
After division by the \(S^n_A\) volume, the density spectrum is
\[
S_{\rm val}^{(n)}=\{1,0,-1,\ldots,-n\}.
\]
Equivalently,
\begin{equation}
\boxed{
P_n(\Dil)\rho=0,\qquad
P_n(z)=(z-1)z(z+1)\cdots(z+n).
}
\end{equation}
All roots are simple and real, hence
\[
\mathbb R[z]/(P_n)\simeq \mathbb R^{\,n+2}.
\]
The boundary-local face removes the bulk intrinsic volume:
\[
S_{\partial}^{(n)}=\{0,-1,\ldots,-n\},
\qquad
P_{\partial,n}(z)=z(z+1)\cdots(z+n).
\]
For a branch \(\rho_s\propto a^s\), the \(n\)-dimensional work-map equation
of state is
\[
w_s=-1-\frac{s}{n}.
\]

\begin{theorem}[Dimension-sensitive boundary quotient]
\label{thm:dimension-quotient}
For \(n\ge3\), quotient the boundary alphabet
\(S_{\partial}^{(n)}=\{0,-1,\ldots,-n\}\) by the standard homogeneous
expansion powers carried by vacuum, geometric curvature, and dust,
\[
0,\qquad -2,\qquad -n .
\]
Then the quotient dimension is \(n-2\).  Hence the condition that there be
exactly one nonstandard boundary-local homogeneous direction selects
\(n=3\).  The surviving direction is then
\[
\boxed{
s=-1,\qquad
\rho_X\propto a^{-1},\qquad
w_X=-\frac23 .
}
\]
\end{theorem}

\begin{proof}
For \(n\ge3\), the three exponents \(0,-2,-n\) are distinct.  Since
\(|S_{\partial}^{(n)}|=n+1\), the quotient dimension is
\((n+1)-3=n-2\).  It equals one if and only if \(n=3\).  The remaining
exponent in \(\{0,-1,-2,-3\}\) after removing \(\{0,-2,-3\}\) is \(s=-1\),
and \(w_s=-1-s/3=-2/3\).
\end{proof}

\subsection{Closed ledger}

The following entries are closed under the hypotheses stated in the
corresponding sections.

First, constant vacuum shifts are removed by the trace-free source
projector:
\[
\Ptr(-\rho_0g_{\mu\nu})=0.
\]
Second, if the physical stress tensor is conserved, the trace-free equation
recovers a Bianchi integration mode,
\[
R+8\pi GT=4\Lambda_{\rm int},
\qquad
G_{\mu\nu}+\Lambda_{\rm int}g_{\mu\nu}=8\pi GT_{\mu\nu}.
\]
Third, the capacity postulate fixes this mode only after the boundary entropy
capacity is supplied:
\[
\Lambda_{\rm cap}=\frac{3\pi}{\ell_P^2N_\partial},
\qquad
N_\partial=\log\dim \mathcal H_{\partial}^{\rm eff}.
\]
Thus the manuscript does not yet compute \(\Lambda\); it turns it into a
boundary entropy-counting problem.

Fourth, positive noninteracting finite mixtures obey the arrow
\[
\Dil w=-\frac13{\rm Var}_\pi(s)\le0.
\]
A reconstructed epoch with \(\Dil w>0\) therefore requires a signed
effective reconstruction, interacting branch exchange, or an enlarged
functional class.  The interacting two-branch current in
\cref{thm:interacting-current} is the local observable diagnostic of the
second possibility.

Fifth, the \(a^{-2}\) coefficient must be split before positivity is read:
\[
\Omega_{2,{\rm val}}=\mathcal C_{-2}[F]-\Omega_K^{\rm geom}.
\]
Sixth, dark matter stability is encoded by a current,
\[
\nabla_\mu J^\mu_{\rm DM}=0,\qquad
\rho_{\rm DM}=m n_{\rm DM}\propto a^{-n},
\]
which gives \(a^{-3}\) in the observed \(n=3\) case.



\subsection{Microscopic closure and open frontier}

The microscopic amplitude problem is no longer diffuse.  It is compressed into
three explicit edge-theorem gates, stated and implemented in
\Cref{sec:edge-theorem-gates}.  None of the gates admits \(q_{\rm obs}\),
\(\Lambda_{\rm obs}\), or \(R_\Lambda\) as an input.  Their data are the edge
recursion, the registry field content, the APS/zeta convention, and the
Harish--Chandra/stringy horizon determinant.

\begin{enumerate}
\item \emph{Recursion Gate.}
The compact Hopf--Bradlow--APS--HC operator \(\mathcal R\) on de Sitter edge
characters has the terminal integer fixed point
\[
N_\varphi=292.
\]

\item \emph{Matching Gate.}
The registry GLSM, Chern--Simons level lock, APS convention, and
Kronecker-limit stringy edge character give the scheme-covariant ledger
\[
b\,g_\partial^2(\ell_P^{-1})
=
0.559754586\ldots,
\qquad
\Xi_{\rm edge}\simeq0.99916928 .
\]

\item \emph{Determinant-Line Gate.}
The projected horizon edge determinant is a capacity-sector function obeying
\[
-\log Z_{\rm edge}(q)
=
\frac12\ln q+\Theta_\partial(q)+S_{\rm HC}(q)+o(1).
\]
\end{enumerate}

With these gates in place, the conditional selector becomes the amplitude
theorem stated in \Cref{thm:edge-selector-conditional}.  The open frontier is
therefore not another fit for \(\Lambda\), but independent reproduction and
falsification: a finite-cutoff spectral evaluator for \(\mathcal R\), a public
matching ledger, a zeta/Harish--Chandra determinant evaluator, full joint
cosmological chains, and laboratory or analog tests of the same edge
operators.  All other entries of the atlas---quantizers, bounds,
obstructions, registries, null tests, and observational diagnostics---serve
only to protect, reproduce, or falsify this edge-selection chain.


\section{Illustrative real-data audit: DESI DR2 BAO}\label{sec:real-data-audit}

The invariant sieve must not be judged by rhetoric.  We therefore include an illustrative, reproducible background audit against published DESI DR2 BAO distance measurements.  The data vector used here is the DESI DR2 BAO Table 4 cosmology vector: \(D_V/r_d\) for BGS and \(D_M/r_d,D_H/r_d\) for the six combined high-redshift tracers, including the combined LRG3+ELG1 bin rather than the two correlated sub-bins \cite{DESIDR2BAO2025}.  The calculation is a BAO-only background audit.  It is not a joint CMB+BAO+SNe+growth likelihood, not a substitute for the official DESI chains, and not a full covariance public-likelihood reproduction.

We compare the following finite models:
\[
\Lambda{\rm CDM},\qquad
X:{F(a)=\Omega_m a^{-3}+\Omega_Xa^{-1}+\Omega_\Lambda},
\]
\[
P:{F(a)=\Omega_m a^{-3}+\Omega_Pa+\Omega_\Lambda},
\qquad
X+P,
\]
and the empirical \(w_0w_a\)CDM benchmark.  The nuisance scale \(C=c/(H_0r_d)\) is profiled.  Positivity is imposed in the physical valuation fits unless explicitly stated.

\begin{table*}[t]
\caption{DESI DR2 BAO-only audit of the capacity/valuation background families.  The \(\chi^2\) values use the published BAO distance vector and the diagonal/limited-correlation approximation implemented in the accompanying audit files.  AIC and BIC are quoted relative to \(\Lambda\)CDM.  The table is an illustrative reproducibility check of the invariant sieve, not a full cosmological model comparison.}
\label{tab:bao-audit}
\begin{ruledtabular}
\begin{tabular}{lcccccc}
model & \(k\) & \(\chi^2\) & dof & \(\Delta\chi^2\) & \(\Delta{\rm AIC}\) & \(\Delta{\rm BIC}\)\\
\hline
\(w_0w_a\)CDM & 4 & 5.812 & 9 & -4.727 & -0.727 & 0.403\\
\(\Lambda\)CDM & 2 & 10.539 & 11 & 0 & 0 & 0\\
\(P\) signed & 3 & 8.971 & 10 & -1.568 & 0.432 & 0.997\\
\(X+P\) signed & 4 & 7.085 & 9 & -3.453 & 0.547 & 1.677\\
\(X\) positive & 3 & 9.530 & 10 & -1.009 & 0.991 & 1.556\\
\(P\) positive & 3 & 10.539 & 10 & 0 & 2.000 & 2.565\\
\(X+P\) positive & 4 & 9.530 & 9 & -1.009 & 2.991 & 4.121
\end{tabular}
\end{ruledtabular}
\end{table*}

The best positive \(X\)-branch fit gives
\[
\boxed{
\Omega_X=0.146840,\qquad
\Omega_m=0.293136,\qquad
\chi^2=9.5299.
}
\]
This is a consistency-level improvement over \(\Lambda\)CDM in this BAO-only audit, not a detection.  The positive \(s=+1\) branch is driven to
\[
\boxed{\Omega_P\simeq0,}
\]
so these BAO data do not select a positive increasing-\(\Lambda_{\rm eff}\propto a\) branch.  Signed \(P\) and signed \(X+P\) fits can improve \(\chi^2\), but their negative branch weights violate the positive-branch interpretation and must be treated as effective reconstructions rather than positive valuation fluids.

\begin{figure*}[!htbp]
\centering
\includegraphics[width=.92\textwidth]{fig_model_comparison.png}
\caption{BAO-only comparison of the tested background families.  The flexible \(w_0w_a\) benchmark has the lowest \(\chi^2\); BIC still favors the simpler \(\Lambda\)CDM model over this small data vector.  The positive \(X:a^{-1}\) branch is viable but not superior.}
\label{fig:model-comparison}
\end{figure*}

\begin{figure*}[!htbp]
\centering
\includegraphics[width=.92\textwidth]{fig_capacity_flow_DESI_w0wa.png}
\caption{Capacity-flow representation of the official DESI-style \(w_0w_a\) trend.  A time-dependent dark-energy density is represented as a flow of \(N_\partial(a)=3\pi/[\ell_P^2\Lambda_{\rm eff}(a)]\), not as a varying bare cosmological constant.}
\label{fig:capacity-flow}
\end{figure*}

The absolute normalization check uses \(H_0=67.4\,{\rm km\,s^{-1}\,Mpc^{-1}}\), \(\Omega_\Lambda=0.685\), and \(\ell_P=1.616255\times10^{-35}\,{\rm m}\) \cite{Planck2018,NISTPlanckLength2022}.  It gives
\[
\boxed{
\Lambda_0=1.091\times10^{-52}\,{\rm m}^{-2},\qquad
N_{\partial0}=3.31\times10^{122}.
}
\]
The exact capacity-flow null remains
\[
\boxed{
\mathcal R_{\rm cap}(a)
=
\Dil\ln N_\partial(a)-3[1+w_{\rm eff}(a)]=0.
}
\]
The illustrative audit conclusion is therefore precise:
\statementbox{DESI DR2 BAO in this audit does not falsify the invariant sieve; it does not yet prove superiority of the simple positive \(a^{-1}\) valuation branch.  The data favor using the dynamical capacity-flow formalism and free-spectrum reconstruction before making an identification claim.}


\section{Quantum covariance algebra}

The capacity in \cref{hyp:capacity} is not a raw dimension of an arbitrary vector space.  It is counted after the gravitational constraints are imposed.  Classically, general relativity is a constrained Hamiltonian system:
\begin{equation}
\mathcal H_\perp(x)=0,
\qquad
\mathcal H_i(x)=0,
\end{equation}
with the hypersurface-deformation algebra \cite{Dirac1958,ADM1962,HenneauxTeitelboim1992,Dirac1964}.

\begin{definition}[Anomaly-free quantum covariance algebra]\label{def:agrav}
A quantum gravitational completion supplies a domain, state space, and constraint operators
\begin{equation}
\widehat{\mathcal H}_A=(\widehat{\mathcal H}_\perp,\widehat{\mathcal H}_i)
\end{equation}
such that
\begin{equation}
\boxed{
\widehat{\mathcal H}_A\Psi=0,
\qquad
[\widehat{\mathcal H}_A,\widehat{\mathcal H}_B]
=
i\hbar\,\widehat f_{AB}^{\ \ C}\widehat{\mathcal H}_C
}
\end{equation}
on the physical domain, with no anomalous obstruction.  This structure is denoted \(\Agrav\).
\end{definition}

\begin{nogotheorem}[Anomalous constraint algebras do not fix capacity]\label{nogo:anomaly}
If the quantum constraints do not close on the physical domain, then \(\mathcal H_\partial\) is not a physical boundary state space for spacetime covariance and its dimension cannot define \(\Ncap\).
\end{nogotheorem}

\begin{proof}
The constraints generate changes of spatial coordinates and refoliations.  If their quantum representatives fail to close, the supposed physical state space depends on the order of gauge projection.  Such a space cannot carry an invariant state count.
\end{proof}


\section{Boundary-constraint quantum kernel}\label{sec:quantum-kernel}

The previous section gives the algebraic admissibility condition.  To avoid leaving the quantum layer as a slogan, we state the minimal physical-space construction.  The kinematical space is
\begin{equation}
\boxed{
\Hkin=
\mathcal H_{\rm geom}(\Sigma)
\otimes
\mathcal H_{\rm matter}
\otimes
\Hbd
\otimes
L^2(\R_\Lambda).
}
\end{equation}
The last factor carries the global integration mode.  A conjugate four-volume variable may be introduced by
\begin{equation}
[\widehat{\Tfour},\widehat\Lambda]=i\hbar,
\end{equation}
which is the quantum form of the unimodular fact that \(\Lambda\) is conjugate to spacetime four-volume.

\begin{definition}[Invariant physical space]\label{def:physical-projector}
The admissible physical space is the common kernel, or equivalently a rigging-map image when the constraints are implemented by group averaging,
\begin{equation}
\boxed{
\begin{aligned}
\Hphys={}&
\ker\widehat{\mathcal H}_A
\cap
\ker\widehat C_{\rm tf}
\cap
\ker\widehat C_{\rm cap}
\\
&{}\cap
\ker P_{\rm val}(\widehat\Dil)
\cap
\ker(\widehat Q_{\rm DM}-q_{\rm DM}) .
\end{aligned}}
\end{equation}
When the corresponding projectors are mutually well defined and commute on a common invariant domain \(\mathfrak D\), this common-kernel notation may be represented formally as
\[
\Hphys=
\Pi_{\rm HD}\Pi_{\rm tf}\Pi_{\rm cap}\Pi_{\rm val}\Pi_Q\,\Hkin .
\]
Without the common-domain and commutation property, the ordered product is not a definition.
\end{definition}

\begin{remark}[Projection order and common domain]
The capacity, valuation, and dark-charge conditions are imposed as abelian superselection conditions and are required to commute with the hypersurface-deformation constraints on the common domain \(\mathfrak D\).  This is a hypothesis on the microscopic completion, not a theorem of the kinematical notation.  If it fails, different projection orders define different spaces and the candidate model fails the quantum part of the invariant sieve.
\end{remark}

The hypersurface-deformation projector is admissible only if the regulated constraints close on a common physical domain:
\begin{align}
[\widehat D[\vec N],\widehat D[\vec M]]
&=i\hbar\,\widehat D[[\vec N,\vec M]],\\
[\widehat D[\vec N],\widehat H[N]]
&=i\hbar\,\widehat H[\Lie_{\vec N}N],\\
[\widehat H[N],\widehat H[M]]
&=i\hbar\,\widehat D[\widehat q^{ab}
(N\partial_bM-M\partial_bN)].
\end{align}
The displayed relations are the hypersurface-deformation algebra in canonical form \cite{HKT1976}.  The last relation is an operator-algebroid relation because the structure function is the inverse spatial metric operator.  This is the point at which many candidate quantum gravities fail: good homogeneous dynamics is insufficient if refoliation covariance is anomalous.

\begin{criterion}[Weak anomaly-free completion]\label{crit:weak-anomaly}
A regulator or microscopic model passes the quantum part of the invariant sieve only if
\begin{equation}
\Delta_{AB}
:=
[\widehat C_A,\widehat C_B]
-i\hbar f_{AB}^{\ \ C}(\widehat q)\widehat C_C
\end{equation}
vanishes on the common physical domain, at least weakly:
\begin{equation}
\langle\Phi|\Delta_{AB}|\Psi\rangle=0,
\qquad
\Phi,\Psi\in\mathfrak D.
\end{equation}
\end{criterion}

\begin{theorem}[Capacity eigenvalue in the projected theory]\label{thm:capacity-eigenvalue-quantum}
If the capacity projector imposes
\begin{equation}
\widehat C_{\rm cap}:=\ell_P^2\widehat\Lambda\,\widehat N_\partial-3\pi=0
\end{equation}
on \(\Hphys\), then every simultaneous \((\widehat\Lambda,\widehat N_\partial)\)-eigen-sector satisfies
\begin{equation}
\boxed{
\Lambda_{\rm cap}=\frac{3\pi}{\ell_P^2N_\partial}.
}
\end{equation}
\end{theorem}

\begin{proof}
On a joint eigenstate,
\[
(\ell_P^2\widehat\Lambda\widehat N_\partial-3\pi)\Psi=0
\]
is exactly
\[
\ell_P^2\Lambda N_\partial=3\pi.
\]
For superpositions the statement holds sector by sector in the joint spectral decomposition whenever the two operators commute on the capacity domain.
\end{proof}

\begin{remark}[Status of the quantum kernel]
This construction is not a completed regulator, spin-foam sum, lattice path integral, or ultraviolet fixed-point calculation.  Concrete microscopic programs such as causal dynamical triangulations and functional-renormalization-group gravity supply examples of the kind of ultraviolet completion that would have to reproduce this kernel \cite{Loll2020,Saueressig2023}.  The present construction is the minimal invariant kernel that any such microscopic realization must reproduce: trace-free locality, boundary capacity, anomaly-free refoliation covariance, dark-charge superselection, and finite homogeneous \(S^3\)-valuation closure.
\end{remark}


\section{Dark matter as a protected charge sector}

A cold dark component is not explained merely by adding \(a^{-3}\) to the Friedmann ledger.  It must be stable, gravitating, and approximately collisionless.

\begin{hypothesis}[Dark charge]\label{hyp:dark-charge}
There exists an exact conserved current
\begin{equation}
\boxed{
\nabla_\mu J^\mu_{\rm DM}=0
}
\end{equation}
and a charge
\begin{equation}
\boxed{
Q_{\rm DM}=\int_\Sigma J^\mu_{\rm DM}n_\mu\,\dd\Sigma
}
\end{equation}
or an equivalent discrete or topological protection.  In the homogeneous cold limit,
\begin{equation}
\boxed{
T^{\rm DM}_{\mu\nu}=\rho_{\rm DM}u_\mu u_\nu,
\qquad
\rho_{\rm DM}\propto a^{-3},
\qquad
p_{\rm DM}\simeq0.
}
\end{equation}
\end{hypothesis}

\begin{remark}[Separation from the valuation dark-energy branch]
The \(s=-1\) valuation branch has \(w=-2/3\), not cold dust.  It must not be double counted as dark matter.  The cold \(s=-3\) sector is protected by \(Q_{\rm DM}\); the nonstandard homogeneous boundary branch \(s=-1\) is controlled by \(P_{\rm val}(\Dil)\).
\end{remark}


\begin{nogotheorem}[No linear density charge for nonzero exponents]\label{nogo:linear-density-charge}
Let
\[
Q[\rho](a)=\sum_s q_sC_sa^s
\]
be a linear functional of independent homogeneous density branches and require \(\Dil Q[\rho]=0\) for all \(a\) on an interval.  Then every active nonzero exponent in the support has \(q_sC_s=0\).  Hence a conserved dark charge cannot in general be a bare linear combination of \(\Omega_\Lambda\) and \(\Omega_X\).
\end{nogotheorem}

\begin{proof}
One has
\[
\Dil Q=\sum_s s\,q_sC_sa^s.
\]
Distinct powers are linearly independent.  Therefore \(s\,q_sC_s=0\) for each active branch.  Only \(s=0\) can survive without further structure.  A physical dark charge must instead be represented by a conserved current, a discrete charge, or a topological protection law.
\end{proof}

\section{\texorpdfstring{\(S^3\)}{S3} valuation closure}

Let a closed homogeneous spatial section be \(\Sigma\simeq\Sthree\) with round metric \(g_A=A^2\dd\Omega_3^2\).  Let \(\Bfour_A\) be the auxiliary Euclidean filling.  Under continuous, rigid-motion invariant, additive Euclidean valuation closure on convex bodies in \(\R^4\), Hadwiger theory gives the intrinsic-volume span \cite{Hadwiger1957,Schneider2014}.  On the round family this implies
\begin{equation}
E_4^{\rm val}(A)
=
\frac{\pi^2}{2}\gamma A^4
+
2\pi^2
\left(
\alpha_0A^3+3\alpha_1A^2+3\alpha_2A+\alpha_3
\right).
\end{equation}
Since \(V_3(A)=2\pi^2A^3\),
\begin{equation}
\begin{split}
\rho_{\rm val}(A)
&=\frac{E_4^{\rm val}(A)}{V_3(A)}\\
&=\frac{\gamma}{4}A+\alpha_0+3\alpha_1A^{-1}\\
&\quad+3\alpha_2A^{-2}+\alpha_3A^{-3}.
\end{split}
\end{equation}
Thus
\begin{equation}
\boxed{
\Spec_\Dil(\rho_{\rm val})\subseteq \{1,0,-1,-2,-3\},
\qquad
\Dil:=A\frac{\dd}{\dd A}.
}
\end{equation}

\begin{theorem}[Finite valuation annihilator]\label{thm:val-ann}
For \(C^5\) profiles on a positive interval,
\begin{equation}
\boxed{\begin{gathered}
P_{\rm val}(\Dil)\rho_{\rm val}=0,\\
P_{\rm val}(z)=(z-1)z(z+1)(z+2)(z+3).
\end{gathered}}
\end{equation}
if and only if
\begin{equation}
\rho_{\rm val}(A)\in
\spanop\{A,1,A^{-1},A^{-2},A^{-3}\}.
\end{equation}
\end{theorem}

\begin{proof}
With \(x=\ln A\), \(\Dil=\dd/\dd x\).  The equation becomes a constant-coefficient ODE with simple roots \(1,0,-1,-2,-3\), hence the stated exponential/power basis.
\end{proof}

\begin{corollary}[Boundary quotient]\label{cor:boundary-quotient}
In the boundary-local class the bulk \(s=1\) branch is removed:
\begin{equation}
S_\partial=\{0,-1,-2,-3\}.
\end{equation}
At the expansion level, quotienting by vacuum, curvature-degenerate, and dust powers \(\{1,a^{-2},a^{-3}\}\) leaves the unique nonstandard direction
\begin{equation}
\boxed{
\rho_X\propto a^{-1},
\qquad
w_X=-1-\frac{-1}{3}=-\frac23.
}
\end{equation}
\end{corollary}

\begin{remark}[What valuation does and does not fix]
The valuation algebra fixes the allowed homogeneous powers.  It does not fix the amplitude of the \(s=0\) branch.  In the present manuscript that amplitude is supplied by the capacity mode:
\[
\rho_\Lambda^{\rm cap}=\frac{\Lambda_{\rm cap}}{8\pi G}.
\]
Thus the \(s=0\) power is valuation-allowed, while its amplitude is capacity-fixed.
\end{remark}


\section{Nonlinear equation-of-state closure}\label{sec:w-closure}

The density closure is linear.  The equation-of-state closure is not.

Let
\begin{equation}
\rho(a)=\sum_{s\in S}C_sa^s,\qquad \rho>0,
\end{equation}
and define
\begin{equation}
\begin{aligned}
w&:=-1-\frac{1}{3}\frac{\Dil\rho}{\rho},\\
\kappa_1&:=\Dil\ln\rho=-3(1+w),\\
\kappa_{n}&:=\Dil^n\ln\rho\qquad(n\ge2).
\end{aligned}
\end{equation}
Then \(\kappa_n\) are the cumulants of the tilted finite spectral distribution.  Since
\begin{equation}
\frac{\Dil^n\rho}{\rho}=B_n(\kappa_1,\ldots,\kappa_n),
\end{equation}
with \(B_n\) the complete exponential Bell polynomial, the valuation annihilator gives the amplitude-free nonlinear identity
\begin{equation}
\boxed{
B_5+5B_4+5B_3-5B_2-6B_1=0.
}
\end{equation}
The boundary-local face gives
\begin{equation}
\boxed{
B_4+6B_3+11B_2+6B_1=0.
}
\end{equation}
For the endpoint support \(\{0,-3\}\),
\begin{equation}
\boxed{
\Dil w=3w(1+w).
}
\end{equation}

\begin{criterion}[Correct \(w\)-null]
The correct null is not \(P_{\rm val}(\Dil)w=0\).  Except for single branches and accidental degeneracies, \(w(a)\) is a ratio of finite exponential sums and satisfies the nonlinear cumulant closure above.  Observational tests based on \(w(a)\) must therefore propagate derivative covariance through this nonlinear identity.
\end{criterion}

\section{Capacity-fixed EDSC}

Let
\begin{equation}
F(a):=\frac{H^2(a)}{H_0^2},
\qquad
\Omega_K^{\rm geom}:=-\frac{1}{A_0^2H_0^2}
\end{equation}
in units \(c=1\).  Define the physical density and pressure represented by the homogeneous Einstein tensor:
\begin{equation}
\boxed{
E_F(a):=F(a)-\Omega_K^{\rm geom}a^{-2},
}
\end{equation}
\begin{equation}
\boxed{
P_F(a):=-E_F(a)-\frac13\Dil E_F(a).
}
\end{equation}
This is the homogeneous Bianchi identity in dilation form.

The capacity-fixed vacuum slot is
\begin{equation}
\boxed{
\Omega_\Lambda^{\rm cap}
=
\frac{\Lambda_{\rm cap}}{3H_0^2}
=
\frac{\pi}{H_0^2\ell_P^2\Ncap}.
}
\end{equation}
Therefore a finite capacity-fixed homogeneous ledger has the form
\begin{equation}
\boxed{\begin{aligned}
E_F(a)
={}&\Omega_\Lambda^{\rm cap}
+\Omega_Xa^{-1}
+\Omega_{2,\rm val}a^{-2}\\
&+\Omega_ma^{-3}
+\Omega_ra^{-4}.
\end{aligned}}
\label{eq:capacity-ledger}
\end{equation}
with
\begin{equation}
F(a)=E_F(a)+\Omega_K^{\rm geom}a^{-2}.
\end{equation}

\begin{theorem}[Homogeneous closure with capacity-fixed vacuum]\label{thm:capacity-edsc}
Assume closed FLRW symmetry, trace-free local dynamics, conserved physical stress, the boundary capacity postulate, and the finite valuation support for the nonstandard homogeneous sector.  Then the homogeneous gravitational sector is closed by \eqref{eq:capacity-ledger}.  The vacuum coefficient is not free:
\[
\Omega_\Lambda=\Omega_\Lambda^{\rm cap}.
\]
The remaining physical coefficients are separated from geometric curvature by
\[
\boxed{
\Omega_{2,\rm val}=\mathcal C_{-2}[F]-\Omega_K^{\rm geom},
}
\]
where \(\mathcal C_s\) are the exact dilation coefficient projectors.
\end{theorem}

\begin{proof}
\Cref{thm:bianchi-integration,thm:capacity-mode} give the Einstein equation with \(\Lambda_{\rm cap}\).  The first Friedmann equation gives \(F=E_F+\Omega_K^{\rm geom}a^{-2}\).  The Bianchi identity gives \(P_F=-E_F-\Dil E_F/3\).  The finite valuation closure gives the allowed nonstandard powers, while the standard cold and radiation sectors supply \(a^{-3}\) and \(a^{-4}\).  The coefficient projector statement follows from linear independence of the finite power alphabet.
\end{proof}

\section{Projectors, spectral recovery, and falsifiable nulls}

For a finite expansion alphabet
\[
S_{\rm cos}:=\{0,-1,-2,-3,-4\},
\qquad
F(a)=\sum_{s\in S_{\rm cos}}\Omega_s a^s,
\]
define
\begin{equation}
\boxed{
\mathcal C_s[F](a)
:=
a^{-s}
\prod_{\substack{t\in S_{\rm cos}\\t\ne s}}
\frac{\Dil-t}{s-t}F(a).
}
\end{equation}
Then \(\mathcal C_s[F]=\Omega_s\).

The exact nulls are:
\begin{align}
\mathcal R_\Lambda
&=
\Lambda_{\rm obs}
-
\frac{3\pi}{\ell_P^2\Ncap},
\\
\mathcal R_{\rm val}[\rho]
&=
(\Dil-1)\Dil(\Dil+1)(\Dil+2)(\Dil+3)\rho,
\\
\mathcal R_{\partial}[\rho]
&=
\Dil(\Dil+1)(\Dil+2)(\Dil+3)\rho,
\\
\mathcal R_X[F]
&=
\Dil\mathcal C_{-1}[F],
\\
\mathcal R_{2}[F]
&=
\mathcal C_{-2}[F]-\Omega_K^{\rm geom}-\Omega_{2,\rm val},
\\
\mathcal R_{\rm DM}
&=
\nabla_\mu J^\mu_{\rm DM}.
\end{align}
A capacity-fixed, boundary-local, one-parameter \(X\) execution has
\begin{equation}
\boxed{
\mathcal R_\Lambda=\mathcal R_{\partial}=\mathcal R_X=\mathcal R_2=\mathcal R_{\rm DM}=0,
\qquad
\Omega_{2,\rm val}=0.
}
\end{equation}

\begin{theorem}[Hankel-pencil recovery under exact support hypotheses]\label{thm:hankel-hyp}
Let \(\rho\in C^{2k-1}\) on an interval and suppose
\[
\rho(a)=\sum_{j=1}^k C_j a^{s_j},
\qquad
C_j\ne0,\qquad
s_i\ne s_j\ (i\ne j).
\]
Define derivative moments \(M_n=\Dil^n\rho\) and Hankel matrices
\[
H_0=[M_{i+j}]_{i,j=0}^{k-1},
\qquad
H_1=[M_{i+j+1}]_{i,j=0}^{k-1}.
\]
Then
\[
\det H_0=
\left(\prod_{i<j}(s_j-s_i)^2\right)
\prod_{j=1}^k C_ja^{s_j}.
\]
Hence \(H_0\) is invertible exactly when the active exponents are distinct and all active amplitudes are nonzero, and
\begin{equation}
\boxed{
\frac{\det(\lambda H_0-H_1)}{\det H_0}
=
\prod_{j=1}^k(\lambda-s_j).
}
\end{equation}
\end{theorem}

Thus the active weights must be inferred freely before comparison with \(\{1,0,-1,-2,-3\}\).  In noisy applications the theorem is not a raw-differentiation prescription: covariance, model order, smoothing, and the condition number of \(H_0\) must be propagated.  A pinned \(w_X=-2/3\) calculation tests consistency, not independent identification.



\section{Complete invariant closure and conditional uniqueness}\label{sec:complete-closure}

The previous sections define an invariant sieve.  This section records the
precise sense in which the sieve gives an algebraically forced solution.
The statement is conditional: exact uniqueness follows after a finite
alphabet, a curvature split, a capacity value, and an exact expansion
function have been declared.  With noisy data it becomes a constrained
statistical reconstruction and falsification test, not a proof that nature
has been solved without uncertainty.

\begin{definition}[Core invariant closure vector]\label{def:closure-vector}
For a declared homogeneous support and a declared perturbative lift, define
the core residual vector
\begin{equation}
\boxed{
\vec{\mathcal I}_{\rm core}
:=
\begin{pmatrix}
\mathcal R_{\rm val}[\rho]\\
\mathcal R_{\partial}[\rho]\\
\mathcal R_\Lambda\\
\mathcal R_{\rm DM}[J]\\
\mathcal R_{\rm cap}(a)\\
\mathcal R_{2}[F]\\
\mathcal R_X[F]
\end{pmatrix}.
}
\end{equation}
Here
\begin{align}
\mathcal R_{\rm val}[\rho]
&=(\Dil-1)\Dil(\Dil+1)(\Dil+2)(\Dil+3)\rho,\\
\mathcal R_{\partial}[\rho]
&=\Dil(\Dil+1)(\Dil+2)(\Dil+3)\rho,\\
\mathcal R_\Lambda
&=\Lambda_{\rm obs}-\frac{3\pi}{\ell_P^2\Ncap},\\
\mathcal R_{\rm DM}[J]
&=\nabla_\mu J^\mu_{\rm DM},\\
\mathcal R_{\rm cap}(a)
&=\Dil\ln N_\partial(a)-3[1+w_{\rm eff}(a)],\\
\mathcal R_2[F]
&=\mathcal C_{-2}[F]-\Omega_K^{\rm geom}-\Omega_{2,\rm val},\\
\mathcal R_X[F]
&=\Dil\mathcal C_{-1}[F].
\end{align}
If a smooth subhorizon growth lift is additionally declared, the augmented
closure vector is
\begin{equation}
\boxed{
\vec{\mathcal I}_{\rm lift}
:=
\left(\vec{\mathcal I}_{\rm core},\,
\mathcal R_{\rm growth}[D;F]\right),
}
\end{equation}
where
\begin{equation}
\mathcal R_{\rm growth}[D;F]
=
\Dil^2D+
\left(2+\frac12\Dil\ln F\right)\Dil D
-\frac32\Omega_m(a)D.
\end{equation}
The growth residual is therefore not part of the source-independent
homogeneous theorem.  It is a null of the chosen perturbative lift.
\end{definition}

\begin{definition}[Complete invariant closure on a declared stratum]
A background model is completely closed on a declared invariant stratum if
\begin{equation}
\boxed{\vec{\mathcal I}_{\rm core}=0}
\end{equation}
and, when a perturbative lift is declared,
\begin{equation}
\boxed{\vec{\mathcal I}_{\rm lift}=0.}
\end{equation}
The phrase ``on a declared stratum'' means that the finite alphabet,
boundary-local restriction, curvature split, capacity value, and positivity
class have already been specified.
\end{definition}

\begin{theorem}[Exact finite-alphabet reconstruction]\label{thm:finite-reconstruction}
Let
\[
S_{\rm cos}=\{0,-1,-2,-3,-4\}
\]
and suppose the exact expansion function has the form
\[
F(a)=\sum_{s\in S_{\rm cos}}\Omega_s a^s
\]
on an interval.  Given exact values \(F(a_i)\) at five distinct positive
scale factors \(a_i\), the coefficient vector
\[
\vec\Omega=(\Omega_0,\Omega_{-1},\Omega_{-2},\Omega_{-3},\Omega_{-4})
\]
is unique.
\end{theorem}

\begin{proof}
Put \(x_i=a_i^{-1}\).  Then
\[
F(a_i)=\Omega_0+\Omega_{-1}x_i+\Omega_{-2}x_i^2
+\Omega_{-3}x_i^3+\Omega_{-4}x_i^4.
\]
The reconstruction matrix is the ordinary Vandermonde matrix
\[
V_{ij}=x_i^{\,j},\qquad i=1,\ldots,5,\quad j=0,\ldots,4.
\]
Since the \(a_i\) are distinct and positive, the \(x_i\) are distinct, and
\[
\det V=\prod_{i<j}(x_j-x_i)\ne0.
\]
Thus \(V\vec\Omega=\vec F\) has exactly one solution.
\end{proof}

\begin{corollary}[Projector form of the same solution]\label{cor:projector-unique}
If \(F\) is differentiable enough and lies in the finite alphabet
\(S_{\rm cos}\), the unique coefficients are the exact dilation projectors
\begin{equation}
\boxed{
\Omega_s=\mathcal C_s[F]
=
a^{-s}
\prod_{\substack{t\in S_{\rm cos}\\t\ne s}}
\frac{\Dil-t}{s-t}F(a).
}
\end{equation}
Thus the sample Vandermonde reconstruction and the local differential
projectors are the same finite-basis statement in two representations.
\end{corollary}

\begin{criterion}[Capacity, positivity, and curvature as compatibility tests]
The finite-basis solution of \cref{thm:finite-reconstruction} satisfies the
capacity-fixed positive boundary model if and only if
\begin{align}
\Omega_0&=\Omega_\Lambda^{\rm cap}
=\frac{\Lambda_{\rm cap}c^2}{3H_0^2},\\
\Omega_s&\ge0
\qquad (s\in S_{\rm cos}\ {\rm declared\ physical}),\\
\Omega_{2,\rm val}
&=\Omega_{-2}-\Omega_K^{\rm geom}\ge0 .
\end{align}
The last inequality is imposed only when the physical \(s=-2\) branch is
declared positive.
If any of these tests fails, the declared stratum is falsified.  The failure
does not destroy the finite reconstruction theorem; it says that the data
require either signed effective weights, interacting branch exchange, an
enlarged support, or a different perturbative/source lift.
\end{criterion}

\begin{theorem}[Constrained variational uniqueness]\label{thm:variational-unique}
Let \(y_i=F_{\rm obs}(a_i)\), let \(V_{is}=a_i^s\) for \(s\in S_{\rm cos}\),
and let \(W\) be a positive-definite inverse covariance matrix.  Define the
closed feasible set
\[
\mathcal K
=
\left\{
\vec\Omega:\ 
\Omega_0=\Omega_\Lambda^{\rm cap},\ 
\Omega_s\ge0,\ 
\Omega_{-2}-\Omega_K^{\rm geom}\ge0
\right\}.
\]
The constrained quadratic functional
\begin{equation}
\boxed{
\mathcal J(\vec\Omega)
=
(y-V\vec\Omega)^T W (y-V\vec\Omega)
+\iota_{\mathcal K}(\vec\Omega)
}
\end{equation}
has at most one minimizer whenever \(V\) has full column rank on the free
coordinates of \(\mathcal K\).  If the minimizer exists, it is the unique
statistical analogue of the algebraic solution on the declared stratum.
\end{theorem}

\begin{proof}
The first term is a strictly convex quadratic on the free coordinate
subspace when \(V\) has full column rank and \(W\succ0\).  The indicator
\(\iota_{\mathcal K}\) of a closed convex set preserves convexity and
restricts the domain.  A strictly convex function on a convex domain has at
most one minimizer.  Existence follows when the feasible set is nonempty and
the quadratic is coercive on the free coordinates.
\end{proof}

\begin{remark}[Noiseless uniqueness versus empirical superiority]
Exact uniqueness is a theorem about a declared finite alphabet and exact
function values.  Real BAO, CMB, supernova, lensing, and growth measurements
come with covariance, calibration choices, and possible model-order
uncertainty.  Therefore the empirical question is not whether the algebra
has a unique coefficient vector once the alphabet is fixed; it is whether
free-spectrum reconstruction and joint likelihoods select that alphabet
over alternatives.  This is why the DESI DR2 audit in \cref{sec:real-data-audit}
is treated as an illustrative background test rather than a proof of
superiority.
\end{remark}

\statementbox{The complete closure theorem is conditional but sharp: once the invariant
stratum is declared, the coefficients are algebraically forced; if the
capacity, positivity, curvature, or lift residuals fail, that stratum is
falsified rather than retuned.}



\section{Deep completion layer: spectral capacity, covariant exchange, free-spectrum diagnostics, and perturbative lift}
\label{sec:deep-completion}

This section records the additional closure layer obtained after the real-data audit and the frontier analysis.  It does not claim a final microscopic quantum-gravity derivation.  It sharpens the remaining tasks into explicit spectral, covariant, and perturbative objects.

\begin{definition}[Spectral horizon-capacity count]
Let \(R_\Lambda=\sqrt{3/\Lambda_{\rm eff}}\) be the de Sitter horizon radius associated with an effective dark-energy scale.  On the two-sphere horizon cross-section, impose a scalar spherical-harmonic band limit \(L\).  The cumulative degeneracy is
\begin{equation}
D_L=\sum_{\ell=0}^{L}(2\ell+1)=(L+1)^2 .
\end{equation}
The spectral-capacity model defines the logarithmic boundary capacity by
\begin{equation}
\boxed{
\Ncap=D_L,\qquad
L+1=\sqrt{\Ncap}=\sqrt{\pi}\frac{R_\Lambda}{\ell_P}.
}
\end{equation}
Equivalently,
\begin{equation}
\boxed{
\Lambda_{\rm eff}
=
\frac{3\pi}{\ell_P^2\Ncap}.
}
\end{equation}
The ordinary Hilbert-space dimension is therefore \(\dim\Hbd\sim\exp(\Ncap)\), not \(\Ncap\).
\end{definition}

\begin{remark}[Status of the spectral count]
This is a concrete entropy/spectral realization of the capacity postulate, not a derivation of the integer band limit from a microscopic Hamiltonian.  With \(H_0=67.4\,{\rm km\,s^{-1}\,Mpc^{-1}}\), \(\Omega_\Lambda=0.685\), and \(\ell_P=1.616255\times10^{-35}{\rm m}\), the calculation gives
\begin{equation}
\Lambda_0=1.091\times10^{-52}{\rm m}^{-2},\qquad
\Ncap=3.31\times10^{122},
\end{equation}
and hence
\begin{equation}
L+1=1.82\times10^{61}.
\end{equation}
The missing microscopic problem is now the derivation of this boundary band limit, or equivalently this code-subspace entropy.  The construction uses the standard horizon entropy normalization and the holographic state-counting viewpoint \cite{GibbonsHawking1977,Bousso2002}.
\end{remark}


\begin{definition}[\texorpdfstring{\({\rm spin}^c\)}{spin-c} index capacity on the horizon]
\label{def:spinc-index-capacity}
Let the causal horizon cross-section be \(S^2_{R_\Lambda}\simeq\mathbb{CP}^1\).  Let
\(\mathcal L_k\to S^2\) be a complex line bundle of degree \(k\ge0\), with curvature
two-form \(F_\partial\) normalized by
\begin{equation}
\frac{1}{2\pi}\int_{S^2}F_\partial=k .
\end{equation}
Equivalently, put \(q:=k+1\).  Riemann--Roch on \(\mathbb{CP}^1\), or the
index theorem for the associated \({\rm spin}^c\) Dirac/\(\bar\partial\) operator,
gives \cite{AtiyahSinger1963,Abrikosov2002}
\begin{equation}
\boxed{
\dim H^0(\mathbb{CP}^1,\mathcal O(q-1))
=
\operatorname{ind}\bar\partial_{\mathcal O(q-1)}
=
q .
}
\end{equation}
The index-capacity model defines the logarithmic horizon capacity by
\begin{equation}
\boxed{
\Ncap:=q .
}
\end{equation}
The ordinary many-state boundary Hilbert space is then a code built over
these \(q\) index cells.  If each cell carries a finite local code
\(\mathcal K\), then
\[
\log\dim(\mathcal K^{\otimes q})=q\log\dim\mathcal K .
\]
The convention \(\Ncap=q\) is the one-nat normalization of this logarithmic
capacity.  Thus \(\Ncap\) is an entropy/index count, not the ordinary
Hilbert-space dimension itself.
\end{definition}

\begin{theorem}[Index capacity, area quantization, and de Sitter levels]
\label{thm:index-capacity}
Assume the Gibbons--Hawking entropy capacity is saturated by the index
capacity of \cref{def:spinc-index-capacity}:
\begin{equation}
\Ncap=q=\frac{A_H}{4\ell_P^2}.
\end{equation}
Then the horizon area, radius, and capacity-fixed cosmological scale are
quantized as
\begin{equation}
\boxed{
A_q=4\ell_P^2q,\qquad
R_q=\sqrt{\frac{q}{\pi}}\,\ell_P,\qquad
\Lambda_q=\frac{3\pi}{q\,\ell_P^2}.
}
\end{equation}
Moreover, the scalar spectral cutoff model and the index model coincide
when
\begin{equation}
\boxed{
q=(L+1)^2.
}
\end{equation}
Thus the band limit \(L\) of the horizon spectral model is equivalently
the square-root of the \({\rm spin}^c\) index capacity.
\end{theorem}

\begin{proof}
For a de Sitter horizon in four spacetime dimensions,
\(A_H=4\pi R_\Lambda^2=12\pi/\Lambda\).  Saturating
\(\Ncap=A_H/(4\ell_P^2)\) with \(\Ncap=q\) gives
\(A_q=4\ell_P^2q\) and \(R_q=\sqrt{q/\pi}\ell_P\).  Solving
\(q=3\pi/(\Lambda\ell_P^2)\) gives the displayed \(\Lambda_q\).
The scalar spectral model has \(\Ncap=\operatorname{Tr}\Pi_L=(L+1)^2\);
identifying both capacities gives \(q=(L+1)^2\).
\end{proof}

\begin{proposition}[Rotational uniqueness of the horizon flux]
\label{prop:rotational-flux-uniqueness}
On the round horizon \(S^2\), every \(SO(3)\)-invariant curvature two-form
is proportional to the area form.  Imposing the flux quantization
\[
(2\pi)^{-1}\int_{S^2}F_\partial=q
\]
therefore fixes it uniquely:
\begin{equation}
\boxed{
F_\partial=\frac{q}{2}\sin\theta\,\dd\theta\wedge\dd\phi .
}
\end{equation}
Consequently, once the entropy normalization \(q=A_H/(4\ell_P^2)\) is
imposed, there is no remaining continuous horizon-capacity parameter.
\end{proposition}

\begin{proof}
An \(SO(3)\)-invariant two-form on \(S^2\) must be a constant multiple of
the invariant volume form \(\sin\theta\,\dd\theta\wedge\dd\phi\).  Since
\[
\int_{S^2}\sin\theta\,\dd\theta\wedge\dd\phi=4\pi,
\]
the condition \((2\pi)^{-1}\int F_\partial=q\) fixes the coefficient to
\(q/2\).
\end{proof}

\begin{remark}[What the index lift does and does not prove]
The index-capacity lift replaces the earlier symbolic
\(\Capacity\Hbd\) by a precise integer object:
\[
\boxed{
\Ncap=\operatorname{ind}\bar\partial_{\mathcal O(q-1)}=q .
}
\]
It also explains why the spectral count \((L+1)^2\) is natural on a
spherical horizon: it is the same integer capacity when \(q=(L+1)^2\), matching the familiar monopole/lowest-Landau-level origin of fuzzy-sphere matrix geometry \cite{Hasebe2010}.
What remains open is the deeper Hamiltonian problem: why the physical
boundary code subspace chooses the enormous present value
\(q\simeq3.31\times10^{122}\), or equivalently
\(L+1\simeq1.82\times10^{61}\).  The present section turns that problem
into an index/flux quantization problem, rather than leaving it as a
dimensionful cosmological constant.
\end{remark}


\begin{theorem}[Hopf-index bridge from spatial \texorpdfstring{\(S^3\)}{S3} to horizon capacity]
\label{thm:hopf-index-bridge}
Let \(S^3\to S^2\simeq\CP^1\) be the Hopf fibration and let
\(\mathcal O(1)\) be the associated hyperplane line bundle on \(\CP^1\).
For each integer \(q\ge1\), define the horizon code space
\begin{equation}
\boxed{
\mathcal H_{\partial,q}:=
H^0\!\left(\CP^1,\mathcal O(q-1)\right).
}
\end{equation}
Then \(\mathcal H_{\partial,q}\) is the unique irreducible \(SU(2)\)
representation of spin \(j=(q-1)/2\), and
\begin{equation}
\boxed{
\dim \mathcal H_{\partial,q}=2j+1=q.
}
\end{equation}
Equivalently, the same integer is the first Chern flux plus one,
\begin{equation}
\boxed{
q=1+\frac{1}{2\pi}\int_{S^2}F_{\partial},
\qquad
F_\partial\in2\pi\,c_1\!\left(\mathcal O(q-1)\right).
}
\end{equation}
Thus the capacity integer is simultaneously a Chern-flux datum, an
index, and an irreducible \(SU(2)\) code-space dimension.
\end{theorem}

\begin{proof}
The Hopf fibration is the unit-circle bundle of \(\mathcal O(-1)\), up to
the conventional dual sign.  Its tensor powers generate all holomorphic
line bundles \(\mathcal O(k)\) on \(\CP^1\).  Taking \(k=q-1\), the
Borel--Weil/Riemann--Roch theorem gives
\[
H^0(\CP^1,\mathcal O(q-1))\simeq {\rm Sym}^{q-1}(\mathbb C^2)^*,
\]
which is the irreducible \(SU(2)\) representation with
\(j=(q-1)/2\) and dimension \(q\).  Its curvature represents
\(2\pi c_1(\mathcal O(q-1))\), so the flux is \(q-1\).  Hence
\(q=1+(2\pi)^{-1}\int F_\partial\).
\end{proof}

\begin{corollary}[Rotationally invariant boundary Hamiltonian]
\label{cor:rotational-boundary-hamiltonian}
On a fixed capacity sector \(\mathcal H_{\partial,q}\), every
\(SU(2)\)-invariant boundary Hamiltonian is scalar:
\begin{equation}
\boxed{
[H_{\partial,q},U_q(g)]=0\ \forall g\in SU(2)
\Rightarrow
H_{\partial,q}=E_q\,{\bf 1}_{q}.
}
\end{equation}
Consequently, a nontrivial capacity dynamics cannot occur inside one
irreducible horizon sector.  It must be an inter-sector flow
\[
q\longrightarrow q+\Delta q,
\]
or equivalently a boundary--bulk exchange process.
\end{corollary}

\begin{proof}
This is Schur's lemma applied to the irreducible representation
\(\mathcal H_{\partial,q}\).
\end{proof}

\begin{definition}[Fuzzy-horizon coordinate algebra]
Let \(J_i\) be the spin-\(j=(q-1)/2\) generators on
\(\mathcal H_{\partial,q}\), and put
\begin{equation}
\boxed{
X_i:=R_q\frac{J_i}{\sqrt{j(j+1)}} ,
\qquad
R_q=\sqrt{\frac{q}{\pi}}\,\ell_P .
}
\end{equation}
Then
\begin{equation}
\boxed{
\sum_iX_i^2=R_q^2{\bf 1},
\qquad
[X_i,X_j]=
i\,\frac{R_q}{\sqrt{j(j+1)}}\,\epsilon_{ijk}X_k .
}
\end{equation}
Thus the same capacity integer \(q\) defines a finite noncommutative
horizon algebra whose commutative limit is \(S^2_{R_q}\).
\end{definition}

\begin{remark}[Why this is a sharper microscopic target]
The previous capacity postulate required a number \(N_\partial\).  The
Hopf-index lift sharpens the target to an integer representation label:
\[
\boxed{
N_\partial=q=2j+1.
}
\]
The present cosmological value corresponds to
\[
q\simeq3.31\times10^{122},\qquad
j\simeq1.65\times10^{122}.
\]
The remaining microscopic problem is no longer the vague request to
``explain a small \(\Lambda\)''.  It is the precise request to derive the
selected Hopf/Dirac-monopole charge \(q\), or the equivalent spin label
\(j\), from the boundary Hamiltonian and the allowed code subspace.  This
is the next nontrivial theorem-level target.  The geometric-quantization
reading is standard: quantizing \(\CP^1\) with \(\mathcal O(k)\) gives a
finite \(SU(2)\) representation and the fuzzy sphere in the large-\(k\)
limit \cite{Hasebe2010,Nair2020,Kirwin2005}.
\end{remark}

\begin{figure}[!htbp]
\centering
\includegraphics[width=\columnwidth]{fig_hopf_index_capacity_bridge.png}
\caption{Hopf-index bridge.  The spatial topology supplies a natural
\(S^3\to S^2\) fibration; the causal-horizon \(S^2\simeq\CP^1\) carries
the line bundle \(\mathcal O(q-1)\).  Its index, Chern flux plus one,
and \(SU(2)\) code-space dimension all equal \(q=N_\partial\), fixing
\(\Lambda_q=3\pi/(q\ell_P^2)\).}
\label{fig:hopf-index-bridge}
\end{figure}


\begin{theorem}[Weyl-capacity quantization of the horizon]
\label{thm:weyl-capacity-quantization}
Let \(\Delta_{S^2}\) be the scalar Laplacian on the round horizon cross-section
\(S^2_{R_\Lambda}\), and define the invariant spectral projector
\begin{equation}
\Pi_L:=
{\bf 1}_{[0,L(L+1)]}\!\left(-R_\Lambda^2\Delta_{S^2}\right),
\qquad L\in\mathbb N .
\end{equation}
Then
\begin{equation}
\boxed{
\operatorname{Tr}\Pi_L=(L+1)^2 .
}
\end{equation}
If the entropy capacity is identified with this cutoff trace,
\(\Ncap=\operatorname{Tr}\Pi_L\), the allowed capacity-fixed de Sitter
values form the discrete spectrum
\begin{equation}
\boxed{
\Lambda_L=
\frac{3\pi}{\ell_P^2(L+1)^2},
\qquad
R_L=\frac{L+1}{\sqrt{\pi}}\,\ell_P .
}
\end{equation}
The relative spacing of adjacent large-\(L\) capacity levels is
\begin{equation}
\boxed{
\frac{\Lambda_{L+1}-\Lambda_L}{\Lambda_L}
=
-\frac{2}{L+2}+O(L^{-2}).
}
\end{equation}
For the present horizon value quoted above, this spacing is of order
\(10^{-61}\), far below current cosmological precision.
\end{theorem}

\begin{proof}
The eigenvalues of \(-R_\Lambda^2\Delta_{S^2}\) are
\(\ell(\ell+1)\) with degeneracy \(2\ell+1\), \(\ell=0,1,\ldots\).
The projector \(\Pi_L\) therefore counts all modes with \(\ell\le L\):
\[
\operatorname{Tr}\Pi_L
=
\sum_{\ell=0}^{L}(2\ell+1)
=
(L+1)^2 .
\]
Substitution of \(\Ncap=(L+1)^2\) into
\(\Lambda_{\rm cap}=3\pi/(\ell_P^2\Ncap)\) gives the displayed spectrum.
The spacing formula follows by expanding
\[
\Lambda_{L+1}/\Lambda_L=\frac{(L+1)^2}{(L+2)^2}.
\]
\end{proof}

\begin{corollary}[Nearest capacity level and falsification residual]
\label{cor:nearest-capacity-level}
Given an observed late-time effective scale \(\Lambda_{\rm obs}>0\), the
nearest spectral-capacity level is obtained by
\begin{equation}
\boxed{
L_{\rm near}
=
\operatorname{round}\!\left[
\sqrt{\frac{3\pi}{\Lambda_{\rm obs}\ell_P^2}}-1
\right].
}
\end{equation}
The exact discreteness residual is
\begin{equation}
\boxed{
\mathcal R_{\rm level}
:=
\frac{\Lambda_{\rm obs}-\Lambda_{L_{\rm near}}}{\Lambda_{\rm obs}} .
}
\end{equation}
For \(L_{\rm near}\sim10^{61}\), the unavoidable level uncertainty is
Planck-suppressed by \(O(10^{-61})\).  Hence capacity quantization is a
microscopic constraint rather than a presently resolvable late-time
cosmological oscillation.
\end{corollary}

\begin{remark}[Two boundaries, two different invariants]
The horizon capacity count lives on the causal horizon cross-section
\(S^2\).  The valuation annihilator lives on the spatial round filling
\(B^4/S^3\).  These are not the same boundary.  Their connection in this
paper is through the closed-FLRW homogeneous geometry and the capacity
postulate, not through an identity \(S^2=S^3\).  This separation prevents
double counting of horizon entropy and homogeneous valuation branches.
\end{remark}

\begin{theorem}[Covariant capacity-exchange law]
\label{thm:covariant-exchange-law}
Let
\begin{equation}
T^{\rm cap}_{\mu\nu}:=-\frac{\Lambda_{\rm eff}(x)}{8\pi G}g_{\mu\nu}.
\end{equation}
If the total stress tensor \(T^{\rm tot}_{\mu\nu}=T^{\rm bulk}_{\mu\nu}+T^{\rm cap}_{\mu\nu}\) is covariantly conserved, then the bulk sector obeys
\begin{equation}
\boxed{
\nabla_\mu T_{\rm bulk}^{\mu}{}_{\nu}
=
\frac{1}{8\pi G}\nabla_\nu\Lambda_{\rm eff}
=:Q_\nu^{\partial\to{\rm bulk}} .
}
\end{equation}
For a homogeneous FLRW history this is equivalent to
\begin{equation}
\boxed{
\Dil\ln\Ncap(a)=3[1+w_{\rm eff}(a)],
}
\end{equation}
provided
\[
\Lambda_{\rm eff}(a)=\frac{3\pi}{\ell_P^2\Ncap(a)}
\]
is interpreted as the effective dark-energy density scale, not as the constant Bianchi integration mode.
\end{theorem}

\begin{proof}
Since \(\nabla_\mu g^{\mu}{}_{\nu}=0\),
\[
\nabla_\mu T_{\rm cap}^{\mu}{}_{\nu}
=
-\frac{1}{8\pi G}\nabla_\nu\Lambda_{\rm eff}.
\]
The condition \(\nabla_\mu T_{\rm tot}^{\mu}{}_{\nu}=0\) gives the displayed exchange law.  In FLRW, \(\rho_{\rm eff}=\Lambda_{\rm eff}/(8\pi G)\) satisfies
\[
\Dil\ln\rho_{\rm eff}=-3(1+w_{\rm eff}).
\]
Because \(\Ncap\propto\Lambda_{\rm eff}^{-1}\), the stated capacity-flow equation follows.
\end{proof}

\begin{criterion}[Boundary-exchange lift]
A dynamical dark-energy reconstruction must specify \(Q_\nu^{\partial\to{\rm bulk}}\).  Without this covariant exchange current, \(\Ncap(a)\) is only a diagnostic dictionary.  With it, capacity flow becomes a conservation law for the enlarged bulk-plus-boundary system.
\end{criterion}

\begin{theorem}[Observable two-root exchange current, executed form]
For the interacting pair \((s_-,s_+)=(-1,1)\), define
\[
\kappa_1=-3(1+w),\qquad
\kappa_2=(1-\kappa_1)(1+\kappa_1).
\]
The capacity-exchange current is recovered from \((w,\Dil w)\) by
\begin{equation}
\boxed{
J_+
=
\frac{-3\Dil w-\kappa_2}{2}.
}
\end{equation}
Thus a DESI-type quadrant with \(\Dil w>0\), when it lies in \(-1<\kappa_1<1\), requires \(J_+<-\kappa_2/2<0\).  The sign is the invariant statement that weight flows from the upper branch to the lower branch while positivity is maintained.
\end{theorem}

\begin{figure}[!htbp]
\centering
\includegraphics[width=\columnwidth]{fig_interacting_capacity_current.png}
\caption{Observable interacting capacity current \(J_+\) reconstructed from the DESI DR2 \(w_0w_a\) curves used in the audit.  The negative sign in the admissible interval is the invariant signature that a positive two-root representation needs exchange; a noninteracting positive mixture cannot produce \(\Dil w>0\).}
\label{fig:interacting-current}
\end{figure}


\begin{proposition}[Executed free-spectrum smoke test]
For the exact four-branch test profile
\begin{align}
F_{\rm test}(a)={}&0.541216+0.174541a^{-1}
+0.284243a^{-3}
\nonumber\\
&+9.2\times10^{-5}a^{-4}.
\end{align}
the logarithmically sampled Hankel pencil with \(k=4\) recovers
\[
s=(-4.0000,-3.0000,-1.0000,0.0000).
\]
up to numerical roundoff.  Applying the same local pencil to the DESI \(w_0w_a\) capacity-flow curves does not return a stable subset of \(\{1,0,-1,-2,-3\}\).  Hence present \(w_0w_a\) curves are evidence for a dynamical diagnostic trend, not a free-spectrum identification of the valuation roots.
\end{proposition}

\begin{criterion}[GitHub data interface for the next audit]
The next numerical layer shall use public data products rather than hand-entered summaries.  The intended inputs are the DESI DR2 BAO data in the Cobaya BAO repository, Pantheon+SH0ES distance and covariance files, compressed Planck likelihood products, and public growth compilations \cite{CobayaBAOData,PantheonPlusDataRelease,PlanckCompressedGitHub,GrowthDataGitHub}.  The acceptance condition remains unchanged: infer \(s_j\) freely before comparing with \(S_{\rm val}\).
\end{criterion}

\begin{definition}[Gauge-invariant perturbative valuation lift]
\label{def:gi-perturbative-valuation-lift}
Let \(K\) denote the compact \(S^3\) harmonic label.  After either solving the
linearized constraints or keeping them as constraint components, collect the
gauge-invariant scalar, vector, tensor, and matter amplitudes in
\begin{equation}
\boxed{
\mathbf X_K
=
(\Phi_K,\Psi_K,V^{\rm T}_K,h^{\rm TT}_K,
\delta_{\alpha K},\theta_{\alpha K},\ldots).
}
\end{equation}
Here \(\Phi,\Psi\) denote scalar gauge-invariant metric potentials,
\(V^{\rm T}\) transverse vector amplitudes, \(h^{\rm TT}\) transverse-traceless
tensor amplitudes, and \(\delta_{\alpha},\theta_\alpha\) matter-source
perturbations.  Define the gauge-invariant dilation evolution operator by
\begin{equation}
\boxed{
\Dil_{\rm GI}\mathbf X_K
=
\mathsf A_K(a)\mathbf X_K+\mathbf J_K(a),
\qquad
\Dil:=a\frac{\dd}{\dd a}.
}
\end{equation}
The perturbative valuation lift is the spectral decomposition of
\(\mathsf A_K\) relative to the background valuation roots
\[
S_{\rm val}=\{1,0,-1,-2,-3\}.
\]
For an isolated root \(s\), define the perturbed valuation projector by the
Riesz formula
\begin{equation}
\boxed{
\mathsf P^{\rm pert}_{s,K}
=
\frac{1}{2\pi i}
\oint_{\Gamma_s}
(z-\mathsf A_K)^{-1}\,\dd z ,
}
\end{equation}
where \(\Gamma_s\) encloses \(s\) and no other background root.
\end{definition}

\begin{criterion}[Gauge-invariant perturbative closure]
\label{crit:gi-perturbative-closure}
The homogeneous valuation closure has a successful perturbative lift only if
the compact-harmonic mode system admits a controlled decomposition
\begin{equation}
\boxed{
\mathsf A_K(a)=\mathsf S+\mathsf M_K(a),
\qquad
\mathsf S=\operatorname{diag}(1,0,-1,-2,-3),
}
\end{equation}
on the declared gauge-invariant perturbation space.  With the background
projectors \(\mathsf P_s^0\), define the mixing blocks
\begin{equation}
\boxed{
\mathsf M_{ss'}(K,a)
:=
\mathsf P_s^0\,\mathsf M_K(a)\,\mathsf P_{s'}^0 .
}
\end{equation}
The lift is valuation-diagonal only when
\[
\mathsf M_{ss'}(K,a)=0
\qquad (s\ne s').
\]
Otherwise the theory predicts branch mixing.  Sound speed, anisotropic stress,
momentum exchange, compact \(S^3\) harmonic effects, and tensor propagation
belong to the matrix \(\mathsf M_K\), not to the homogeneous polynomial
\(P_{\rm val}(\Dil)\) alone.
\end{criterion}

\begin{proposition}[Perfect-fluid Newtonian-gauge reduction]
\label{prop:perfect-fluid-reduction}
A noninteracting perfect-fluid source with constant
\(w_s=-1-s/3\) gives, in Newtonian gauge, the familiar scalar subsystem
\cite{MaBertschinger1995}
\begin{align}
\delta'={}&-(1+w)(\theta-3\Phi')
-3\mathcal H(c_s^2-w)\delta
\nonumber\\
&-9\mathcal H^2(1+w)(c_s^2-c_a^2)\frac{\theta}{k^2},
\\
\theta'={}&-\mathcal H(1-3c_s^2)\theta
+\frac{k^2c_s^2}{1+w}\delta
+k^2\Psi-k^2\sigma .
\end{align}
For the boundary-local branch \(s=-1\),
\begin{equation}
\boxed{
w=-\frac23,\qquad c_a^2=w .
}
\end{equation}
This system is a reduction of \cref{def:gi-perturbative-valuation-lift}, not
the definition of the perturbative lift.  A smooth canonical realization sets
\(c_s^2=1\) and \(\sigma=0\); a clustered or interacting realization requires
the corresponding entries of \(\mathsf M_K\) and \(\mathbf J_K\).
\end{proposition}

\begin{criterion}[Perturbative superiority test]
A background success is insufficient.  To claim empirical superiority over
\(\Lambda{\rm CDM}\) or \(w_0w_a{\rm CDM}\), a lift must predict the entries of
\(\mathsf M_K\) and \(\mathbf J_K\): sound speed, anisotropic stress, momentum
transfer, compact \(S^3\) harmonic corrections, tensor propagation, and the
coupled evolution of \(\Phi,\Psi,\delta,\theta,V^{\rm T},h^{\rm TT}\).  BAO
tests the background; CMB, lensing, and growth test the source.
\end{criterion}



\section{Microscopic capacity selection: no-go theorem and surviving selectors}
\label{sec:capacity-selection}

The previous sections turn the cosmological constant into a horizon-capacity integer \(q=\Ncap\).  This is a genuine narrowing, but it is not yet a selection theorem.  The remaining problem is not to write another expression for \(q\); it is to explain why the boundary theory occupies the observed superselection sector
\[
q_{\rm obs}\simeq 3.31\times 10^{122}.
\]

\begin{nogotheorem}[Topology and index quantize capacity but do not select its observed value]
\label{nogo:q-selection}
The data
\[
S^2\simeq \CP^1,\qquad
\mathcal H_{\partial,q}=H^0(\CP^1,\mathcal O(q-1)),
\qquad
\Ncap=q
\]
quantize the allowed horizon capacities, but they do not select a unique value of \(q\).  In particular, topology, \(SU(2)\) covariance, the spin-\(c\) index, and the fuzzy-sphere representation theory alone do not force \(q\simeq 3.31\times10^{122}\).
\end{nogotheorem}

\begin{proof}
Holomorphic line bundles over \(\CP^1\) are classified by their first Chern class \(m\in\mathbb Z\).  For \(m=q-1\ge0\),
\[
\operatorname{ind}\bar\partial_{\mathcal O(m)}
=\dim H^0(\CP^1,\mathcal O(m))=m+1=q.
\]
Thus the construction permits one admissible irreducible \(SU(2)\) code sector for every positive integer \(q\).  The symmetry algebra fixes the representation once \(q\) is chosen, but it does not prefer one integer over another.  Therefore any selection of the observed large integer requires an additional dynamical, statistical, or boundary-condition principle.
\end{proof}

\begin{definition}[Boundary selector functional]
A microscopic completion selects a capacity sector if it supplies a dimensionless functional
\[
\Gamma_\partial(q;\lambda)
=
-\ln Z_\partial(q;\lambda),
\]
where \(\lambda\) denotes all dimensionless microscopic couplings and boundary conditions, such that the physical sector satisfies the discrete stationarity condition
\[
\boxed{
\Delta_q\Gamma_\partial(q;\lambda)\big|_{q=q_*}=0
}
\]
together with stability
\[
\Delta_q^2\Gamma_\partial(q_*;\lambda)>0
\]
or, for a microcanonical construction, an equivalent extremal entropy condition.  The observed cosmological constant is then
\[
\Lambda_*=\frac{3\pi}{\ell_P^2q_*}.
\]
\end{definition}

\begin{criterion}[What would count as a derivation of \(\Lambda\)]
The amplitude problem is solved only when a microscopic boundary theory computes \(q_*\) without using \(\Lambda_{\rm obs}\) as an input.  Equivalently,
\[
q_*=\arg\min_q\Gamma_\partial(q;\lambda)
\]
or the corresponding index/entropy extremum must yield
\[
q_*\simeq 3.31\times10^{122}
\]
from independently specified dimensionless data.
\end{criterion}

\begin{widetext}
\begin{center}
\footnotesize
\renewcommand{\arraystretch}{0.92}
\begin{tabular}{@{}p{0.18\linewidth}p{0.29\linewidth}p{0.25\linewidth}p{0.20\linewidth}@{}}
\toprule
Candidate selector & Invariant output & What it achieves & Remaining obstruction\\
\midrule
Spin-\(c\) index on \(S^2\) &
\(q=\operatorname{ind}\bar\partial_{\mathcal O(q-1)}\) &
Quantizes capacity as an integer and realizes the horizon code as a single \(SU(2)\) irrep &
Every \(q\in\mathbb N\) is allowed; no sector selection\\
Spectral cutoff / fuzzy sphere &
\(\operatorname{Tr}\Pi_L=(L+1)^2=q\) &
Gives a finite noncommutative horizon algebra and angular resolution scale &
The cutoff \(L\) remains a sector label\\
Loop-quantum-gravity horizon counting &
Area-proportional entropy from horizon degrees of freedom &
Provides an independent microscopic area-counting paradigm &
Requires choice/renormalization of microscopic parameters and horizon data\\
Causal-set fluctuation picture &
Ever-present \(\Lambda\) of order inverse spacetime volume &
Explains the scale of a small fluctuating term qualitatively &
Gives stochastic dynamics, not the fixed de Sitter capacity integer\\
RG or dimensional transmutation &
Large numbers from exponentials of inverse couplings &
Can generate \(10^{122}\)-size hierarchies from dimensionless inputs &
Requires the correct boundary coupling and beta function\\
\bottomrule
\end{tabular}
\end{center}
\end{widetext}

\begin{proposition}[Minimal remaining microscopic data]
A complete capacity derivation must provide at least one of the following equivalent selectors:
\begin{enumerate}
\item a boundary Hamiltonian spectrum \(E_q\) and ensemble rule selecting \(q_*\);
\item a flux-sector action whose discrete saddle occurs at \(q_*\);
\item an RG transmutation law producing \(q_*\) from independently measured dimensionless couplings;
\item a dynamical boundary-exchange law whose late-time attractor is \(q_*\).
\end{enumerate}
Without one of these, the construction remains an exact invariant sieve rather than a full microscopic derivation of \(\Lambda\).
\end{proposition}

\begin{remark}[Relation to existing quantum-gravity directions]
Loop-quantum-gravity black-hole entropy, causal-set ``ever-present'' \(\Lambda\), causal dynamical triangulations, functional renormalization group approaches, and fuzzy-sphere matrix geometries each touch one side of the selector problem \cite{Rovelli1996,Ashtekar1998,Jacobson2007,DiazPoloPranzetti2011,BarberoPranzetti2022,Sorkin2007,Zwane2017,Loll2020,Saueressig2023,Hasebe2025}.  None of these ingredients alone selects the observed de Sitter capacity sector inside the present invariant construction.  Their value here is diagnostic: they show which kind of missing microscopic datum would be sufficient.
\end{remark}



\section{Corpus-assisted selector audit and surviving boundary functional}
\label{sec:corpus-selector}

The supplied FBS3R discovery corpus was audited as an internal source of possible
amplitude selectors.  Four structures survive the audit as invariant ingredients,
but none by itself fixes the observed capacity sector:
\begin{equation}
\boxed{
n_\varphi(R),\qquad
\Delta\theta_{\rm BO},\qquad
Q_\xi[\partial U],\qquad
\mathcal B_{\rm P}.
}
\end{equation}
Here \(n_\varphi(R)=\log_\varphi(R/\ell_P)\) is a dimensionless scale
coordinate, not a force or coupling; \(\Delta\theta_{\rm BO}\) is the
Born--Oppenheimer/CMC phase drift; \(Q_\xi[\partial U]\) is a quasilocal
observer-domain charge; and \(\mathcal B_{\rm P}\) is the Planck-barrier
admissibility condition.  This is consistent with the topological decoupling
principle: topology and index data quantize admissible sectors, but amplitude
selection requires additional non-topological dynamical information.

\begin{definition}[\(\varphi\)-shell capacity diagnostic]
For a horizon-capacity sector \(q\), define
\begin{align}
n_\partial(q)&:=\log_\varphi\sqrt{\frac{q}{\pi}},\\
\mathcal R_\varphi(q;N)&:=n_\partial(q)-N,\qquad N\in\mathbb Z .
\end{align}
The corresponding \(\varphi\)-shell candidate is
\begin{equation}
q_\varphi(N):=\pi\varphi^{2N}.
\end{equation}
This is only a shell diagnostic unless an independent variational principle
selects \(N\).
\end{definition}

For the Planck-2018/CODATA normalization used in this manuscript,
\begin{equation}
q_{\rm obs}=3.3072\times10^{122},
\qquad
n_\partial(q_{\rm obs})=291.936672 .
\end{equation}
The nearest integer shell is \(N=292\), for which
\begin{align}
q_\varphi(292)&=3.5151\times10^{122},\\
\epsilon_\varphi&:=
\frac{q_\varphi(292)-q_{\rm obs}}{q_{\rm obs}}
=6.2844\times10^{-2}.
\end{align}
Thus the \(\varphi\)-ladder is numerically close enough to be a useful
dimensionless registry, but not exact enough to be the missing selector.

\begin{proposition}[\(\varphi\)-shell residual audit]\label{prop:phi-residual-audit}
Let \(N=\operatorname{round} n_\partial(q_{\rm obs})=292\) and
\(\delta_\varphi:=N-n_\partial(q_{\rm obs})\).  Then
\begin{equation}
1+\epsilon_\varphi=\varphi^{2\delta_\varphi},
\qquad
\delta_\varphi=6.3328\times10^{-2}.
\end{equation}
For the Planck-2018/CODATA normalization used here,
\begin{equation}
\epsilon_\varphi=0.0628440,
\qquad
\frac{q_{\rm obs}}{q_\varphi(292)}=0.940872 .
\end{equation}
The numerical proximity
\begin{equation}
\epsilon_\varphi-\frac{\pi}{50}=1.22\times10^{-5}
\end{equation}
is a registry coincidence unless it follows from a boundary action.  At fixed
\(H_0\), the nearest exact \(\varphi\)-shell would lower the inferred
\(\Lambda\) by the factor \(0.940872\), which is far too large to be hidden
inside the quoted precision of the Planck-anchored background.  Therefore the
\(N=292\) shell is an important clue, but not a derivation.
\end{proposition}

\begin{center}
\includegraphics[width=\columnwidth]{fig_phi_residual_audit.png}

\refstepcounter{figure}\label{fig:phi-residual-audit}
{\small\textbf{Figure~\thefigure.} Audit of the nearest \(\varphi\)-shell residual.  The \(6.28\%\)
offset is close to simple dimensionless numbers such as \(\pi/50\), but these
matches are not invariant mechanisms.  A final theory must derive the residual
or force it to vanish from a boundary selector functional.}
\end{center}

\begin{center}
\includegraphics[width=\columnwidth]{fig_phi_shell_selector.png}

\refstepcounter{figure}\label{fig:phi-shell}
{\small\textbf{Figure~\thefigure.} \(\varphi\)-shell diagnostic for the horizon capacity.  The observed
capacity sector lies near the integer shell \(N=292\), but the residual is not
zero.  Hence the \(\varphi\)-coordinate is a registry, not an amplitude
derivation.}
\end{center}

The corpus therefore suggests the following sharpened selector ansatz:
\begin{equation}
\boxed{
q_*=\arg\min_{q\in\mathbb N}\Gamma_\partial^{\rm sel}(q),
}
\end{equation}
where
\begin{align}
\Gamma_\partial^{\rm sel}(q)
={}&
\lambda_\varphi\,\operatorname{dist}
   \!\left(n_\partial(q),\mathbb Z\right)^2
+\lambda_{\rm ph}\int_I
\left[
H_q-\frac13\dot{\Delta\theta}_q
\right]^2\dd\tau \nonumber\\
&+\lambda_Q
\left\|Q_\xi^{(q)}[\partial U]-Q_{\xi,{\rm obs}}\right\|^2_{\mathcal C^{-1}}
+\lambda_{\rm P}\,\mathcal B_{\rm P}(q).
\label{eq:corpus-selector-functional}
\end{align}
Here \(H_q^2=\Lambda_q/3=\pi/(q\ell_P^2)\), \(\mathcal C\) is the
observational covariance assigned to the quasilocal-charge reconstruction, and
\(\mathcal B_{\rm P}\) is an infinite barrier for sectors whose admissible
history violates the Planck-curvature bound.  The constants
\(\lambda_\varphi,\lambda_{\rm ph},\lambda_Q,\lambda_{\rm P}\) are not allowed
phenomenological knobs in a final theory; they must be supplied by the
microscopic boundary action.

\begin{criterion}[True selector criterion]
A microscopic completion solves the capacity-selection problem only if it
derives a boundary action or Hamiltonian whose semiclassical reduction yields
\cref{eq:corpus-selector-functional}, or a sharper replacement, and whose
integer minimizer satisfies
\begin{equation}
q_*=3.3072\times10^{122}
\end{equation}
without using the observed \(\Lambda\) as an input.  The corpus audit therefore
does not complete the \(\Lambda\) derivation, but it eliminates three false
routes: topology alone, \(\varphi\)-shell numerology alone, and bounded forcing
alone.
\end{criterion}

\begin{criterion}[Admission rules for further proposed invariants]\label{crit:admission-rules}
The corpus and the subsequent external proposals generate many possible
connections.  In the present manuscript only the following versions are
admitted.
\begin{enumerate}
\item The identity \(\chi(B^4)-\chi(S^3)=1\) is an Euler checksum for the
filling pair, not a new valuation index.  The alternating expression
\(\sum_s(-1)^s\) has no invariant meaning for dilation weights.  The fact that
the boundary-local restriction removes one bulk branch is already encoded by
\(\mathsf P_1\rho=0\).
\item The \(S^3\) Dirac spectrum is a compact-topology perturbative datum,
not the homogeneous dilation spectrum.  It may enter a future perturbative
lift, but it cannot replace \(P_{\rm val}(\mathscr D)\rho=0\)
\cite{Baer1996}.
\item Zeta regularization of the scalar Laplacian on \(S^2\) is not a positive
capacity count.  Excluding the zero mode, heat-kernel theory gives
\(\zeta_{\Delta_{S^2}}(0)=-2/3\), not a state number.  The accepted capacity
normalization remains the entropy/index count
\(q=A/(4\ell_P^2)\), while zeta data may only enter determinant corrections
\cite{Vassilevich2003}.
\item Perturbative valuation projectors are admissible only after a gauge-fixed
or gauge-invariant perturbative phase space is declared.  Acting with
\(e_s(\mathscr D)\) on background densities is exact; acting on metric
perturbations is a separate lift problem.
\item A boundary stress tensor obtained by varying a capacity functional must be
a codimension-two screen or Brown--York type object with a specified embedding.
The correct covariant statement at this stage is the exchange balance law,
not the conservation of \(N_\partial u^\mu\) when \(N_\partial(a)\) varies.
\end{enumerate}
These rules are deliberately restrictive.  They prevent numerical coincidences
or gauge-dependent constructions from being mistaken for new invariants.
\end{criterion}

\begin{remark}[What the corpus added]
The useful contribution of the corpus is not a hidden numerical answer for
\(q\).  It supplies a narrower search space: (i) use \(n_\varphi\) only as a
dimensionless scale coordinate; (ii) test whether CMC/Born--Oppenheimer phase
drift can act as a Hamiltonian selector; (iii) use quasilocal charges on
observer-domain boundaries as the measurable exchange channel; and (iv) impose
the Planck barrier as an admissibility wall.  The remaining problem is therefore
well posed: derive \(\Gamma_\partial^{\rm sel}\) rather than invent another
dark-energy parameter.
\end{remark}



\section{Rare-law invariant sweep and phase-defect selector}

This section records a negative and positive result from a wider
search through less standard invariant mechanisms.  The purpose is not
to decorate the model with numerology.  A candidate is admitted only if
it supplies a dimensionless, gauge-independent, and normalization-free
quantity that can enter a falsifiable selector for the capacity integer
\(q=N_\partial\).

\begin{criterion}[Admission test for rare-law invariants]
A proposed additional invariant \(I\) is admissible only if it satisfies
all of the following conditions:
\begin{enumerate}
\item it is dimensionless, or it fixes a dimensionless ratio after the
declared Planck normalization;
\item it is invariant under diffeomorphisms and under changes of the
scale pivot used in the dilation algebra;
\item it is not merely a regularized determinant with sign or phase
unless that sign or phase enters a stated boundary action;
\item it either quantizes \(q\), selects \(q\), or gives a covariant
exchange law for \(q(a)\).
\end{enumerate}
Numerical proximity to a familiar constant is not sufficient.
\end{criterion}

\begin{proposition}[Rare-law sweep]
Among the surveyed mechanisms, the following survive only in restricted
roles:
\begin{itemize}
\item Bousso-type entropy bounds give a capacity ceiling, not a selector
for \(q\) \cite{Bousso2002}.
\item Holographic dark energy gives an inverse-area dictionary, but its
standard form still contains an IR cutoff choice and a parameter
\cite{Li2004}.
\item Causal-set everpresent \(\Lambda\) gives a stochastic
\(V^{-1/2}\)-scale residue, not a deterministic index selector
\cite{Sorkin2007,Zwane2017}.
\item Gauss--Bonnet topology change can provide a covariant exchange
channel through topology-change density, but not a fixed \(q\) by itself
\cite{Tsilioukas2024}.
\item Heat-kernel, zeta, and analytic-torsion data give determinant and
anomaly data, not a positive horizon state count without an action
\cite{Vassilevich2003,RaySinger1971}.
\item Selberg, automorphic, and modular trace data become possible
spectral selectors only after a compact quotient and an operator are
declared.
\item Renormalization-group transmutation can generate hierarchies, but
must output the discrete boundary index \(q\), not merely a running scale
\cite{Saueressig2023}.
\item The \(\varphi\)-shell registry is admitted as a scale registry and
phase-defect clue, not as an independent law.
\end{itemize}
Thus none of these mechanisms alone replaces the boundary selector.
The surviving target remains a functional \(\Gamma_\partial(q)\) whose
minimum is the observed capacity integer without using \(\Lambda_{\rm obs}\)
as input.
\end{proposition}

\begin{theorem}[\(\varphi\)-shell phase-defect identity]
Let
\[
q_\varphi(N):=\pi\varphi^{2N},
\qquad
n_\partial(q):=\log_\varphi\sqrt{q/\pi}.
\]
For the present capacity ledger one finds
\[
n_\partial(q_{\rm obs})=291.936672\ldots,
\qquad
N=292,
\]
and hence
\[
\epsilon_\varphi:=\frac{q_\varphi(292)}{q_{\rm obs}}-1
=0.062844029\ldots .
\]
Equivalently,
\begin{align}
1+\epsilon_\varphi&=\varphi^{2\delta_\varphi},\\
\delta_\varphi&:=292-n_\partial(q_{\rm obs})
=0.063327996\ldots .
\end{align}
Moreover
\[
\epsilon_\varphi-\frac{\pi}{50}=1.2176\times10^{-5}.
\]
Therefore the closest simple phase interpretation of the residual is
\[
q_{\rm obs}
\simeq
\frac{\pi\varphi^{584}}{1+\pi/50}.
\]
This is a phase-defect clue, not a theorem selecting \(q\).  It becomes
a theorem only if the boundary selector action derives the factor
\(1+\pi/50\), or an indistinguishable correction, from a covariant
boundary phase.
\end{theorem}

\begin{proof}
The first two identities are definitions.  The numerical values follow
by direct substitution of \(q_{\rm obs}=3\pi/(\Lambda_0\ell_P^2)\) into
\(n_\partial(q)\).  The comparison with \(\pi/50\) is an arithmetic
identity at the stated precision.  No dynamical conclusion follows from
this arithmetic closeness alone.
\end{proof}

\begin{figure}[!htbp]
\includegraphics[width=\columnwidth]{fig_rare_law_phi_phase_audit.png}
\caption{Numerical audit of simple constants near the golden-shell residual
\(\epsilon_\varphi\).  The proximity to \(\pi/50\) is real, but it is
not promoted to a physical law unless a boundary phase functional derives it.}
\label{fig:rare-law-phi-audit}
\end{figure}

\begin{criterion}[Phase-corrected selector target]
The next admissible microscopic improvement is a boundary action whose
stationarity condition contains the dimensionless defect
\[
\Theta_\partial(q)
:=
\ln\frac{q_\varphi(N(q))}{q}
\]
and derives
\[
\Theta_\partial(q_*)=\ln(1+\pi/50)+O(10^{-5})
\]
without inserting \(q_*\) or \(\Lambda_{\rm obs}\).  Equivalently, the
selector may be written schematically as
\begin{align}
\Gamma_\partial^{(15)}(q)
={}&
\lambda_\varphi d_\varphi(q)^2
+\lambda_\theta
\left[
\Theta_\partial(q)-\theta_{\rm bd}
\right]^2 \nonumber\\
&+\Gamma_{\rm exch}(q)
+\Gamma_{\rm data\;nulls}.
\end{align}
where \(d_\varphi(q)\) is the distance of \(n_\partial(q)\) from an
integer and \(\theta_{\rm bd}\) must be derived from the boundary theory,
not fitted.  The manuscript therefore admits the \(\pi/50\) residual as
a target for the selector, while rejecting it as a completed derivation.
\end{criterion}



\section{Scale-selector theorem and transmutation ledger}

The rare-law sweep leaves one hard question: why the admitted boundary
capacity sector has
\[
q_{\rm obs}\simeq3.31\times10^{122}.
\]
The preceding sections show that topology, index theory, and the fuzzy
\(S^2\) algebra quantize \(q\).  They do not select it.  This section
records the next invariant obstruction and the remaining viable class of
selectors.

Define the logarithmic capacity coordinate
\begin{equation}
\eta_q:=\frac12\ln\!\left(\frac{q}{\pi}\right)
       =\ln\!\left(\frac{R_q}{\ell_P}\right),
\qquad
n_\varphi(q):=\frac{\eta_q}{\ln\varphi}.
\end{equation}
For the present ledger,
\begin{align}
\ln q_{\rm obs}&=282.111\pm0.015,\\
\eta_q&=140.48338\ldots,\\
n_\varphi&=291.93667\ldots .
\end{align}
Thus the huge number \(q\) becomes an order-\(10^2\) logarithmic
coordinate.  Any microscopic selector that uses only order-one
dimensionless data must act on this logarithmic variable, not on \(q\)
itself.

\begin{nogotheorem}[Polynomial naturalness obstruction for \(q\)-selection]
Let a proposed selector \(\Gamma_\partial(q)\) be built from finitely
many order-one dimensionless coefficients, topology/index data, and a
finite combination of powers of \(q\) and powers of \(\ln q\), with no
exponential or transmutation mechanism and with no large integer
inserted by hand.  Then an isolated stable minimum at
\(q\sim10^{122}\) is not natural: either the minimum is controlled by
a large input already comparable to \(\ln q\) or \(q\), or by a
coefficient tuned to exponential precision.  Therefore a genuine
selector for \(q_{\rm obs}\) must use one of the following structures:
\begin{enumerate}
\item an exponential hierarchy, such as dimensional transmutation or
localization;
\item a quantized large flux or large representation sector;
\item a boundary phase condition acting on \(\eta_q\) or
\(n_\varphi(q)\);
\item a data-falsifiable exchange law that dynamically drives the
system to a stable \(q\)-sector.
\end{enumerate}
\end{nogotheorem}

\begin{proof}
A finite polynomial/rational expression in \(q\) with order-one
coefficients has stationary points at order-one algebraic combinations
of those coefficients unless large coefficients or cancellations are
inserted.  Replacing \(q\) by \(\ln q\) moves the observed target from
\(10^{122}\) to \(282.1\), but a polynomial in \(\ln q\) still requires
a large dimensionless input of that order unless the variable itself is
selected by a phase, flux, or transmutation law.  Exponential maps
evade this obstruction because order-one changes in a coupling or phase
can produce exponentially large \(q\).  The statement is therefore a
naturalness obstruction, not a logical impossibility theorem.
\end{proof}

\begin{criterion}[Transmutation-compatible capacity selector]
A boundary theory gives an admissible capacity selector only if it
outputs \(q\) in one of the invariant forms
\begin{align}
q&=\exp\!\left(\frac{16\pi^2}{b\,g_\partial^2}\right)\,\Xi,
\label{eq:rg-q}\\
q&=\pi\varphi^{2N}\exp[-\Theta_{\rm bd}],\label{eq:phi-phase-q}\\
q&=\dim \mathcal C_{\rm code}(Q_{\rm DM},\Agrav,P_{\rm val}),
\label{eq:code-q}
\end{align}
where \(b g_\partial^2\), \(\Theta_{\rm bd}\), or the code data are
computed inside the boundary theory rather than fitted to
\(\Lambda_{\rm obs}\).  Equation~\eqref{eq:rg-q} is the standard
dimensional-transmutation pattern; \eqref{eq:phi-phase-q} is the
phase-corrected shell pattern found above; \eqref{eq:code-q} is the
holographic code-space form.  These are not three different
cosmological constants.  They are three mathematically possible
representations of the same integer \(q\).
\end{criterion}

For the transmutation form with \(\Xi=1\), the present ledger requires
\begin{equation}
b\,g_\partial^2=\frac{16\pi^2}{\ln q_{\rm obs}}
=0.55976\ldots .
\end{equation}
Equivalently,
\[
\alpha_\partial:=\frac{g_\partial^2}{4\pi}
=\frac{4\pi}{b\ln q_{\rm obs}}.
\]
For example, \(b=10\) would require
\(\alpha_\partial^{-1}=224.5\).  This is not a derivation; it is a
target ledger.  Its value is in the ordinary weak-coupling range, so
the transmutation route is not numerically absurd.  Dimensional
transmutation in gravity/cosmology has been studied in scale-invariant
and higher-derivative settings \cite{EinhornJones2014,Salvio2020}, and
edge-mode boundary actions provide a natural arena in which boundary
degrees of freedom can acquire their own dynamics \cite{GeillerJaiAkson2019}.
What is still missing is the specific boundary beta function or
edge-mode Hamiltonian that computes \(b\,g_\partial^2\).

\begin{center}
\includegraphics[width=\columnwidth]{fig_transmutation_selector.png}

\refstepcounter{figure}\label{fig:transmutation-selector}
{\small\textbf{Figure~\thefigure.} Capacity transmutation ledger.  If
\(q=\exp[16\pi^2/(b g_\partial^2)]\), the observed capacity requires
\(b g_\partial^2=0.55976\).  The plot shows the corresponding
\(\alpha_\partial^{-1}=4\pi/g_\partial^2\) for representative beta
coefficients \(b\).  This does not derive \(q\); it identifies the
weak-coupling target any boundary transmutation mechanism must hit.}
\end{center}

\begin{theorem}[Reduced final selector problem]
After the trace-free, index, valuation, dark-current, and observational
closures are imposed, the remaining microscopic problem is equivalent
to deriving one of the following three quantities without using
\(\Lambda_{\rm obs}\) as input:
\begin{equation}
\boxed{
\begin{gathered}
b\,g_\partial^2=0.55976\ldots,\\
\Theta_{\rm bd}=\ln(1+\pi/50)+O(10^{-5}),\\
\dim\mathcal C_{\rm code}=q_{\rm obs}.
\end{gathered}}
\end{equation}
Any successful derivation of one of these three, together with the
already proved equivalences
\[
\Lambda_q=\frac{3\pi}{q\ell_P^2},
\qquad
q=\operatorname{ind}\bar\partial_{\mathcal O(q-1)},
\qquad
q=(L+1)^2,
\]
is a microscopic derivation of the capacity sector.  Conversely, a
model that computes none of these quantities has not solved the
amplitude problem; it has only renamed it.
\end{theorem}

\begin{proof}
The previous sections show that \(\Lambda\) is fixed once \(q\) is
fixed, and that topology/index/fuzzy geometry allow every integer
\(q\).  The no-go theorem above excludes purely polynomial, scale-free
selection of \(q_{\rm obs}\) without a large input or fine tuning.  The
remaining invariant possibilities are precisely exponential
transmutation, phase-corrected shell selection, or a direct code-space
dimension.  Each gives \(q\), and \(q\) gives \(\Lambda_q\).
\end{proof}



\section{\texorpdfstring{Boundary Chern--Simons level lock and the \(N\)-bound}{Boundary Chern--Simons level lock and the N-bound}}

The preceding selector ledger left one large integer to be selected,
\[
q_{\rm obs}=N_\partial\simeq 3.31\times10^{122}.
\]
There is, however, a sharper structural identification than the
spectral and spin-\(c\) readings alone.  In the positive-\(\Lambda\)
loop-gravity boundary treatment, the relevant boundary term is a
Chern--Simons theory and its level is related to the cosmological
constant by \cite{Smolin2002}
\begin{equation}
k_{\rm CS}=\frac{6\pi}{\lambda},
\qquad
\lambda=\ell_P^2\Lambda .
\end{equation}
Combining this with the capacity definition
\[
N_\partial=\frac{3\pi}{\ell_P^2\Lambda}
\]
gives the exact lock
\begin{equation}
\boxed{
N_\partial=\frac{k_{\rm CS}}{2}.
}
\label{eq:cs-level-lock}
\end{equation}
Thus the capacity integer is not only a de Sitter entropy count and a
spin-\(c\) index.  It is also one half of the boundary Chern--Simons
level in the positive-cosmological-constant quantum-group description.

In the Euclidean root-of-unity form the quantum-group deformation is
\begin{equation}
\mathfrak q_{\rm QG}
=
\exp\!\left(\frac{2\pi i}{k_{\rm CS}+2}\right),
\end{equation}
so the observed capacity corresponds to
\begin{equation}
k_{\rm CS}=6.61\times10^{122},
\qquad
\theta_{\rm QG}:=\frac{2\pi}{k_{\rm CS}+2}
=9.50\times10^{-123}.
\end{equation}
This gives a second exact formulation of the remaining problem:
selecting \(\Lambda\) is equivalent to selecting the boundary
Chern--Simons level \(k_{\rm CS}\).

\begin{theorem}[Chern--Simons capacity lock]
Assume the positive-\(\Lambda\) boundary Chern--Simons relation
\(k_{\rm CS}=6\pi/(\ell_P^2\Lambda)\) and the horizon-capacity relation
\(N_\partial=3\pi/(\ell_P^2\Lambda)\).  Then \(N_\partial=k_{\rm CS}/2\).
Consequently every microscopic selector for the capacity sector is
equivalently a selector for the Chern--Simons level, for the maximum
quantum-group spin \(j_{\max}=k_{\rm CS}/2\), and for the de Sitter
capacity \(N_\partial\).
\end{theorem}

\begin{proof}
The statement follows by division of the two displayed formulae:
\[
\frac{k_{\rm CS}}{N_\partial}
=
\frac{6\pi/(\ell_P^2\Lambda)}
     {3\pi/(\ell_P^2\Lambda)}
=2.
\]
The identification with the largest allowed quantum-group spin is the
standard root-of-unity truncation: for \(SU(2)_{k_{\rm CS}}\) the
admissible representations are cut off at \(j_{\max}=k_{\rm CS}/2\), up
to the conventional shift in \(k+2\) used in the deformation parameter.
\end{proof}

\begin{remark}[Status of the lock]
Equation~\eqref{eq:cs-level-lock} does not by itself select the
observed value of \(k_{\rm CS}\).  It eliminates an ambiguity: the
previously introduced capacity integer is the same large integer that
controls the finite quantum-group boundary theory.  The remaining
frontier is therefore not ``what is \(N_\partial\)?'' but
\[
\boxed{
\text{what boundary dynamics selects }k_{\rm CS}=2N_\partial?
}
\]
\end{remark}

The \(\varphi\)-registry residual now has a natural valuation-scale
comparison.  Let
\[
d_{\rm val}:=\dim\mathfrak A_{\mathscr D}=5.
\]
For the \(N=292\) shell one finds
\begin{equation}
\epsilon_\varphi
:=
\frac{\pi\varphi^{584}-N_\partial}{N_\partial}
=0.0628338\ldots,
\end{equation}
while
\begin{equation}
\frac{\pi}{2d_{\rm val}^2}
=
\frac{\pi}{50}
=0.0628319\ldots .
\end{equation}
Thus
\begin{equation}
\boxed{
N_\partial
=
\frac{\pi\varphi^{584}}
 {1+\pi/(2d_{\rm val}^2)}
\left[1+O(10^{-6})\right].
}
\label{eq:valuation-phase-clue}
\end{equation}
This is not promoted to a theorem.  It is a sharpened selector target:
a successful boundary phase action would have to derive the small
dimensionless phase defect \(\pi/(2d_{\rm val}^2)\) from the
five-idempotent valuation algebra, or replace it by a better exact
invariant.

\begin{figure}[!htbp]
\includegraphics[width=\columnwidth]{fig_cs_level_lock.png}
\caption{Chern--Simons level lock.  The same large integer appears as
the de Sitter capacity \(N_\partial\), half the boundary Chern--Simons
level \(k_{\rm CS}\), the maximum quantum-group spin \(j_{\max}\), and
the inverse cosmological scale.  The lock does not yet select the level;
it makes the remaining selector problem discrete and boundary-theoretic.}
\label{fig:cs-level-lock}
\end{figure}

\begin{criterion}[Level-selector completion]
A microscopic completion of the capacity sector must provide a boundary
Hamiltonian, beta function, phase action, or code-space rule that selects
the integer
\[
k_{\rm CS,*}=2N_{\partial,*}
\]
without inserting \(\Lambda_{\rm obs}\).  After
\cref{eq:cs-level-lock}, any such level selector is automatically a
capacity selector and hence fixes \(\Lambda_{\rm cap}\).
\end{criterion}



\section{Hopfion--helicity bridge and Schwinger admissibility}
\label{sec:hopf-schwinger}

The Hopf sector gives a real physical bridge between the spatial
\(\Sthree\) language and observed three-dimensional topological textures.
It must, however, be read with the correct logical strength.  Laboratory
hopfions show that the map \(S^3\to S^2\) and its linking invariant are
physically realizable in fields and materials; they do not by themselves
prove that the global spatial section of the universe is \(\Sthree\).
They do require that the Hopf channel be kept in the invariant atlas
whenever the \(\Sthree\) hypothesis is being tested.

\begin{definition}[Hopf field sector]
Let \(\mathbf n:\mathbb R^3\to S^2\) be a finite-energy orientation
field with a fixed value at spatial infinity.  Compactification gives
\[
\mathbb R^3\cup\{\infty\}\simeq S^3,
\qquad
\mathbf n:S^3\to S^2 .
\]
The Hopf charge is the integer
\begin{align}
Q_H[\mathbf n]
&=
\frac{1}{(4\pi)^2}
\int_{S^3} A\wedge F,
\\
F&=\mathbf n^*\omega_{S^2},
\qquad dA=F .
\end{align}
with the normalization chosen so that the standard Hopf fibration has
\(Q_H=1\).
\end{definition}

\begin{theorem}[Hopf charge as Chern--Simons helicity]
For a smooth finite-energy Hopf field, \(Q_H\in\mathbb Z\) equals the
linking number of two regular fibers of \(\mathbf n:S^3\to S^2\).  In
abelian gauge language the same quantity is the normalized
Chern--Simons/helicity integral
\[
\int A\wedge dA .
\]
Thus the observed Hopfion sector belongs to the same invariant family as
magnetic helicity and the boundary Chern--Simons level \(k_{\rm CS}\),
but it does not fix \(k_{\rm CS}\) or \(N_\partial\) by itself.
\end{theorem}

\begin{proof}
For a regular value \(p\in S^2\), \(\mathbf n^{-1}(p)\) is a closed
curve in \(S^3\).  The Hopf invariant is the linking number of the
preimages of two regular values.  Since \(F=\mathbf n^*\omega_{S^2}\)
is closed and exact on \(S^3\), \(F=dA\), and the Chern--Simons integral
\(A\wedge F\) computes the same integer after the stated normalization \cite{Ranada1989,Moffatt1969,Urbantke2003}.
\end{proof}

\begin{remark}[Experimental status]
Magnetic hopfions have been created and observed in multilayers and
chiral magnets, and hopfion/toron-type textures are also studied in
chiral liquid crystals.  These experiments confirm the physical
realizability of the Hopf invariant in matter fields.  They support the
inclusion of a Hopf channel in the invariant sieve; they do not by
themselves determine the cosmological spatial topology
\cite{Kent2021,Zheng2023,Ackerman2015}.
\end{remark}

The energy side of the same topological sector is supplied by helicity
bounds.  In fluid and magnetic systems,
\[
\mathcal H_B=\int_{\Sigma} \mathbf A\cdot\mathbf B\,d^3x
             =\int_{\Sigma} A\wedge dA
\]
is the magnetic helicity.  For thin flux tubes it reduces to linking
data, for example \(2\,{\rm Lk}\,\Phi_1\Phi_2\) for two linked tubes
up to internal twist and writhe terms.  In Faddeev--Skyrme type models
the Hopf charge produces a topological energy lower bound of
Vakulenko--Kapitanskii type,
\begin{equation}
\boxed{
E_{\rm Hopf}\ge C\,|Q_H|^{3/4},
}
\end{equation}
with the constant \(C\) depending on the model couplings \cite{Vakulenko1979,Harland2014}.  This supplies
an admissible energy penalty for a boundary selector but not a unique
cosmological capacity level.

The Schwinger scale supplies a second independent admissibility wall.  It
is not a capacity selector, but it excludes boundary electromagnetic
exchange configurations whose local field strength would trigger
unsuppressed pair creation \cite{Schwinger1951,Dunne2009}:
\begin{align}
E_{\rm Sch}
&=
\frac{m_e^2c^3}{e\hbar}
=
1.3233\times10^{18}\ {\rm V\,m^{-1}},
\\
B_{\rm Sch}&=E_{\rm Sch}/c .
\end{align}
For the present capacity anchor \(\Lambda_0=1.0909\times10^{-52}\,{\rm
m^{-2}}\), the dark-energy density may be written as an equivalent
electric-field energy density,
\[
\rho_\Lambda=\frac{\Lambda_0c^4}{8\pi G},
\qquad
E_\Lambda=\sqrt{\frac{2\rho_\Lambda}{\epsilon_0}} .
\]
Numerically,
\begin{align}
E_\Lambda&=10.89\ {\rm V\,m^{-1}},
\\
\frac{E_\Lambda}{E_{\rm Sch}}
&=8.23\times10^{-18},
\\
\pi\frac{E_{\rm Sch}}{E_\Lambda}
&=3.82\times10^{17}.
\end{align}
Thus the mean dark-energy scale is far below the QED Schwinger wall.
A boundary electric field at the Schwinger limit on a de Sitter sphere
would carry
\begin{equation}
\boxed{
N_e^{\rm Sch}
=
\frac{4\pi\epsilon_0R_\Lambda^2E_{\rm Sch}}{e}
=
2.53\times10^{79}
}
\end{equation}
elementary charges, many orders below
\[
N_\partial=3.31\times10^{122}.
\]
Therefore Schwinger physics is an admissibility constraint on exchange
fields, not a derivation of the horizon capacity.




\begin{criterion}[Hopf--Schwinger selector admissibility]
A microscopic boundary selector may include Hopf and electromagnetic
terms only through invariant quantities:
\begin{equation}
\Gamma_\partial^{\rm HS}
=
\Gamma_\partial
+
\lambda_H\,\mathcal E_H(Q_H)
+
\lambda_{\rm hel}\,\mathcal H_{\rm CS}
+
\lambda_{\rm Sch}\,\mathcal B_{\rm Sch}[E],
\end{equation}
where \(\mathcal E_H\) is a Hopf-charge energy penalty,
\(\mathcal H_{\rm CS}=\int A\wedge dA\) is a helicity/Chern--Simons
functional, and \(\mathcal B_{\rm Sch}\) penalizes fields approaching
\(E_{\rm Sch}\).  These terms may constrain the selector, but none of
them alone selects \(q_{\rm obs}\).
\end{criterion}





\section{\texorpdfstring{Ginzburg--Landau vortex flux and admissible \(U(1)\) winding}{Ginzburg--Landau vortex flux and admissible U(1) winding}}
\label{sec:gl-vortex-admissibility}

The projected slide supplied for this revision is useful because it shows the
correct topological level of superconducting vortices.  The complex
Ginzburg--Landau order parameter is a scalar phase field,
\(\psi=|\psi|e^{i\theta}\); it is not a spinor.  The topological integer is
therefore a \(U(1)\) winding number around a vortex core, not a spin
representation label.  This gives a laboratory-realized flux quantization
channel which must be kept separate from both the Hopf charge
\(\pi_3(S^2)\) and the horizon capacity index \(q\).

For a Cooper-pair charge \(q_{\rm cp}=2e\), the static Ginzburg--Landau
functional may be written schematically as
\begin{align}
\mathcal F_{\rm GL}
=&
\int d^3x\,
\bigg[
\alpha |\psi|^2+\frac{\beta}{2}|\psi|^4
+\frac{|\mathbf B|^2}{2\mu_0}
\nonumber\\
&\hspace{7mm}
+\frac{1}{2m^*}
\left|
\left(-i\hbar\nabla-q_{\rm cp}\mathbf A\right)\psi
\right|^2
\bigg].
\end{align}
Its type-II phase admits Abrikosov vortices and vortex lattices
\cite{GinzburgLandau1950,Abrikosov2004,Tinkham2004}.

\begin{theorem}[Ginzburg--Landau flux integer]
\label{thm:gl-flux}
Let \(\Gamma\) be a closed loop in a superconducting region where
\(|\psi|\ne0\), enclosing vortex cores but lying where the gauge-invariant
supercurrent vanishes.  Then
\begin{equation}
n_\Gamma
:=
\frac{1}{2\pi}\oint_\Gamma d\theta
\in\mathbb Z,
\end{equation}
and the magnetic flux through any spanning surface \(\Sigma_\Gamma\) obeys
\begin{equation}
\boxed{
\Phi_\Gamma
=
\int_{\Sigma_\Gamma}\mathbf B\cdot d\mathbf S
=
n_\Gamma\,\Phi_0,
\qquad
\Phi_0=\frac{h}{2e}.
}
\end{equation}
Equivalently, \(n_\Gamma\) is the first Chern number of the \(U(1)\) bundle
seen by the condensate around the vortex.
\end{theorem}

\begin{proof}
Far from the vortex core the vanishing of the gauge-invariant current gives
\[
\nabla\theta-\frac{q_{\rm cp}}{\hbar}\mathbf A=0
\]
along \(\Gamma\).  Single-valuedness of \(\psi\) gives
\(\oint_\Gamma d\theta=2\pi n_\Gamma\).  Therefore
\[
\frac{q_{\rm cp}}{\hbar}\oint_\Gamma\mathbf A\cdot d\mathbf l
=
2\pi n_\Gamma .
\]
By Stokes' theorem \(\oint\mathbf A\cdot d\mathbf l=\Phi_\Gamma\), hence
\[
\Phi_\Gamma=n_\Gamma\,\frac{2\pi\hbar}{q_{\rm cp}}
=n_\Gamma\,\frac{h}{2e}.
\]
\end{proof}

\begin{remark}[Three distinct topological integers]
The manuscript now contains three different integers which must not be
identified:
\[
\begin{gathered}
n_\Gamma\in\pi_1(U(1)),
\qquad
Q_H\in\pi_3(S^2),
\\
q=N_\partial
=\dim H^0(\CP^1,\mathcal O(q-1)).
\end{gathered}
\]
The first is a superconducting vortex winding, the second is a Hopf linking
charge, and the third is the horizon capacity sector.  Linked or knotted
vortex tubes can contribute to the helicity functional
\(\int A\wedge dA\), as in the Moffatt helicity invariant, but this does not
make the three integers equal \cite{Moffatt1969}.
\end{remark}

\begin{criterion}[Ginzburg--Landau contribution to the selector]
The vortex sector can enter the boundary selector only through admissible
gauge-invariant data:
\begin{equation}
\Gamma_\partial^{\rm GL}
=
\Gamma_\partial
+
\lambda_{\rm flux}
\sum_{\Gamma}
\left(
\frac{\Phi_\Gamma}{\Phi_0}-n_\Gamma
\right)^2
+
\lambda_{\rm GL}\,\mathcal F_{\rm GL}
+
\lambda_{\rm link}\,\mathcal H_{\rm hel}.
\end{equation}
Here \(\mathcal H_{\rm hel}\) is a helicity/Chern--Simons functional.
The first term enforces flux quantization, while the second and third terms
penalize energetic and linking configurations.  This supplies a real
laboratory topological channel, but not a standalone selector for
\(q_{\rm obs}\).
\end{criterion}

\begin{proposition}[Schwinger magnetic wall versus vortex flux count]
\label{prop:schwinger-gl-wall}
Let
\[
B_{\rm Sch}:=\frac{E_{\rm Sch}}{c}
=
\frac{m_e^2c^2}{e\hbar}.
\]
The maximum number of superconducting flux quanta that a field of order
\(B_{\rm Sch}\) could thread through the de Sitter horizon area
\(A_\Lambda=4\pi R_\Lambda^2\) is
\begin{equation}
N_\Phi^{\rm Sch}
=
\frac{A_\Lambda B_{\rm Sch}}{\Phi_0}.
\end{equation}
Using the same Planck-2018/CODATA anchor as above gives
\begin{equation}
\boxed{
\begin{gathered}
\Phi_0=2.067833848\times10^{-15}\ {\rm Wb},
\\
B_{\rm Sch}=4.4140\times10^{9}\ {\rm T}.
\end{gathered}}
\end{equation}
\begin{equation}
\boxed{
N_\Phi^{\rm Sch}=7.38\times10^{77},
\qquad
\frac{N_\partial}{N_\Phi^{\rm Sch}}
=4.48\times10^{44}.
}
\end{equation}
Hence Schwinger-limited electromagnetic flux quanta cannot directly count
the horizon capacity states.  The Ginzburg--Landau/Schwinger sector is an
admissibility and helicity sector, not the missing microscopic derivation of
\(q_{\rm obs}\).
\end{proposition}

\begin{remark}[Critical-field translation]
The Ginzburg--Landau upper critical field satisfies
\[
B_{c2}=\frac{\Phi_0}{2\pi\xi^2}.
\]
If one formally sets \(B_{c2}=B_{\rm Sch}\), the corresponding coherence
length is
\[
\xi_{\rm Sch}
=
\sqrt{\frac{\Phi_0}{2\pi B_{\rm Sch}}}
=
2.73\times10^{-13}\ {\rm m}
\simeq1.69\times10^{22}\ell_P.
\]
This is a useful scale comparison, not a quantum-gravity derivation.
\end{remark}



\section{Bradlow horizon packing and Planck-vortex capacity bound}
\label{sec:bradlow-packing}
\label{sec:bradlow-capacity}

The Ginzburg--Landau vortex layer has one more consequence which is more
useful than the raw Schwinger comparison.  On a compact two-dimensional
surface the Abelian Higgs equations do not admit an arbitrary number of
vortices at fixed area.  The integrated vortex equation gives a Bradlow-type
area bound: the vortex number is limited by the available area.  This is a
packing theorem, not a new force law.

In a common critical-coupling normalization the vortex equations on a compact
Riemann surface \(\Sigma\) contain
\begin{equation}
F_A+\frac12\bigl(|\psi|^2-\tau\bigr)\omega_\Sigma=0,
\qquad
\frac{1}{2\pi}\int_\Sigma F_A=N_{\rm v}\in\mathbb Z ,
\end{equation}
with \(\tau>0\).  Integrating gives
\begin{equation}
2\pi N_{\rm v}
=\frac12\int_\Sigma(\tau-|\psi|^2)\,\omega_\Sigma
\le
\frac{\tau}{2}\operatorname{Area}(\Sigma),
\end{equation}
and therefore
\begin{equation}
\boxed{
N_{\rm v}\le
\frac{\tau\,\operatorname{Area}(\Sigma)}{4\pi}.
}
\label{eq:bradlow-bound}
\end{equation}
This is the Bradlow bound for Abelian Higgs vortices on compact surfaces,
up to the conventional normalization of \(\tau\) \cite{Bradlow1990,AdamSpeightWereszczynski2017,ZhangAbrikosov}.

For the de Sitter horizon \(\Sigma=S^2_{R_\Lambda}\), put
\begin{equation}
A_\Lambda=4\pi R_\Lambda^2.
\end{equation}
There is a distinguished Planck--Bradlow normalization
\begin{equation}
\boxed{
\tau_P:=\frac{\pi}{\ell_P^2}.
}
\end{equation}
Then \cref{eq:bradlow-bound} becomes
\begin{equation}
\boxed{
N_{\rm v}^{\max}
=
\frac{\tau_P A_\Lambda}{4\pi}
=
\frac{A_\Lambda}{4\ell_P^2}
=
\Ncap .
}
\end{equation}
Thus the same integer which appears as the horizon entropy capacity is the
maximum number of Planck--Bradlow vortex cells that can be packed on the
horizon.  Equivalently, one elementary Bradlow cell has area
\begin{equation}
\boxed{
A_{\rm cell}^{P}=\frac{4\pi}{\tau_P}=4\ell_P^2,
\qquad
\xi_P:=\tau_P^{-1/2}=\frac{\ell_P}{\sqrt{\pi}}.
}
\end{equation}

This result does not say that the cosmological horizon is an ordinary
superconductor.  It says something weaker and sharper: whenever a boundary
selector is represented by a compact \(U(1)\) Higgs/vortex sector, its
topological vortex number obeys the same area-capacity law, and saturation at
the Planck--Bradlow normalization reproduces precisely the Bekenstein--Hawking
capacity.  Hence the allowed selector action may contain a Bradlow packing
term,
\begin{equation}
\boxed{
\Gamma_{\partial}^{\rm Br}
=
\lambda_{\rm Br}
\left[
\frac{N_{\rm v}}{\Ncap}
-
\frac{\tau}{\tau_P}
\right]^2
+
\lambda_{\rm sat}
\left(1-\frac{N_{\rm v}}{\Ncap}\right)^2 ,
}
\end{equation}
where the second term is imposed only in a saturation hypothesis.  Without
such a saturation principle, the Bradlow theorem supplies an upper bound and a
packing channel, not a derivation of the observed \(q=\Ncap\).

Numerically, for the present capacity anchor,
\begin{align}
\Ncap&\simeq3.31\times10^{122},\\
\tau_P&\simeq1.20\times10^{70}\ {\rm m}^{-2},\\
\xi_P&\simeq9.12\times10^{-36}\ {\rm m}.
\end{align}
Thus the vortex-packing interpretation is exactly Planckian, whereas the
electromagnetic Schwinger flux count of \cref{sec:gl-vortex-admissibility}
is smaller by more than forty orders of magnitude.  The two counts are
therefore different invariants: Schwinger is an admissibility wall;
Bradlow saturation is an area-capacity packing law.



\section{Bogomolny--Bradlow saturation as a boundary selector}
\label{sec:bogomolny-bradlow-selector}

The previous section gave the Bradlow inequality.  The remaining question is
whether saturation can be imposed by a variational principle rather than by an
extra slogan.  In the Abelian Higgs model at critical coupling this is exactly
what the Bogomolny decomposition supplies: it separates an integer topological
term from a sum of nonnegative squares.

Let \(\Sigma\) be a compact oriented Riemann surface with area form
\(\omega_\Sigma\), let \(L\to\Sigma\) be a Hermitian line bundle, and let
\((A,\psi)\) be a unitary connection and section.  In the same normalization as
\cref{sec:bradlow-capacity}, consider the nonnegative BPS functional
\begin{align}
\mathcal E_{\rm sq}[A,\psi;\tau]
:={}&
\int_\Sigma 2\,|\bar\partial_A\psi|^2 \nonumber\\
&+
\int_\Sigma
\left|
F_A+\frac12(|\psi|^2-\tau)\omega_\Sigma
\right|^2 .
\end{align}
Its zero set is the vortex system
\begin{equation}
\bar\partial_A\psi=0,
\qquad
F_A+\frac12(|\psi|^2-\tau)\omega_\Sigma=0.
\end{equation}
The second equation gives the exact identity
\begin{equation}
\boxed{
\int_\Sigma |\psi|^2\omega_\Sigma
=
\tau\,{\rm Area}(\Sigma)-4\pi N_{\rm v},
\qquad
N_{\rm v}:=\frac1{2\pi}\int_\Sigma F_A .
}
\label{eq:bradlow-identity}
\end{equation}
Thus the Bradlow inequality is not merely an estimate.  It is the positivity
condition for the integrated Higgs norm.

\begin{theorem}[Planck--Bradlow saturation selector]
\label{thm:planck-bradlow-selector}
Assume that the boundary sector is described by a compact critical Abelian
Higgs system on the de Sitter horizon \(S^2_{R_\Lambda}\), with Planck--Bradlow
coupling
\[
\tau=\tau_P:=\pi/\ell_P^2,
\]
and that the null horizon capacity screen carries no local Higgs condensate,
\[
\int_{S^2_{R_\Lambda}}|\psi|^2\omega=0.
\]
Then the vortex number is forced to be
\begin{equation}
\boxed{
N_{\rm v}
=
\frac{\tau_P A_\Lambda}{4\pi}
=
\frac{A_\Lambda}{4\ell_P^2}
=
\Ncap .
}
\end{equation}
Conversely, within the BPS equations and \(\tau=\tau_P\), the equality
\(N_{\rm v}=\Ncap\) is equivalent to the vanishing of the integrated Higgs
norm.
\end{theorem}

\begin{proof}
Substitute \(\tau=\tau_P\) and \(\Sigma=S^2_{R_\Lambda}\) in
\cref{eq:bradlow-identity}.  Since the integrand \(|\psi|^2\) is
nonnegative, the condition \(\int|\psi|^2\omega=0\) is equivalent to
\[
4\pi N_{\rm v}=\tau_P A_\Lambda.
\]
Using \(\tau_P=\pi/\ell_P^2\) gives \(N_{\rm v}=A_\Lambda/(4\ell_P^2)\), which
is \(\Ncap\).  The converse is the same identity read backwards.
\end{proof}

This theorem removes one previously open ambiguity.  Bradlow saturation need
not be a manually imposed equality; it follows from three invariant boundary
conditions: critical coupling, Planck normalization, and the absence of a
local Higgs condensate on the capacity screen.  It still does not determine
\(A_\Lambda\) or the observed numerical sector by itself.  It determines the
capacity once the horizon area is fixed.

Equivalently, define the boundary selector
\begin{equation}
\boxed{
\begin{aligned}
\Gamma_{\partial}^{\rm BPS-Br}
={}&
\lambda_{\rm BPS}\,\mathcal E_{\rm sq}[A,\psi;\tau]
+\lambda_P(\tau-\tau_P)^2\\
&+
\lambda_0
\left(
\int_\Sigma |\psi|^2\omega_\Sigma
\right)^2 .
\end{aligned}}
\end{equation}
In the limit of positive infinite weights, the minimizers satisfy the
hypotheses of \cref{thm:planck-bradlow-selector}.  Hence
\[
\boxed{
\Gamma_{\partial}^{\rm BPS-Br}\ {\rm minimized}
\quad\Longrightarrow\quad
N_{\rm v}=\Ncap .
}
\]
The missing microscopic step is therefore sharpened again: a final boundary
theory must derive the three selector conditions above from the edge-mode
Hamiltonian, not assume them.

\begin{corollary}[Capacity dictionary after the BPS--Bradlow selector]
On the Planck--Bradlow saturated boundary,
\begin{equation}
\boxed{
\Ncap
=
N_{\rm v}
=
\frac{A_\Lambda}{4\ell_P^2}
=
q
=
\frac{k_{\rm CS}}2 .
}
\end{equation}
Thus the same integer is simultaneously a vortex-packing number, a horizon
entropy capacity, a holomorphic-index dimension, and half of the boundary
Chern--Simons level.  These equalities are conditional on the common boundary
sector; they are not separate numerical coincidences.
\end{corollary}




\section{Boundary edge-mode Hamiltonian and entropy-cell locking}
\label{sec:edge-mode-hamiltonian}

The previous Bradlow argument still left three conditions external to the
functional: critical coupling, Planck normalization, and the null-condensate
condition at the capacity screen.  They can be written as one local boundary
edge-mode selector.  This is not a new bulk matter field.  It is a boundary
admissibility functional on the horizon two-sphere.

Let \(\Sigma=S^2_R\), let \(A\) be a \(U(1)\) connection on a complex line
bundle over \(\Sigma\), and let \(\psi\) be a boundary section.  In units where
the round metric on \(\Sigma\) is fixed, define
\begin{equation}
\begin{aligned}
B_A&:=*F_A+\frac12(|\psi|^2-\tau),\\
\Gamma_{\partial}^{\rm edge}
={}&
\int_\Sigma
\left(|D_{\bar z}\psi|^2+\frac12 B_A^2\right)\omega_\Sigma\\
&+\lambda_P(\tau\ell_P^2-\pi)^2\\
&+\lambda_0
\left(\int_\Sigma|\psi|^2\omega_\Sigma\right)^2 .
\end{aligned}
\label{eq:edge-selector}
\end{equation}
The first line is the Bogomolny form of the compact Abelian-Higgs boundary
energy.  The second line imposes, in the sharp-selector limit
\(\lambda_P,\lambda_0\to\infty\), the Planck entropy-cell normalization and
the null-condensate condition.  Thus the critical coupling is no longer a
loose phenomenological choice: it is the unique self-dual point at which the
compact vortex number is controlled by a first-order topological bound.

\begin{theorem}[Edge-mode derivation of Planck--Bradlow capacity]
\label{thm:edge-mode-capacity}
Let \(A,\psi,\tau\) minimize \(\Gamma_{\partial}^{\rm edge}\) in the sharp
selector limit on a fixed topological sector
\[
\frac{1}{2\pi}\int_\Sigma F_A=N_{\rm v}\in\mathbb Z .
\]
If the square term in \cref{eq:edge-selector} is saturated, then
\begin{equation}
\begin{aligned}
D_{\bar z}\psi&=0,\\
*F_A+\frac12(|\psi|^2-\tau)&=0,\\
\tau&=\frac{\pi}{\ell_P^2},\\
\int_\Sigma|\psi|^2\omega_\Sigma&=0 .
\end{aligned}
\end{equation}
Consequently
\begin{equation}
\boxed{
N_{\rm v}
=
\frac{\tau A_\Sigma}{4\pi}
=
\frac{A_\Sigma}{4\ell_P^2}.
}
\end{equation}
For \(\Sigma=S^2_{R_\Lambda}\), this gives
\begin{equation}
\boxed{
N_{\rm v}=N_\partial=q=\frac{k_{\rm CS}}2 .
}
\end{equation}
\end{theorem}

\begin{proof}
The two nonnegative terms in the first line of
\cref{eq:edge-selector} are minimized by the stated first-order equations.
The sharp Planck-lock term gives \(\tau\ell_P^2=\pi\), while the sharp
null-condensate term gives \(\int|\psi|^2\omega_\Sigma=0\).
Integrating the second Bogomolny equation over \(\Sigma\) gives
\[
\int_\Sigma F_A
=
\frac{\tau A_\Sigma}{2}
-
\frac12\int_\Sigma|\psi|^2\omega_\Sigma .
\]
Using \(\int F_A=2\pi N_{\rm v}\) and the null-condensate condition yields
\(N_{\rm v}=\tau A_\Sigma/(4\pi)\).  Substituting
\(\tau=\pi/\ell_P^2\) proves the result.  The equality with
\(N_\partial\) is precisely the Gibbons--Hawking entropy normalization, and
the equality with \(k_{\rm CS}/2\) is the Chern--Simons level lock already
derived above.
\end{proof}

\begin{criterion}[What remains to be derived]
\label{crit:edge-remaining}
\Cref{thm:edge-mode-capacity} closes the local boundary-selector logic:
critical coupling, Planck normalization, and Bradlow saturation can be
implemented by a single gauge-invariant edge functional.  It does not yet
derive why the universe occupies the particular large topological sector
\[
q\simeq3.31\times10^{122}.
\]
The remaining microscopic problem is therefore sharpened to the selection of
the integer sector of the edge Hamiltonian,
\begin{equation}
q_*=\arg\min_{q\in\mathbb N}\Gamma_{\partial}^{\rm edge,eff}(q),
\end{equation}
without inserting \(\Lambda_{\rm obs}\) or \(q_{\rm obs}\) as input.
\end{criterion}



\section{Boundary-selector Hamiltonian: common edge action, no-go theorem, and the only surviving transmutation route}
\label{sec:edge-selector}

The preceding sections identified several exact readings of the same
capacity integer,
\begin{equation}
q=\Ncap=\frac{k_{\rm CS}}2=\operatorname{ind}\bar\partial
      =\operatorname{Tr}\Pi_L=N_{\rm v}^{\rm Bradlow}.
\end{equation}
This does not yet choose the observed sector.  The selection problem is
therefore not the definition of \(q\), but the derivation of an effective
edge action whose minimum lies at the realized sector.

\begin{definition}[Edge-sector Hilbert decomposition]
Let the boundary Hilbert space decompose into capacity sectors
\begin{equation}
\mathcal H_\partial=\bigoplus_{q\in\mathbb N}\mathcal C_q .
\end{equation}
Here \(q\) is the entropy capacity, or number of Planck--Bradlow cells.
The holomorphic space \(H^0(\CP^1,\mathcal O(q-1))\) supplies the
one-particle \(SU(2)\)-equivariant cell basis of dimension \(q\), while
the full admissible code sector \(\mathcal C_q\) has entropy
\begin{equation}
\log\dim \mathcal C_q=q .
\end{equation}
This separates the one-particle index count from the many-body entropy
capacity and removes the otherwise misleading identification
\(\Ncap=\dim\mathcal H_\partial\).
\end{definition}

Let \(H_{\partial,q}\) be the edge-mode Hamiltonian restricted to
\(\mathcal C_q\).  The sector weight is
\begin{equation}
Z_q=\operatorname{Tr}_{\mathcal C_q}e^{-\beta_q H_{\partial,q}}
    =\exp[-\Gamma_\partial(q)] .
\end{equation}
Thus the selector is
\begin{equation}
q_*=\arg\min_{q\in\mathbb N}\Gamma_\partial(q).
\end{equation}

\begin{theorem}[Reduction of the sharp edge Hamiltonian]
\label{thm:edge-selector-reduction}
Assume the boundary edge theory is in the sharp Planck--Bradlow phase:
\begin{align}
D_{\bar z}\psi&=0,
&
*F_A+\frac12(|\psi|^2-\tau)&=0,
\\
\tau&=\frac{\pi}{\ell_P^2},
&
\int_\Sigma|\psi|^2\omega_\Sigma&=0 .
\end{align}
Then the Bogomolny--Bradlow part of the Hamiltonian vanishes on every
admissible \(q\)-sector and fixes only
\begin{equation}
q=\frac{A_\Sigma}{4\ell_P^2}.
\end{equation}
It does not select which admissible value of \(q\) is realized.
The remaining selector has the universal sector form
\begin{equation}
\boxed{
\Gamma_\partial(q)
=
\Gamma_{\rm cl}(q)
+\Gamma_{\zeta\eta}(q)
+\Gamma_{\rm CS}(q)
+\Gamma_{\rm exch}(q)
-q .
}
\end{equation}
The final term is the entropy of the code sector.
\end{theorem}

\begin{proof}
The Bogomolny decomposition is a sum of nonnegative square terms plus a
topological term.  In the sharp phase the square terms vanish and the
integrated Bradlow identity gives \(q=A_\Sigma/(4\ell_P^2)\).  Since the
same cancellation occurs in every sector satisfying the sharp conditions,
the BPS/Bradlow part alone is sector-degenerate.  After tracing over
\(\mathcal C_q\), the entropy contribution is \(-\log\dim\mathcal C_q=-q\),
while the non-BPS determinant, eta, Chern--Simons, and exchange pieces
remain in the effective action.  This gives the displayed form.
\end{proof}

\begin{nogotheorem}[No parameter-free topological selector]
\label{nogo:parameter-free-selector}
No selector built only from the topology of \(S^2\simeq\CP^1\), the
integer flux \(q-1\), \(SU(2)\)-equivariance, the index
\(\operatorname{ind}\bar\partial_{\mathcal O(q-1)}=q\), and sharp
Bogomolny--Bradlow saturation can isolate a finite large sector
\(q_*\gg1\).  It either permits all \(q\), selects a boundary endpoint,
or is exactly sector-degenerate.
\end{nogotheorem}

\begin{proof}
Those data specify the admissible sector labels and their representation
content, but introduce no dimensionless scale at which a large integer
can be distinguished from a neighboring large integer.  The index is
linear in the sector label, \(\operatorname{ind}=q\), and equivariance
makes the intra-sector Hamiltonian proportional to the identity on the
irreducible cell space.  The sharp Bradlow equations fix the area carried
by a chosen sector, not the chosen sector itself.  Therefore a finite
interior minimum requires an additional dimensionless selector datum:
a coupling, a phase, a determinant scale, a chemical potential, or an
exchange boundary condition.  Without such a datum the selector cannot
distinguish \(q_*\) from \(q_*+1\), nor can it generate the hierarchy
\(q_*\sim10^{122}\).
\end{proof}

The surviving possibility is dimensional transmutation in the edge
sector, in the Coleman--Weinberg sense of a scale generated by quantum
corrections \cite{ColemanWeinberg1973}.  Let a marginal boundary coupling \(g_\partial\) obey
\begin{equation}
\frac{\dd g_\partial}{\dd\ln\mu}
=
-\frac{b}{16\pi^2}g_\partial^3+O(g_\partial^5),
\qquad b>0 .
\end{equation}
Then the renormalization-group invariant scale is
\begin{equation}
\Lambda_{\rm edge}
=
\mu\exp\!\left[-\frac{8\pi^2}{b g_\partial(\mu)^2}\right],
\end{equation}
and an area-like capacity can be generated as
\begin{equation}
\boxed{
q_\star
=
\Xi\,
\exp\!\left[\frac{16\pi^2}{b g_\partial(\mu)^2}\right],
}
\label{eq:qstar-transmutation}
\end{equation}
where \(\Xi\) is a dimensionless determinant/phase prefactor.

\begin{theorem}[Coleman--Weinberg edge selector]
\label{thm:cw-edge-selector}
Suppose the determinant-line and exchange terms reduce, after integrating
out nonzero boundary modes, to the sector action
\begin{multline}
\Gamma_\partial^{\rm CW}(q)
=
\chi\,q\left[\ln\!\left(\frac{q}{q_\star}\right)-1\right]
+\nu\ln q
\\
+\Theta_{\rm bd}(q)+O(q^{-1}).
\end{multline}
with \(\chi>0\), slowly varying \(\Theta_{\rm bd}\), and \(q_\star\) given
by \cref{eq:qstar-transmutation}.  Then the minimum satisfies
\begin{equation}
\boxed{
q_{\rm min}
=
q_\star
\left[
1-\frac{\nu+\Theta'_{\rm bd}(q_\star)q_\star}{\chi q_\star}
+O(q_\star^{-2})
\right].
}
\end{equation}
Hence a naturally huge \(q_{\rm min}\) requires a transmuted edge scale;
topology and index quantization alone do not suffice.
\end{theorem}

\begin{proof}
Differentiate:
\[
\Gamma_\partial'(q)
=
\chi\ln(q/q_\star)+\nu/q+\Theta'_{\rm bd}(q)+O(q^{-2}).
\]
Put \(q=q_\star(1+\delta)\) and expand
\(\ln(1+\delta)=\delta+O(\delta^2)\).  Solving to first order gives
\[
\delta=
-\frac{\nu+\Theta'_{\rm bd}(q_\star)q_\star}{\chi q_\star}
+O(q_\star^{-2}).
\]
The second derivative is \(\Gamma_\partial''(q)=\chi/q+O(q^{-2})>0\) at
large \(q\), so the stationary point is a local minimum.
\end{proof}

For the observed capacity target used in the numerical ledgers,
\begin{equation}
q_{\rm obs}=3.3072\times10^{122},
\qquad
\ln q_{\rm obs}=282.111\pm0.015 .
\end{equation}
If \(\Xi=1\), the required transmutation combination is
\begin{equation}
\boxed{
b\,g_\partial^2
=
\frac{16\pi^2}{\ln q_{\rm obs}}
=
0.559756\ldots .
}
\end{equation}
This number is close to \(e^{-\gamma_{\rm E}}\), differing by
\(3.03\times10^{-3}\) fractionally, but the manuscript does not identify
the two.  Such a numerical proximity is admissible only as a selector
target until a boundary determinant calculation derives it.

\begin{criterion}[What would count as the solution of the selector problem]
The capacity problem is solved only when a concrete boundary operator
\(D_{\partial,q}\) and edge Hamiltonian \(H_{\partial,q}\) give
\begin{equation}
\Gamma_\partial(q)
=
-\ln\operatorname{Tr}_{\mathcal C_q}e^{-\beta_qH_{\partial,q}}
\end{equation}
and its determinant-line, eta-phase, Chern--Simons, Bradlow, Hopf, and
exchange contributions reduce to a selector with
\begin{equation}
q_{\rm min}=3.31\times10^{122}
\end{equation}
without using \(\Lambda_{\rm obs}\), \(H_0\), or \(q_{\rm obs}\) as an
input.  Until that calculation is supplied, the theory has converted the
cosmological constant into a sharply posed edge-mode selection problem,
not into a completed numerical derivation.
\end{criterion}




\section{Explicit minimal edge-mode Hamiltonian}
\label{sec:explicit-edge-mode-hamiltonian}

The preceding section isolated the missing object.  I now write the minimal
Hamiltonian which has the required logical form, with the determinant correction
counted once.  It is not yet a microscopic
derivation of the observed number.  It is the first closed selector ansatz in
which every retained connector has a precise operator role.

For each positive integer \(q\), let
\begin{equation}
R_q=\sqrt{\frac{q}{\pi}}\,\ell_P,
\qquad
\beta_q=2\pi R_q
\end{equation}
in units \(c=\hbar=k_B=1\).  The boundary Hilbert space is decomposed as
\begin{equation}
\boxed{
\mathcal H_\partial
=
\bigoplus_{q\ge1}\mathcal C_q,
\qquad
\widehat\Pi_q:\mathcal H_\partial\to\mathcal C_q,
\qquad
\log\dim\mathcal C_q=q .
}
\end{equation}
The last equation is the entropy-capacity convention.  The one-particle
index space \(H^0(\CP^1,\mathcal O(q-1))\) has dimension \(q\); the many-body
code sector \(\mathcal C_q\) has entropy \(q\).

The explicit edge Hamiltonian is
\begin{equation}
\boxed{
\widehat H_\partial
=
\bigoplus_{q\ge1}
\left[
\widehat H_{{\rm adm},q}
+
\widehat H_{{\rm run},q}
+
\widehat H_{{\rm osc},q}
\right].
}
\label{eq:explicit-edge-H}
\end{equation}
The three summands have different logical roles.

First, the admissibility Hamiltonian is a positive common-kernel operator:
\begin{align}
\beta_q\widehat H_{{\rm adm},q}
={}&
\lambda_{\rm Br}\,\widehat{\mathcal E}^{\rm Br}_q
+\lambda_P(\widehat\tau\ell_P^2-\pi)^2
+\lambda_0
\left(\int_{\Sigma_q}|\widehat\psi|^2\omega_q\right)^2
\nonumber\\
&+
\lambda_{\rm CS}(\widehat k_{\rm CS}-2q)^2
+\lambda_{\rm HD}\sum_A\widehat{\mathcal H}_A^\dagger
\widehat{\mathcal H}_A
\nonumber\\
&+
\lambda_{\rm tf}\widehat C_{\rm tf}^{\,\dagger}\widehat C_{\rm tf}
+\lambda_{\rm cap}
(\ell_P^2\widehat\Lambda\,q-3\pi)^2
\nonumber\\
&+
\lambda_Q(\widehat Q_{\rm DM}-q_{\rm DM})^2
\nonumber\\
&+
\lambda_{\rm val}
\left[
P_{\rm val}(\widehat\Dil)^\dagger
P_{\rm val}(\widehat\Dil)
\right],
\label{eq:Hadm-explicit}
\end{align}
with all \(\lambda_i>0\).  Here \(\Sigma_q=S^2_{R_q}\), and
\(\widehat{\mathcal E}^{\rm Br}_q\) is the compact Abelian-Higgs
Bogomolny square on \(\Sigma_q\),
\begin{equation}
\widehat{\mathcal E}^{\rm Br}_q
=
\int_{\Sigma_q}
\left(
2|D_{\bar z}\widehat\psi|^2+
\left|
\widehat F_A+\frac12(|\widehat\psi|^2-\widehat\tau)\omega_q
\right|^2
\right).
\end{equation}
Thus \(\widehat H_{{\rm adm},q}\) enforces the role-\({\rm Q}\),
role-\({\rm B}\), and role-\({\rm O}\) constraints: Bradlow saturation,
Planck cell normalization, Chern--Simons level lock, quantum covariance,
trace-free locality, capacity fixation, dark-charge superselection, and
valuation closure.  It does not select \(q\).  In the sharp limit
\(\lambda_i\to\infty\), inadmissible states are projected out.

Second, the selector Hamiltonian is the leading nonconformal running zero mode
\begin{equation}
\boxed{
\beta_q\widehat H_{{\rm run},q}
=
q\ln\!\left(\frac{q}{q_\star}\right)\mathbf 1_{\mathcal C_q}.
}
\label{eq:Hrun-explicit}
\end{equation}
The transmuted capacity scale is
\begin{equation}
\boxed{
q_\star
=
\Xi\,(\mu\ell_P)^{-2}
\exp\!\left[
\frac{16\pi^2}{b\,g_\partial(\mu)^2}
\right],
}
\label{eq:qstar-explicit-H}
\end{equation}
where \(g_\partial\) is a marginal boundary coupling with
\begin{equation}
\mu\frac{\dd g_\partial}{\dd\mu}
=
-\frac{b}{16\pi^2}g_\partial^3+O(g_\partial^5),
\qquad b>0 .
\end{equation}
The dimensionless prefactor \(\Xi\) is to be computed from the same edge
renormalization problem.  It is not an observational input.

Third, \(\widehat H_{{\rm osc},q}\) contains nonzero oscillators and
determinant-line fluctuations around the admissible edge background.  Its
renormalized contribution is counted here, and not again in
\(\widehat H_{{\rm run},q}\):
\begin{equation}
-\ln\operatorname{Tr}'_{\mathcal C_q}
\exp[-\beta_q\widehat H_{{\rm osc},q}]
=
\nu\ln q+\Theta_\partial(q)+O(q^{-1}),
\end{equation}
after zero modes and the entropy capacity have been separated.  This is the
operator location of the APS/zeta/eta data.  It is spectral information, not
a positive state count.

\begin{theorem}[Explicit edge Hamiltonian selector]
\label{thm:explicit-edge-selector}
Assume that \(\widehat H_{{\rm adm},q}\) has a nonempty zero space in each
admissible sector and a positive gap away from it.  Assume also that the
oscillator determinant has the asymptotic form stated above and that
\(\Theta_\partial\) varies slowly on the scale of one capacity cell.  Then
the sector partition function of \cref{eq:explicit-edge-H} has the form
\begin{equation}
Z_q=
\operatorname{Tr}_{\mathcal C_q}e^{-\beta_q\widehat H_{\partial,q}}
=
\exp[-\Gamma_\partial(q)]
\end{equation}
with
\begin{equation}
\boxed{
\Gamma_\partial(q)
=
q\left[\ln\!\left(\frac{q}{q_\star}\right)-1\right]
+\nu\ln q+\Theta_\partial(q)+O(q^{-1}).
}
\label{eq:Gamma-explicit-H}
\end{equation}
Hence the stable large-\(q\) selector is
\begin{equation}
\boxed{
q_{\rm min}
=
q_\star
\left[
1-\frac{\nu+q_\star\Theta_\partial'(q_\star)}{q_\star}
+O(q_\star^{-2})
\right],
}
\label{eq:qmin-explicit-H}
\end{equation}
with the physical integer chosen by the nearest admissible integer rule.
If \(\Theta_\partial'=O(q_\star^{-2})\), then
\[
q_{\rm min}=q_\star-\nu+O(q_\star^{-1}).
\]
\end{theorem}

\begin{proof}
On the admissible subspace, \(\widehat H_{{\rm adm},q}\) vanishes in the
sharp limit.  The entropy convention gives
\[
\operatorname{Tr}_{\mathcal C_q}{\bf 1}=e^q .
\]
The running zero mode contributes
\[
\exp[-q\ln(q/q_\star)],
\]
and the renormalized oscillator determinant contributes
\[
\exp[-\nu\ln q-\Theta_\partial(q)+O(q^{-1})]
\]
after zero modes are separated.  Therefore
\[
-\ln Z_q
=
q[\ln(q/q_\star)-1]+\nu\ln q+\Theta_\partial(q)+O(q^{-1}).
\]
Differentiating gives
\[
\Gamma_\partial'(q)
=
\ln(q/q_\star)+\nu/q+\Theta_\partial'(q)+O(q^{-2}).
\]
Put \(q=q_\star(1+\delta)\) and expand
\(\ln(1+\delta)=\delta+O(\delta^2)\).  Solving to first order yields
\cref{eq:qmin-explicit-H}.  The second derivative is
\[
\Gamma_\partial''(q)=q^{-1}-\nu q^{-2}
+\Theta_\partial''(q)+O(q^{-3}),
\]
which is positive at large \(q\) under the slow-variation assumption.
Thus the stationary point is a stable local minimum, and discreteness selects
the nearest admissible integer.
\end{proof}

\begin{corollary}[Observed-sector target for the explicit Hamiltonian]
If the microscopic edge theory uses \(\mu=\ell_P^{-1}\), has
\(\Xi=1+o(1)\), and has negligible \(\Theta_\partial'\) at the selected
sector, then matching the present capacity target requires
\begin{equation}
\boxed{
b\,g_\partial^2
=
\frac{16\pi^2}{\ln q_{\rm obs}}
=
0.559756\ldots .
}
\end{equation}
More generally,
\begin{equation}
\boxed{
b\,g_\partial^2
=
\frac{16\pi^2}
{\ln q_{\rm obs}-\ln\Xi+2\ln(\mu\ell_P)}
}
\end{equation}
up to the determinant displacement in \cref{eq:qmin-explicit-H}.  A final
theory must compute the right-hand side from the boundary beta function and
determinant data.  Choosing it after observing \(\Lambda\) is only a fit.
\end{corollary}

\begin{criterion}[Explicit-Hamiltonian completion test]
The capacity-selection problem is solved by \cref{eq:explicit-edge-H} only
when the following data are derived internally:
\[
\boxed{
b,\qquad g_\partial(\mu),\qquad \Xi,\qquad
\nu,\qquad \Theta_\partial(q).
}
\]
The admissibility Hamiltonian closes the invariant atlas; the running
Hamiltonian supplies the first genuine selector.  If these quantities are
not computed, this is a precise Hamiltonian ansatz and a falsification target,
not a numerical derivation of \(\Lambda\).
\end{criterion}

\begin{remark}[Why this is the narrowest honest Hamiltonian]
A purely topological Hamiltonian is sector-degenerate.  A purely zeta or eta
Hamiltonian is scheme-dependent and does not count positive states.  A purely
Bradlow Hamiltonian fixes area per vortex but not the realized vortex number.
The term \(q\ln(q/q_\star)\) is therefore not decorative: it is the minimal
nonconformal running term that can generate a stable large sector from
order-one microscopic data.  All other retained structures enter
\(\widehat H_{{\rm adm},q}\), \(\widehat H_{{\rm osc},q}\), or
\(\Theta_\partial(q)\).
\end{remark}


\section{APS eta and determinant-line selector}
\label{sec:aps-eta-selector}

The preceding Hopf--Chern--Simons layer supplies a topological and
helicity channel.  It still does not select the observed level.  The
next admissible refinement must be spectral: it must depend on the
boundary operator, not only on the topology, index, or Chern--Simons
integer.

Let \(D_{\partial,q}\) be a self-adjoint Dirac-type operator on the
capacity boundary sector \(\mathcal H_{\partial,q}\), with zero modes
separated as a finite kernel.  Define
\begin{align}
\zeta_q(s)&:=\operatorname{Tr}' |D_{\partial,q}|^{-s},\\
\eta_q(s)&:=\operatorname{Tr}'\!\left[
D_{\partial,q}|D_{\partial,q}|^{-s-1}\right],
\end{align}
by meromorphic continuation.  Here the prime omits zero modes.  The
Atiyah--Patodi--Singer invariant \(\eta_q(0)\) is the nonperturbative
spectral completion of the Chern--Simons phase, while
\(-\zeta_q'(0)\) is the logarithmic determinant magnitude
\cite{AtiyahPatodiSinger1975,DaiFreed1994,Witten2019Eta}.  Thus a
boundary selector can contain a real determinant term and an eta-phase
term without confusing either with the positive capacity count.

\begin{definition}[Spectral determinant/eta selector]
The admissible boundary selector may be refined by
\begin{equation}
\boxed{
\begin{aligned}
\Gamma_{\partial}^{\zeta\eta}(q)
={}&\Gamma_{\partial}^{\rm HS}(q)
+\lambda_{\zeta}\Bigl[
-\frac12\zeta_q'(0)-\Theta_\varphi(q)
\Bigr]^2\\
&+\lambda_{\eta}\sin^2\!\left[
\frac{\pi}{2}\eta_q(0)-\theta_{\rm CS}(q)
\right],
\end{aligned}}
\end{equation}
where
\begin{align}
\Theta_\varphi(q)
&:=
\ln\!\left[
\frac{\pi\varphi^{2N(q)}}{q}
\right],\\
N(q)
&:=\operatorname*{argmin}_{N\in\mathbb Z}
\left|
\log_\varphi\sqrt{\frac{q}{\pi}}-N
\right|.
\end{align}
The \(\zeta\)-term is a real spectral-action penalty.  The \(\eta\)-term
is a phase/framing penalty.  Neither term is a state count by itself.
\end{definition}

For the present capacity anchor, the nearest \(\varphi\)-shell has
\(N=292\), giving
\begin{equation}
\epsilon_\varphi
:=
\frac{\pi\varphi^{584}-q_{\rm obs}}{q_{\rm obs}}
=0.062844\ldots,
\end{equation}
and therefore
\begin{equation}
\Theta_\varphi(q_{\rm obs})
=
\ln(1+\epsilon_\varphi)
=
0.060948\ldots .
\end{equation}
The simple phase-defect target
\begin{equation}
\theta_{50}:=\ln\!\left(1+\frac{\pi}{50}\right)
=0.0609369\ldots
\end{equation}
differs by only
\begin{equation}
\boxed{
\Theta_\varphi(q_{\rm obs})-\theta_{50}
\simeq 1.15\times10^{-5}.
}
\end{equation}
This numerical proximity is not promoted to a theorem.  Its invariant
content is more precise: a microscopic boundary operator \(D_{\partial,q}\)
would solve the selector problem only if its determinant-line data force
\[
-\frac12\zeta_{q_*}'(0)\simeq\Theta_\varphi(q_*),
\qquad
\frac{\pi}{2}\eta_{q_*}(0)\simeq \theta_{\rm CS}(q_*)
\quad(\bmod\ \pi),
\]
with \(q_*\) selected without inserting \(\Lambda_{\rm obs}\).

\begin{nogotheorem}[Zeta value versus determinant selector]
The number \(\zeta_q(0)\) is not a positive horizon capacity.  It is a
regularization datum.  The admissible spectral selector is instead the
determinant-line pair
\[
\left(-\zeta_q'(0),\eta_q(0)\right).
\]
Consequently, a zeta-regularized constant can enter the boundary action
only through a determinant or anomaly phase; it cannot replace the
integer capacity \(q=\dim\mathcal H_{\partial,q}\).
\end{nogotheorem}

\begin{remark}[What this closes]
This section closes another false route.  The APS eta invariant and the
zeta determinant do not compute \(q_{\rm obs}\) alone.  They identify the
minimal spectral data that a genuine boundary Hamiltonian must possess.
The remaining problem is now sharper: construct \(D_{\partial,q}\), or an
equivalent edge-mode Hamiltonian, whose determinant-line functional has
its minimum at
\[
q_*=3.31\times10^{122}.
\]
\end{remark}


\section{Concrete edge Dirac operator and determinant-selector audit}
\label{sec:concrete-edge-dirac}

The preceding selector still contained a symbolic \(D_{\partial,q}\).  The
next invariant test is to choose the most economical boundary operator and
ask whether it selects the observed sector.  The canonical choice on the
capacity two-sphere is the \(SU(2)\)-equivariant spin-\(c\) Dirac operator
on \(S^2_R\simeq\mathbb{CP}^1\), twisted by the monopole line bundle whose
holomorphic zero modes are
\[
\mathcal H_{\partial,q}=H^0(\mathbb{CP}^1,\mathcal O(q-1)),
\qquad
\dim\mathcal H_{\partial,q}=q .
\]
This is the same one-particle cell space used in the index/capacity
construction.  It is also the standard monopole-harmonic setting for
fermions on the sphere \cite{WuYang1976,Dolan2020QHDirac}.

\begin{theorem}[Concrete monopole spin-\(c\) edge spectrum]
\label{thm:concrete-dirac-spectrum}
Let \(D_q\) be the \(SU(2)\)-equivariant spin-\(c\) Dirac operator on
\(S^2_R\) with index \(q\ge1\).  With zero modes omitted from the
determinant, the nonzero spectrum is symmetric and may be written as
\begin{equation}
\boxed{
\begin{gathered}
\lambda_{\ell,\pm}
=
\pm R^{-1}\sqrt{\ell(\ell+q)},\\
d_\ell=2\ell+q,
\qquad
\ell=1,2,\ldots .
\end{gathered}}
\end{equation}
The kernel has dimension \(q\).  Hence
\begin{equation}
\boxed{
\eta_q(0)=0
}
\end{equation}
for this symmetric operator, while the nonzero determinant is encoded by
\begin{equation}
\boxed{
\zeta_{D_q^2}(s)
=
2R^{2s}
\sum_{\ell=1}^{\infty}
(2\ell+q)\,[\ell(\ell+q)]^{-s}.
}
\end{equation}
\end{theorem}

\begin{proof}
The index statement is the Riemann--Roch/Borel--Weil count already used
above.  The nonzero monopole-Dirac spectrum is the standard
\(SU(2)\)-equivariant spectrum: each positive level has a negative partner
with the same degeneracy.  Therefore the APS eta sum over the nonzero
spectrum cancels pairwise after meromorphic continuation.  Squaring the
nonzero eigenvalues gives the displayed zeta function, with a factor of
two for the two signs.
\end{proof}

\begin{corollary}[Eta no-go for the minimal symmetric edge operator]
The minimal \(SU(2)\)-equivariant monopole Dirac operator does not supply a
nontrivial eta selector:
\[
\eta_q(0)=0 .
\]
Therefore a nonzero APS/framing phase capable of selecting \(q\) must come
from additional boundary data: a framing choice, chiral asymmetry,
interaction, spectral-flow path, or a genuinely nontrivial edge-mode
Hamiltonian.  It is not produced by the bare symmetric monopole Dirac
operator alone.
\end{corollary}

\begin{nogotheorem}[Determinant-line no-go for a natural large-\(q\) selector]
\label{nogo:determinant-alone}
The determinant of the concrete operator \(D_q\), by itself, does not give
a scheme-independent minimum at \(q\simeq10^{122}\).  After the capacity
radius relation
\[
R_q^2=\frac{q}{\pi}\ell_P^2
\]
is imposed, the determinant-line effective action is still defined up to
local boundary counterterms and renormalization conditions.  Thus it can
shift by terms built from the heat-kernel coefficients of the same
operator.  Such local terms can move, remove, or create extrema of
\(\Gamma_\partial(q)\).  Consequently the determinant data identify the
correct operator channel, but they do not select the observed sector
without an additional dimensionless boundary coupling or renormalization
condition.
\end{nogotheorem}

\begin{proof}
The zeta determinant is a renormalized determinant of an elliptic family.
Changing the subtraction scheme or determinant-line metric adds local
functionals of the boundary geometry and gauge curvature.  On the
homogeneous capacity sphere these local invariants reduce to functions of
the integer \(q\) after \(R_q^2=q\ell_P^2/\pi\) is imposed.  Since these
additive local terms enter the same one-dimensional variational problem
for \(q\), the position of an extremum is not fixed by the bare determinant
alone.  A genuine large hierarchy must therefore be generated by a
renormalization-group invariant, for example dimensional transmutation, or
by an independently specified edge-mode Hamiltonian.
\end{proof}

The transmutation form is therefore not optional decoration; it is the
surviving mathematical mechanism for a large sector:
\begin{equation}
\boxed{
q_\star
=
\Xi\exp\!\left(\frac{16\pi^2}{b\,g_\partial^2}\right).
}
\end{equation}
For the present capacity target,
\begin{equation}
\boxed{
b\,g_\partial^2
=
\frac{16\pi^2}{\ln q_{\rm obs}}
=
0.55976\pm0.00003 .
}
\end{equation}
A small constants audit is useful but not decisive:
\begin{center}
\begin{tabular}{@{}lcc@{}}
\toprule
candidate & value & relative error\\
\midrule
\(e^{-\gamma_{\rm E}}\) & \(0.56145948\) & \(3.04\times10^{-3}\)\\
\(\sqrt5/4\) & \(0.55901699\) & \(-1.32\times10^{-3}\)\\
\(1/\sqrt\pi\) & \(0.56418958\) & \(7.92\times10^{-3}\)\\
\(\ln\varphi+\pi/40\) & \(0.55975164\) & \(-8.26\times10^{-6}\)\\
\bottomrule
\end{tabular}
\end{center}
None of these equalities is exact.  The last row is deliberately included
as a warning: post-hoc combinations of simple constants can be made very close,
but closeness is not an invariant derivation.  If one uses \(e^{-\gamma_{\rm E}}\),
the resulting \(q\) is \(1.41\times10^{122}\), not the observed
\(3.31\times10^{122}\).  If one uses \(\sqrt5/4\), the resulting \(q\) is
\(4.80\times10^{122}\).  The audit therefore closes another false route:
nearby famous constants do not solve the selector problem.  The boundary
Hamiltonian must compute \(b g_\partial^2\) or an equivalent determinant
condition.

\begin{figure}[!htbp]
\centering
\includegraphics[width=\columnwidth]{fig_edge_dirac_selector.png}
\caption{Concrete edge-operator selector audit.  The target
\(b g_\partial^2=16\pi^2/\ln q_{\rm obs}\) is compared with simple
dimensionless constants.  The closeness of \(e^{-\gamma_{\rm E}}\) and
\(\sqrt5/4\) is a diagnostic clue, not a theorem: neither value reproduces
the observed capacity sector.}
\label{fig:edge-dirac-selector}
\end{figure}

\begin{remark}[What this closes]
The concrete monopole Dirac operator supplies the required edge spectrum
and confirms the index \(q\), but its symmetric eta invariant vanishes and
its determinant does not select a hierarchy without renormalization data.
The remaining target is now exact:
\[
\boxed{
\begin{gathered}
\text{derive } b\,g_\partial^2=0.55976\pm0.00003\\
\text{from the boundary Hamiltonian.}
\end{gathered}}
\]
This is narrower than deriving \(\Lambda\) directly and sharper than any
purely topological, vortex, Hopf, or zeta-value explanation.
\end{remark}



\section{Invariant scale atlas: configuration, phase, field, and capacity spaces}\label{sec:scale-atlas}

The preceding sections reduce the remaining problem to a boundary selector for
the integer \(q=N_\partial\).  This section builds the common scale map on
which every admissible candidate must live.  The purpose is not to add another
analogy, but to identify which numerical similarities are invariant and which
are merely scale coincidences.

\begin{definition}[Four-coordinate invariant scale map]
For a phenomenon with a characteristic length \(L\), energy \(E\), field scale
\(\mathcal E\), and topological integer \(n_{\rm top}\), define
\begin{equation}
\boxed{
\mathfrak X
=
\left(
x_L,\,
N_A,\,
N_{\rm ph},\,
n_\varphi,\,
n_{\rm top},\,
\mathcal E/\mathcal E_{\rm crit}
\right)
}
\end{equation}
with
\begin{equation}
x_L:=\ln(L/\ell_P),\qquad
N_A(L):=\frac{4\pi L^2}{4\ell_P^2}
       =\pi\left(\frac{L}{\ell_P}\right)^2,
\end{equation}
\begin{equation}
N_{\rm ph}(E,L):=\frac{2\pi E L}{\hbar c},
\qquad
n_\varphi(L):=\frac{\ln(L/\ell_P)}{\ln\varphi}.
\end{equation}
Here \(N_A\) is the spherical area-capacity coordinate, \(N_{\rm ph}\)
is the Bekenstein phase/action coordinate, and \(n_\varphi\) is only a
scale registry.  A proposed selector may use these coordinates only through
dimensionless, pivot-independent combinations.
\end{definition}

\begin{theorem}[Area--phase--entropy collapse for a gravitational screen]
Let a spherical gravitational screen of radius \(R\) carry the horizon/Misner--
Sharp energy
\begin{equation}
E_G(R)=\frac{c^4R}{2G}.
\end{equation}
Then its Bekenstein phase coordinate equals its area capacity:
\begin{equation}
\boxed{
N_{\rm ph}(E_G,R)
=
\frac{2\pi E_G R}{\hbar c}
=
\pi\frac{R^2}{\ell_P^2}
=
N_A(R).
}
\end{equation}
For a de Sitter horizon \(R_\Lambda=\sqrt{3/\Lambda}\), this is exactly
\begin{equation}
\boxed{
N_{\rm ph}=N_A=\frac{3\pi}{\Lambda\ell_P^2}=N_\partial.
}
\end{equation}
Thus the configuration-space area count, phase-space action count, and
horizon entropy count collapse to the same invariant number on a gravitational
horizon \cite{Bekenstein1973,GibbonsHawking1977}.
\end{theorem}

\begin{proof}
Using \(\ell_P^2=\hbar G/c^3\),
\[
\frac{2\pi E_G R}{\hbar c}
=
\frac{2\pi}{\hbar c}\frac{c^4R}{2G}R
=
\pi\frac{c^3R^2}{\hbar G}
=
\pi\frac{R^2}{\ell_P^2}.
\]
For \(R_\Lambda^2=3/\Lambda\), the last expression becomes
\(3\pi/(\Lambda\ell_P^2)\).
\end{proof}


\begin{definition}[Radius-locked \texorpdfstring{\(S^3\)}{S3} reference volume]
Let \(R_H\) be the causal horizon radius and let
\[
V_{S^3}(R_3)=2\pi^2R_3^3
\]
be the volume of a round spatial three-sphere of radius \(R_3\).  The
\emph{radius-locked} reference geometry is the special case
\[
\boxed{
R_3=R_H=R_\Lambda .
}
\]
Without this additional identification, the causal horizon determines only
\(R_H\), while the total \(S^3\) spatial volume requires an independent
curvature/topology datum.
\end{definition}

\begin{theorem}[\texorpdfstring{\(S^3\)}{S3} volume--horizon density bridge]
\label{thm:s3-volume-density-bridge}
Assume the radius-locked reference geometry.  Let
\[
\mathcal A_H=4\pi R^2,\qquad
V_{S^3}=2\pi^2R^3,\qquad
E_H=\frac{c^4R}{2G}.
\]
Then the dark-energy density associated with the de Sitter scale is
equivalently
\begin{equation}
\boxed{
\rho_\Lambda
=
\frac{\Lambda c^4}{8\pi G}
=
\frac{3c^4}{8\pi G R^2}
=
\frac{3\pi}{2}\frac{E_H}{V_{S^3}} .
}
\end{equation}
It may also be written purely from the radius-locked \(S^3\) volume,
\begin{equation}
\boxed{
\rho_\Lambda
=
\frac{3c^4}{8\pi G}
\left(\frac{2\pi^2}{V_{S^3}}\right)^{2/3},
}
\end{equation}
or from the paired horizon area and \(S^3\) volume,
\begin{equation}
\boxed{
\rho_\Lambda
=
\frac{3\pi c^4}{32G}
\frac{\mathcal A_H^2}{V_{S^3}^{\,2}} .
}
\end{equation}
Conversely, a measured energy density determines the radius-locked reference
volume by
\begin{equation}
\boxed{
\begin{aligned}
R_\rho&=
\left(\frac{3c^4}{8\pi G\rho_\Lambda}\right)^{1/2},\\
V_{S^3}(\rho_\Lambda)
&=
2\pi^2
\left(\frac{3c^4}{8\pi G\rho_\Lambda}\right)^{3/2}.
\end{aligned}}
\end{equation}
In capacity variables,
\begin{equation}
\boxed{
\begin{aligned}
R&=\ell_P\sqrt{\frac{N_\partial}{\pi}},\\
V_{S^3}&=2\sqrt{\pi}\,\ell_P^3N_\partial^{3/2},\\
\rho_\Lambda&=\frac{3c^4}{8G\ell_P^2N_\partial}.
\end{aligned}}
\end{equation}
\end{theorem}

\begin{proof}
The de Sitter relation is \(\Lambda=3/R^2\), hence
\[
\rho_\Lambda=\Lambda c^4/(8\pi G)=3c^4/(8\pi G R^2).
\]
Using \(E_H=c^4R/(2G)\) and \(V_{S^3}=2\pi^2R^3\),
\[
\frac{3\pi}{2}\frac{E_H}{V_{S^3}}
=
\frac{3\pi}{2}
\frac{c^4R/(2G)}{2\pi^2R^3}
=
\frac{3c^4}{8\pi G R^2}.
\]
Solving \(V_{S^3}=2\pi^2R^3\) for \(R\), or using
\(\mathcal A_H=4\pi R^2\), gives the two alternative expressions.  Solving
the same density formula for \(R\) gives the inverse map.  The capacity form
follows from \(N_\partial=\pi R^2/\ell_P^2\).
\end{proof}

\begin{proposition}[Numerical check of the bridge]
With the capacity anchor used in this manuscript,
\[
H_0=67.4\,{\rm km\,s^{-1}Mpc^{-1}},\qquad
\Omega_\Lambda=0.685,
\]
one obtains
\[
\boxed{
\begin{gathered}
\Lambda_0=1.0909\times10^{-52}\,{\rm m}^{-2},\\
R_\Lambda=1.6583\times10^{26}\,{\rm m},\\
\rho_\Lambda=5.253\times10^{-10}\,{\rm J\,m^{-3}}
=5.845\times10^{-27}\,{\rm kg\,m^{-3}},\\
V_{S^3}=9.002\times10^{79}\,{\rm m^3},\\
E_H=1.0035\times10^{70}\,{\rm J}.
\end{gathered}}
\]
The volume energy is
\[
\rho_\Lambda V_{S^3}=4.729\times10^{70}\,{\rm J}
=
\frac{3\pi}{2}E_H,
\]
as required by \cref{thm:s3-volume-density-bridge}.
\end{proposition}

\begin{remark}[What the inverse calculation means]
The inverse formula does not prove that the whole universe has the
radius-locked volume \(V_{S^3}(\rho_\Lambda)\).  It computes the unique
round \(S^3\) reference volume that would share the de Sitter horizon radius.
In a general closed FLRW model the spatial curvature radius \(R_3(a)\) and
the causal horizon radius \(R_H\) are different data.  Therefore the
\(S^3\)-volume bridge is an exact consistency identity under radius lock, not
an independent measurement of the total spatial volume unless the lock is
derived or observed.
\end{remark}



\begin{corollary}[Dark-energy density length as the logarithmic midpoint]
Let
\begin{equation}
\rho_\Lambda=\frac{\Lambda c^4}{8\pi G},
\qquad
\ell_\rho:=\left(\frac{\hbar c}{\rho_\Lambda}\right)^{1/4}.
\end{equation}
Then
\begin{equation}
\boxed{
\ell_\rho
=
\left(\frac{8\pi}{3}\right)^{1/4}
\sqrt{\ell_P R_\Lambda}.
}
\end{equation}
Hence the dark-energy density length is the geometric mean of the Planck
length and the de Sitter radius, up to the fixed factor \((8\pi/3)^{1/4}\).
In \(\varphi\)-registry coordinates,
\begin{equation}
\boxed{
n_\varphi(\ell_\rho)
=
\frac12 n_\varphi(R_\Lambda)
+
\log_\varphi\!\left[(8\pi/3)^{1/4}\right].
}
\end{equation}
For the Planck-2018/CODATA numerical anchor used in this paper,
\[
R_\Lambda=1.6583\times10^{26}{\rm m},\quad
N_\partial=3.3072\times10^{122},
\]
and
\[
\boxed{
\begin{gathered}
\ell_\rho=8.8078\times10^{-5}{\rm m},\\
n_\varphi(R_\Lambda)=291.9367,\qquad
n_\varphi(\ell_\rho)=147.0726.
\end{gathered}}
\]
Thus the observed dark-energy density scale lies almost exactly on the
middle logarithmic scale between \(\ell_P\) and \(R_\Lambda\), shifted by the
universal density-normalization constant.
\end{corollary}

\begin{figure*}[!htbp]
\centering
\includegraphics[width=0.92\textwidth]{fig_scale_atlas.png}
\caption{Configuration-space invariant scale atlas.  The de Sitter horizon,
the \(\varphi\)-shell \(N=292\), and the dark-energy density length occupy
distinguished positions on the same logarithmic axis.  The density length is
the geometric midpoint scale between \(\ell_P\) and \(R_\Lambda\), not an
independent fundamental length.}
\label{fig:scale-atlas}
\end{figure*}

\begin{figure}[!htbp]
\centering
\includegraphics[width=\linewidth]{fig_phase_capacity_collapse.png}
\caption{Phase-space/area-capacity collapse.  For any gravitational screen
with \(E=c^4R/(2G)\), the Bekenstein action coordinate \(2\pi ER/\hbar c\)
equals the area capacity \(\pi R^2/\ell_P^2\).  The de Sitter capacity is
therefore simultaneously an entropy, area, and phase-action count.}
\label{fig:phase-capacity}
\end{figure}

\begin{table}[t]
\caption{Selected invariant scale coordinates.  \(N_A=\pi(L/\ell_P)^2\).}
\label{tab:scale-atlas}
\begin{ruledtabular}
\begin{tabular}{lccc}
scale & \(L\,[{\rm m}]\) & \(n_\varphi\) & \(\log_{10}N_A\)\\
\hline
\(\ell_P\) & \(1.6163\times10^{-35}\) & \(0\) & \(0.497\)\\
electron \(\bar\lambda_C\) & \(3.8616\times10^{-13}\) & \(107.079\) & \(45.254\)\\
\(\ell_\rho\) & \(8.8078\times10^{-5}\) & \(147.073\) & \(61.970\)\\
\(R_\Lambda\) & \(1.6583\times10^{26}\) & \(291.937\) & \(122.519\)\\
\(\ell_P\varphi^{292}\) & \(1.7096\times10^{26}\) & \(292.000\) & \(122.546\)
\end{tabular}
\end{ruledtabular}
\end{table}

\begin{criterion}[Scale-atlas admissibility]
A proposed boundary selector for \(q=N_\partial\) must pass the following
scale-atlas tests:
\begin{enumerate}
\item it must reproduce the collapse \(N_{\rm ph}=N_A=N_\partial\) on a
      gravitational horizon;
\item it must respect the midpoint relation
      \(\ell_\rho=(8\pi/3)^{1/4}\sqrt{\ell_P R_\Lambda}\);
\item it must not identify Schwinger, Ginzburg--Landau, Hopf, or Dirac
      scales with \(q\) unless a dimensionless invariant map is supplied;
\item it must distinguish registry clues, such as the near \(\varphi\)-shell
      relation \(n_\varphi(R_\Lambda)\simeq292\), from selector theorems.
\end{enumerate}
\end{criterion}

\begin{remark}[What the scale atlas proves]
The atlas reveals a real invariant collapse:
\[
\boxed{
\begin{gathered}
\text{geometry area}
=
\text{Bekenstein phase action}\\
=
\text{horizon entropy}
\end{gathered}}
\]
for gravitational screens.  It also reveals a real midpoint law for the dark
energy density scale.  It does not by itself select \(q_{\rm obs}\).  Therefore
the remaining selector problem is not to invent another scale, but to derive
which edge-mode Hamiltonian chooses the gravitational screen radius, or
equivalently the sector \(q\), without using \(\Lambda_{\rm obs}\) as input.
\end{remark}




\section{Mechanism-scale atlas and necessary selector conditions}
\label{sec:mechanism-scale-atlas}

The previous scale atlas compared lengths, fields, capacities, and phase-space
counts.  This section adds a different axis: the mechanism by which a discrete
integer is selected from a set of admissible integers.  This is the correct
place to use analogies with boxes, oscillators, Hall systems, black holes, and
dimensional transmutation.  The analogies are not proofs of the de Sitter
sector.  They are filters: a proposed \(q\)-selector must have the right
invariant type.

\begin{definition}[Selector mechanism level]
A selector mechanism for an integer \(q\) is assigned to the first applicable
level in the following hierarchy:
\[
\begin{array}{ccl}
0&:&\text{boundary-condition selector},\\
1&:&\text{self-adjoint spectral selector},\\
2&:&\text{topological or index selector},\\
3&:&\text{dimensional-transmutation selector},\\
4&:&\text{entropic or thermodynamic selector},\\
5&:&\text{self-referential boundary-capacity selector}.
\end{array}
\]
Levels \(0\)--\(4\) are familiar in ordinary quantum and statistical systems.
The de Sitter capacity problem belongs to level \(5\), because the system must
select the boundary sector that defines its own horizon size.
\end{definition}

\begin{remark}[What is corrected from the raw analogy list]
Several tempting identifications are not invariant.  Hydrogen degeneracy
\(n^2\) is a spectral degeneracy, not an index in the sense of
Atiyah--Singer.  The disjoint sector space
\[
\mathcal P_{\rm all}:=\bigsqcup_{q\ge1}\mathcal C_q
\]
is a countable disjoint union of finite-dimensional sectors, not an uncountable
union.  Finally, the code-cell dimension
\(\dim H^0(\CP^1,\mathcal O(q-1))=q\) is not by itself the full entropy of the
horizon.  The entropy statement used in this manuscript is
\[
\log\dim\mathcal C_q=q.
\]
Thus the index supplies admissible integer labels, while the edge-mode code
space supplies the capacity.
\end{remark}

\begin{definition}[Two-sided sector invariant]
For each admissible sector,
\[
q=\operatorname{ind}\bar\partial_{\mathcal O(q-1)}
=\dim H^0(\CP^1,\mathcal O(q-1))
=\frac{k_{\rm CS}}2
=\frac{A_q}{4\ell_P^2}
\]
is the same integer when seen from the outside as a sector label and from the
inside as the horizon capacity.  This is the two-sided invariant.  It is not
yet a selection theorem for \(q_{\rm obs}\).
\end{definition}

\begin{theorem}[Necessary conditions for a de Sitter capacity selector]
A boundary theory that truly selects \(q_{\rm obs}\) without using
\(\Lambda_{\rm obs}\) as an input must satisfy the following conditions.
\begin{enumerate}
\item \emph{Discrete admissibility:} \(q\in\mathbb Z_{>0}\).
\item \emph{Topological stability:} small metric or coupling deformations do
not change the selected sector unless a gap closes or a sector transition
occurs.
\item \emph{Classical--quantum consistency:}
\[
\Lambda_{\rm classical}=\Lambda_{\rm quantum}
=\frac{3\pi}{q\ell_P^2}.
\]
\item \emph{Transmutation or equivalent hierarchy generation:} since
\(q_{\rm obs}\sim10^{122}\), an \(O(1)\) microscopic input must generate a
large logarithm,
\[
\ln q_{\rm obs}\simeq 282.1.
\]
Equivalently, for an asymptotically free boundary coupling,
\[
q\sim
\Xi\exp\!\left(\frac{16\pi^2}{b g_\partial^2}\right).
\]
\item \emph{Exchange covariance:} if \(q\) changes along a cosmological
history, the effective bulk stress tensor must obey a covariant exchange law,
not a varying Bianchi integration constant.
\end{enumerate}
\end{theorem}

\begin{proof}
The first condition follows from the index/Chern--Simons/Bradlow readings of
the sector.  The second is the homotopy invariance of the index away from
singular sector changes.  The third is the defining capacity relation.  The
fourth follows from the absence of any intermediate algebraic scale between
\(O(1)\) microscopic data and \(q_{\rm obs}\); a power law in \(O(1)\)
parameters cannot naturally produce \(10^{122}\), while dimensional
transmutation can.  The fifth is forced by the Bianchi identity: a strictly
constant integration mode cannot be the same object as a time-dependent
effective dark-energy density.
\end{proof}

\begin{criterion}[What a final selector must compute]
The remaining unknown is a concrete edge-mode functional
\[
\Gamma_\partial(q)=-\log\operatorname{Tr}_{\mathcal C_q}
e^{-\beta_q H_{\partial,q}}
\]
with a discrete stable minimum
\[
q_*=\arg\min_{q\in\mathbb Z_{>0}}\Gamma_\partial(q).
\]
It must compute either
\[
b g_\partial^2=0.55976\pm0.00003 ,
\]
or an equivalent determinant/phase/code-space condition, without inserting
\(\Lambda_{\rm obs}\), \(q_{\rm obs}\), or the observed horizon radius as
input data.
\end{criterion}


\begin{center}
\includegraphics[width=\linewidth]{fig_mechanism_scale.png}

{\small\textbf{Figure.} Mechanism-scale atlas.  The horizontal axis is \(\log_{10}q\), while the vertical bands encode the kind of selection mechanism.  Ordinary quantum systems select small integers through boundary conditions or spectra; Hall systems and index problems select topological integers; QCD-like theories produce large hierarchies by transmutation.  The de Sitter sector is a level-5 self-referential boundary-capacity problem.}
\end{center}



\begin{center}
\includegraphics[width=\linewidth]{fig_selector_conditions.png}

{\small\textbf{Figure.} Necessary selector conditions.  Index quantization, entropy capacity, Chern--Simons level locking, Bradlow packing, and phase-space collapse are already admissibility constraints.  The missing input is a concrete edge-mode Hamiltonian whose effective action has a stable minimum at the observed sector without using the observed cosmological constant as an input.}
\end{center}



\section{Universal connector sweep across dimensions and phase spaces}\label{sec:universal-connector-sweep}

The preceding mechanism-scale atlas leaves one narrow problem: construct an
edge Hamiltonian whose sector effective action selects \(q_{\rm obs}\) without
using \(\Lambda_{\rm obs}\) as an input.  I now add the all-dimensional
connector sweep.  It is deliberately a filtering statement.  The point is not
to accumulate analogies, but to sort every candidate invariant into the one
role it can actually play.

\begin{definition}[Universal connector test]
A proposed connector between a known physical mechanism and the de Sitter
capacity sector is admitted only if it satisfies all four conditions:
\begin{enumerate}
\item it is dimensionless or becomes dimensionless after division by
the natural Planck, curvature, or flux unit;
\item it is invariant under gauge choice, pivot rescaling, and smooth
deformation within a fixed topological sector;
\item it has a definite role in \(\Gamma_\partial(q)\): quantizer, bound,
obstruction, energy cost, or selector;
\item it does not use \(\Lambda_{\rm obs}\), \(R_\Lambda\), or \(q_{\rm obs}\)
as an input.
\end{enumerate}
Numerical proximity alone fails the test.
\end{definition}

\begin{proposition}[Division-algebra Hopf ladder]\label{prop:hopf-ladder}
The sphere-to-sphere Hopf fibrations supplied by the real normed division
algebras form the ladder
\begin{gather}
S^0\to S^1\to S^1,\\
S^1\to S^3\to S^2,\\
S^3\to S^7\to S^4,\\
S^7\to S^{15}\to S^8 .
\end{gather}
The \(S^3\to S^2\) member is therefore the unique complex Hopf connector
appropriate to three compactified spatial dimensions.  Its charge is the
Chern--Simons/helicity number
\begin{equation}
Q_H=\frac1{(4\pi)^2}\int_{S^3}A\wedge F,\qquad F=\dd A .
\end{equation}
The higher members supply the corresponding quaternionic and octonionic
admissibility channels, not alternative proofs of the observed
\(q_{\rm obs}\).  The division-algebra result is classical; here its role is
only to prevent an arbitrary proliferation of Hopf-like connectors
\cite{DivisionHopfKTheory,GluckWarnerZiller1986}.
\end{proposition}

\begin{corollary}[All-dimensional connector roles]
Every admitted connector found in the sweep falls into exactly one of five
roles:
\begin{enumerate}
\item sector quantizer: \(q=\operatorname{index}D\) or
\(q=(2\pi)^{-1}\int F+1\);
\item saturation bound: \(N_{\rm Bek}=2\pi ER/(\hbar c)\le
N_A=A/(4\ell_P^2)\);
\item obstruction: \(\Omega_{\rm cob}=0\), or a Chern--Weil current is
gauged, broken, or edge-completed;
\item energy cost: \(E_{\rm Hopf}\ge C|Q_H|^{3/4}\) in Faddeev--Skyrme type
models;
\item selector:
\begin{equation}
q_*\sim \Xi\exp\!\left[\frac{16\pi^2}{b\,g_\partial^2}\right],
\end{equation}
or an exactly equivalent edge-mode Hamiltonian minimum.
\end{enumerate}
Only the last role can produce a hierarchy of order \(10^{122}\) from
order-one input data.  The other roles restrict the admissible sector but do
not choose it.
\end{corollary}

\begin{theorem}[Phase-area-entropy saturation as the two-sided invariant]
For a four-dimensional gravitational screen of radius \(R\) whose gravitational
energy is
\begin{equation}
E_G(R)=\frac{c^4R}{2G},
\end{equation}
the Bekenstein phase count equals the area capacity:
\begin{equation}
\boxed{
N_{\rm ph}(R):=\frac{2\pi E_G(R)R}{\hbar c}
=
\pi\frac{R^2}{\ell_P^2}
=
\frac{A(R)}{4\ell_P^2}
=:N_A(R).
}
\end{equation}
Thus the same number is seen from the interior as a phase-space action count
and from the boundary as an entropy-area count.  This is the invariant
collapse that any acceptable \(\Gamma_\partial(q)\) must reproduce
\cite{Bekenstein1981,Susskind1995,Bousso2002}.
\end{theorem}

\begin{proof}
Substitute \(E_G=c^4R/(2G)\) into \(2\pi ER/(\hbar c)\):
\[
\frac{2\pi}{\hbar c}\frac{c^4R}{2G}R
=
\pi\frac{c^3R^2}{\hbar G}
=
\pi\frac{R^2}{\ell_P^2}.
\]
Since \(A=4\pi R^2\), the last expression is \(A/(4\ell_P^2)\).
\end{proof}

\begin{criterion}[Chern--Weil and cobordism admissibility]
Chern--Weil currents and cobordism classes cannot be treated as silent exact
global labels in quantum gravity.  They must either be gauged, broken by
defects, cancelled by anomaly inflow, or represented by explicit edge modes.
This is an admissibility constraint on \(\Gamma_\partial\), not a selector for
the observed value of \(q\) by itself \cite{HeidenreichEtAl2021ChernWeil,McNamaraVafa2019}.
\end{criterion}

\begin{criterion}[Universal selector normal form]
After the sweep, a genuine boundary selector must have the schematic form
\begin{equation}
\boxed{
\begin{aligned}
\Gamma_\partial(q)
={}&
\Gamma_{\rm run}(q;g_\partial)
+\Gamma_{\rm det}(q;D_{\partial,q})\\
&+\Gamma_{\rm anom}(q)
+\Gamma_{\rm sat}(q)
+\Gamma_{\rm exch}(q).
\end{aligned}}
\end{equation}
Here \(\Gamma_{\rm run}\) is the only accepted source of an exponential
hierarchy unless an explicit code Hamiltonian directly computes
\(\log\dim\mathcal C_q=q\).  The target is therefore not a new analogy but the
single number
\begin{equation}
b\,g_\partial^2
=
\frac{16\pi^2}{\ln q_*},
\qquad
q_*=q_{\rm obs},
\end{equation}
or an exactly equivalent determinant/edge-code condition.
\end{criterion}


\begin{center}
\includegraphics[width=\linewidth]{fig_universal_connector.png}

{\small\textbf{Figure.} Universal connector sweep.  Left: the sphere-to-sphere Hopf ladder is discrete and division-algebraic.  Right: admitted invariants have different logical roles.  Index, saturation, anomaly, and energy-cost data restrict admissible sectors; only an RG/edge-Hamiltonian selector can pick the enormous observed sector from order-one input data.}
\end{center}


The practical conclusion is sharp.  I found no all-dimensional invariant that
selects \(q_{\rm obs}\) alone.  The sweep instead proves a stronger
organization: all reliable universal structures enter \(\Gamma_\partial\) as
quantizers, obstructions, costs, or saturation tests, while the actual
selection of \(q\sim10^{122}\) still requires boundary dynamics.


\section{Fibonacci, spinors, Zeckendorf coding, and golden-ratio admissibility}

The golden ratio enters the present atlas only when it has an invariant
role.  I separate three such roles.  First, it is the Perron--Frobenius
eigenvalue of the Fibonacci substitution matrix.  Second, it is the sharp
Diophantine obstruction to global rational commensurability.  Third, it
is the quantum dimension of the Fibonacci anyon.  These are genuine
mathematical and physical appearances of \(\varphi\).  None of them, by
itself, selects the observed capacity sector.

\begin{definition}[Fibonacci matrix and scale eigenvalues]
Let
\begin{equation}
\mathcal F:=
\begin{pmatrix}
1&1\\
1&0
\end{pmatrix}.
\end{equation}
Then
\begin{align}
\det\mathcal F&=-1,\\
\operatorname{Spec}(\mathcal F)&=
\{\varphi,-\varphi^{-1}\},\\
\varphi&=\frac{1+\sqrt5}{2}.
\end{align}
Thus \(\log\varphi\) is the natural scale entropy per Fibonacci
substitution step, while \(\varphi^{-1}\) is the stable contraction
ratio.
\end{definition}

\begin{theorem}[Zeckendorf coding is a unique integer code, not a selector]
Every positive integer \(m\) has a unique expansion
\begin{equation}
m=\sum_i F_{n_i},
\qquad
n_{i+1}\ge n_i+2,
\end{equation}
as a sum of nonconsecutive Fibonacci numbers \cite{Zeckendorf1972}.
Consequently every integer shell label has a canonical Fibonacci code.
For the present capacity registry,
\begin{equation}
n_\varphi(q):=
\log_\varphi\sqrt{\frac{q}{\pi}},
\end{equation}
the observed capacity target gives
\begin{equation}
n_\varphi(q_{\rm obs})=291.936672\ldots ,
\qquad
N_\varphi=292,
\end{equation}
and the integer shell has the Zeckendorf code
\begin{equation}
292=233+55+3+1
=F_{13}+F_{10}+F_4+F_2 .
\end{equation}
This is a canonical registry label.  It is not a derivation of
\(q_{\rm obs}\).
\end{theorem}

\begin{proof}
The first statement is Zeckendorf's theorem.  The displayed shell value
is obtained by applying the dimensionless registry \(n_\varphi\) to the
capacity target already used in the capacity audit.  The decomposition
\(292=233+55+3+1\) uses no consecutive Fibonacci indices and is therefore
the unique Zeckendorf representation of \(292\).
\end{proof}

\begin{proposition}[Golden-ratio incommensurability]
The continued fraction of \(\varphi\) is
\begin{equation}
\varphi=[1;1,1,1,\ldots].
\end{equation}
Equivalently, its rational approximants are the Fibonacci ratios
\(F_{n+1}/F_n\).  In the sharp Hurwitz bound, the constant \(\sqrt5\)
is optimal and the golden ratio is the extremal obstruction to uniformly
improving the rational approximation inequality
\cite{HardyWright2008}.  Thus a \(\varphi\)-registry is the most rigid
one-dimensional irrational registry against global rational phase
locking.
\end{proposition}

\begin{remark}[No-overclaim on standing waves]
The preceding proposition does not say that a Fibonacci or golden-ratio
medium has no modes or no Bragg peaks.  Fibonacci quasicrystals have
sharp diffraction and critical spectral structure \cite{Senechal1995,
Jagannathan2021}.  The precise statement is weaker and invariant:
\(\varphi\) gives maximal resistance to exact rational commensuration in
one dimension.  It frustrates full-spectrum periodic locking; it does
not erase all spectral response.
\end{remark}

\begin{definition}[Spinor and Fibonacci-anyon channel]
The spin group in three spatial dimensions is
\begin{equation}
\operatorname{Spin}(3)\simeq SU(2).
\end{equation}
Spinors therefore enter the boundary atlas through the double cover of
rotations and through possible \(SU(2)\)-equivariant edge sectors.
A distinct and genuinely golden quantum channel is the Fibonacci anyon
fusion rule
\begin{equation}
\tau\times\tau=1+\tau .
\end{equation}
If \(d_\tau\) is the quantum dimension of \(\tau\), then
\begin{equation}
d_\tau^2=1+d_\tau,
\qquad
d_\tau=\varphi .
\end{equation}
The associated fusion spaces grow asymptotically like powers of
\(\varphi\), and Fibonacci anyons are a standard non-Abelian topological
quantum-computing model \cite{NayakEtAl2008}.
\end{definition}

\begin{proposition}[Fibonacci fusion as a boundary-code registry]
If a boundary code sector grows as
\begin{equation}
\dim\mathcal C_m \sim C\,\varphi^m,
\end{equation}
then matching to a capacity sector \(q\) fixes only a code length
\begin{equation}
m(q)=\frac{\ln q-\ln C}{\ln\varphi}
\end{equation}
up to subleading topological and boundary constraints.  For the
\(\varphi\)-shell used here,
\begin{equation}
q_\varphi(N)=\pi\varphi^{2N},
\qquad
N=292,
\end{equation}
the corresponding fusion-growth length is
\begin{equation}
m=2N=584 .
\end{equation}
This converts the near shell relation into a concrete code-length
target.  It still does not determine why Nature selects that length.
\end{proposition}

\begin{criterion}[Fibonacci admissibility term]
A boundary selector may contain a Fibonacci/spinor code term only in the
following invariant form:
\begin{align}
\Gamma_{\partial}^{\rm Fib}(q,m)
={}&
\lambda_{\rm sh}\,
d_\varphi(q)^2
+\lambda_{\rm code}
\left[
\ln q-\ln\dim\mathcal C_m
\right]^2,\\
d_\varphi(q)
:={}&
\operatorname{dist}\!\left(
\log_\varphi\sqrt{q/\pi},\mathbb Z
\right).
\end{align}
It is an admissibility term, not a standalone selector.  A true solution
must derive \(m\), \(C\), and the couplings from the boundary edge
Hamiltonian, not fit them after observing \(\Lambda\).
\end{criterion}

\begin{figure}[!htbp]
\centering
\includegraphics[width=\columnwidth]{fig_fibonacci_spinor_phi_atlas.png}
\caption{
Fibonacci/golden-ratio registry used as an admissibility diagnostic.
The large capacity target lies near the integer shell
\(N=292\), while the relative residual is close to \(\pi/50\).
The plot is deliberately diagnostic, not a proof: the selector action
must derive the phase defect rather than insert it.
}
\label{fig:fibonacci-spinor-atlas}
\end{figure}

\begin{remark}[Admitted and rejected golden-ratio links]
The admitted links are
\begin{align}
\mathcal F&\to\{\varphi,-\varphi^{-1}\},\\
&\text{Zeckendorf uniqueness},\\
&\text{Hurwitz extremality},\\
d_\tau&=\varphi .
\end{align}
The rejected claim is stronger: \(\varphi\) does not by itself select
\(q_{\rm obs}\), does not prove the global topology of the universe, and
does not remove all spectral modes.  Its role in this paper is as a
self-similar, maximally incommensurate, spinorial/topological code
registry that any final boundary selector may exploit, but must still
derive.
\end{remark}




\section{Fibonacci renormalization clock and Diophantine selector audit}

The golden-ratio layer must be treated as an admissibility and registry
mechanism, not as a completed selector.  I therefore separate three facts:
a substitution matrix can generate the natural logarithmic clock, Zeckendorf
coding gives a unique integer registry, and the Diophantine extremality of
\(\varphi\) frustrates rational commensuration.  None of these three facts
selects \(q_{\rm obs}\) without an additional boundary Hamiltonian.

\begin{definition}[Fibonacci transfer clock]
Let
\begin{equation}
{\mathcal F}:=
\begin{pmatrix}1&1\\1&0\end{pmatrix}.
\end{equation}
Then
\begin{equation}
\operatorname{Spec}({\mathcal F})=\{\varphi,-\varphi^{-1}\},
\qquad
\det{\mathcal F}=-1.
\end{equation}
A boundary substitution clock generated by \({\mathcal F}\) has leading
growth \(\varphi^N\).  A two-direction horizon registry therefore has the
natural trial scale
\begin{equation}
q_\varphi(N)=\pi\,\varphi^{2N}.
\end{equation}
\end{definition}

\begin{proposition}[Fibonacci clock no-go]
The Fibonacci transfer clock quantizes the logarithmic registry but does
not select the observed capacity sector.  More precisely, the map
\[
N\longmapsto q_\varphi(N)=\pi\varphi^{2N}
\]
assigns an admissible ordered ladder to every \(N\in\mathbb Z\), and the
substitution dynamics contains no distinguished value of \(N\).  Hence
\(\varphi\)-growth alone cannot prove
\[
q_*=q_{\rm obs}.
\]
A selector term fixing a stopping shell or phase defect is still required.
\end{proposition}

\begin{proof}
The matrix \({\mathcal F}\) is primitive, and
\({\mathcal F}^N\) is defined for every integer \(N\ge0\).  Its
Perron--Frobenius eigenvalue gives the growth rate, while its powers give
a whole orbit of admissible shells.  Multiplication by another power of
\(\varphi^2\) moves from one shell to the next and is not forbidden by the
substitution rule.  Therefore the substitution rule supplies the clock,
not the stopping condition.
\end{proof}

\begin{proposition}[Zeckendorf shell checksum]
For the observed capacity normalization used here,
\[
n_\partial(q_{\rm obs})
=
\log_\varphi\sqrt{\frac{q_{\rm obs}}{\pi}}
=
291.936672\ldots .
\]
The nearest integer shell is \(N=292\).  Zeckendorf's theorem gives the
unique nonconsecutive Fibonacci representation
\begin{equation}
\boxed{
292=233+55+3+1=F_{13}+F_{10}+F_4+F_2 .
}
\end{equation}
This is a rigid code for the shell label, not a derivation of the shell.
\end{proposition}

\begin{proof}
Zeckendorf's theorem gives a unique representation of each positive
integer as a sum of nonconsecutive Fibonacci numbers \cite{Zeckendorf1972,HardyWright2008}.
Applying the greedy algorithm to \(292\) gives
\begin{align*}
292-233&=59,&
59-55&=4,\\
4-3&=1,&
1-1&=0 .
\end{align*}
with no adjacent Fibonacci indices.  Uniqueness is the theorem; physical
selection of the integer is not.
\end{proof}

\begin{theorem}[Diophantine extremality and resonance filtration]
Among irrational rotation numbers, the golden ratio is extremal for
rational approximation in the Hurwitz sense: the constant \(1/\sqrt5\)
is sharp, and the obstruction is realized by the continued fraction
\[
\varphi=[1;1,1,1,\ldots].
\]
Thus a golden-ratio boundary clock is maximally resistant to exact
rational commensuration.  It is a resonance filter.  It is not, by itself,
a capacity selector.
\end{theorem}

\begin{proof}
The classical Hurwitz theorem states that every irrational number admits
infinitely many rational approximants \(p/r\) with
\[
\left|\alpha-\frac pr\right|<\frac{1}{\sqrt5\,r^2},
\]
and the constant \(\sqrt5\) cannot be improved because of the golden-ratio
continued fraction \cite{HardyWright2008}.  Therefore \(\varphi\) is the
most badly approximable rotation number in the precise Diophantine sense.
This prevents exact low-denominator locking more strongly than any other
quadratic-irrational clock, but the theorem concerns commensuration, not
amplitude selection.
\end{proof}

\begin{definition}[Fibonacci-fusion code audit]
For a Fibonacci anyon chain with fusion rule
\[
\tau\times\tau=1+\tau,
\]
the dimension of a fixed-total-charge fusion space grows as a Fibonacci
number, hence as \(\varphi^m\) up to a convention-dependent shift
\cite{NayakEtAl2008}.  Define \(m_q\) by
\begin{equation}
m_q:=\log_\varphi(\sqrt5\,q_{\rm obs}).
\end{equation}
For the present capacity,
\begin{equation}
m_q=587.92446\ldots ,
\end{equation}
so the nearest Fibonacci-chain length is \(m=588\), with
\begin{equation}
\frac{F_{588}}{q_{\rm obs}}=1.03702\ldots .
\end{equation}
\end{definition}

\begin{remark}[Two-sided registry]
The horizon shell and the Fibonacci-fusion chain give two clocks:
\[
N=292,\qquad m\simeq588.
\]
Their relation is
\[
m_q-2N=3.92446\ldots .
\]
The nearby integer \(4\) equals the number of Zeckendorf summands in
\(292=F_{13}+F_{10}+F_4+F_2\).  This is recorded as a registry coincidence,
not as a theorem.  A boundary selector may use it only if an edge-mode
Hamiltonian derives the offset.
\end{remark}

\begin{criterion}[Admissible Fibonacci selector term]
A Fibonacci/Zeckendorf contribution to the boundary selector is admissible
only in the form
\begin{align}
\Gamma_\partial^{\rm Fib}
={}&
\lambda_{\rm sh}
\operatorname{dist}\!\left(
\log_\varphi\sqrt{\frac{q}{\pi}},\mathbb Z
\right)^2 \nonumber\\
&+
\lambda_{\rm code}
\left[
\log q-\log \dim{\mathcal C}_{m(q)}
\right]^2
+\lambda_{\rm Dio}\,{\mathcal R}_{\rm Dio}(q).
\end{align}
where \({\mathcal R}_{\rm Dio}\) penalizes rational locking of the boundary
phase.  This term can constrain the selector.  It cannot replace the
selector:
\[
q_*=\arg\min_q\Gamma_\partial(q)
\]
still requires the full edge-mode Hamiltonian.
\end{criterion}

\begin{remark}[What the golden ratio can and cannot do]
The golden ratio packs recursive data with unique Zeckendorf coding and
frustrates rational standing-wave commensuration.  Quasicrystal spectra
nevertheless remain real spectra, often singular continuous or
multifractal rather than absent \cite{Senechal1995,Jagannathan2021}.
Thus the correct statement is not that \(\varphi\) removes all standing
waves, but that it is the unique extremal clock for suppressing exact
low-denominator resonance.  That is precisely why it belongs in
\(\Gamma_\partial\) as an admissibility term rather than as an independent
derivation of \(q_{\rm obs}\).
\end{remark}






\section{Fibonacci--valuation phase compression of the capacity target}
\label{sec:fib-val-phase-compression}

The previous sections isolate the final problem as the derivation of the logarithmic selector centre \(u_0=\ln q_\star\).  There is, however, a useful compression of this target which uses only two already established discrete structures: the golden-ratio shell clock and the finite valuation-algebra dimension.  This compression is not a microscopic derivation of \(q_\star\).  It is a sharper target for such a derivation.

Let
\begin{equation}
d_{\rm val}:=\dim\mathfrak A_{\mathscr D}=5
\end{equation}
for the \(B^4/S^3\) valuation module, and define the Fibonacci--valuation phase defect
\begin{equation}
\boxed{
\Theta_{\rm FV}(d):=\ln\!\left(1+\frac{\pi}{2d^2}\right).
}
\end{equation}
For an integer shell \(N\), set
\begin{equation}
\boxed{
q_{\rm FV}(N,d)
:=
\pi\varphi^{2N}\exp[-\Theta_{\rm FV}(d)]
=
\frac{\pi\varphi^{2N}}{1+\pi/(2d^2)} .
}
\label{eq:qFV}
\end{equation}
Equivalently,
\begin{equation}
\boxed{
u_{\rm FV}(N,d)
=
\ln q_{\rm FV}
=
\ln\pi+2N\ln\varphi-\ln\!\left(1+\frac{\pi}{2d^2}\right).
}
\end{equation}

For the observed capacity target
\[
q_{\rm obs}=3.31\times10^{122},
\]
the nearest golden shell is
\[
N=292,
\qquad
\log_\varphi\sqrt{\frac{q_{\rm obs}}{\pi}}
=
291.936668\ldots .
\]
The uncorrected shell gives
\begin{equation}
\frac{\pi\varphi^{584}-q_{\rm obs}}{q_{\rm obs}}
=
6.2844029\times10^{-2}.
\end{equation}
With the valuation phase defect at \(d=d_{\rm val}=5\),
\begin{equation}
\boxed{
q_{\rm FV}(292,5)
=
3.3072514607\times10^{122},
}
\end{equation}
and hence
\begin{equation}
\boxed{
\frac{q_{\rm FV}(292,5)-q_{\rm obs}}{q_{\rm obs}}
=
1.14562\times10^{-5}.
}
\label{eq:qFV-residual}
\end{equation}
Thus the capacity amplitude is compressed to a one-integer shell problem:
\begin{equation}
\boxed{
q_{\rm obs}
\simeq
\frac{\pi\varphi^{584}}{1+\pi/50}.
}
\end{equation}

\begin{criterion}[Admissible use of the Fibonacci--valuation compression]
The formula \eqref{eq:qFV} may be used only as a selector target unless a boundary Hamiltonian derives both the integer shell \(N\) and the phase defect \(\Theta_{\rm FV}\).  In particular, the numerical agreement in \eqref{eq:qFV-residual} is not by itself a theorem.  It becomes a theorem only if a microscopic edge-mode calculation gives
\[
N_\star=292,
\qquad
\Theta_\partial=\Theta_{\rm FV}(5),
\]
without using \(q_{\rm obs}\), \(\Lambda_{\rm obs}\), or \(R_\Lambda\) as input.
\end{criterion}

\begin{definition}[Compressed shell selector functional]
The corresponding reduced selector functional is
\begin{equation}
\begin{aligned}
\Gamma_{\partial}^{\rm FV}(q;N,d)
={}&
A\left[
\ln q-\ln q_{\rm FV}(N,d)
\right]^2\\
&+\Gamma_{\rm adm}(q;N,d),
\qquad A>0 .
\end{aligned}
\end{equation}
where \(\Gamma_{\rm adm}\) contains only the already established admissibility constraints: index quantization, Chern--Simons level locking, Bradlow saturation, exchange consistency, and observational nulls.
\end{definition}

\begin{proposition}[What the compression proves and what it cannot prove]
The Fibonacci--valuation compression proves only the following conditional statement.  If the edge-mode Hamiltonian independently selects the shell \(N=292\) and the valuation phase defect \(\Theta_{\rm FV}(5)\), then the resulting capacity agrees with the observed de Sitter capacity at the \(1.2\times10^{-5}\) relative level.  It does not prove that the universe chooses \(N=292\), and it does not compute \(\Lambda\) without a shell-selection dynamics.
\end{proposition}

\begin{proof}
All statements follow by substitution into \eqref{eq:qFV}.  The valuation dimension is \(d_{\rm val}=5\), so
\[
\Theta_{\rm FV}(5)=\ln(1+\pi/50).
\]
For \(N=292\), the formula gives
\[
q_{\rm FV}=\pi\varphi^{584}/(1+\pi/50).
\]
The residual \eqref{eq:qFV-residual} is the direct ratio of this value to the capacity target.  No variational minimum in \(N\) has been derived here; therefore the result is a compressed target, not a microscopic selector.
\end{proof}

This step reduces the remaining problem from a continuous amplitude problem to an integer shell problem plus a boundary phase calculation:
\begin{equation}
\boxed{
\text{derive }N_\star=292
\quad\text{and}\quad
\Theta_\partial=\ln(1+\pi/50)
\text{ from }H_{\partial,q}.
}
\end{equation}


\section{Irreducible selector equation and scale-covariant minimization}

The preceding sweeps leave one honest object:
\[
q_\star=\arg\min_{q\in\mathbb Z_{>0}}\Gamma_\partial(q).
\]
This section records the irreducible form of that problem.  It is a
compression theorem, not a new analogy.

Let
\[
u:=\ln q.
\]
A selector that does not insert an external length scale must be a
dimensionless function of \(u\), topological phases, and running
couplings.  Any smooth large-\(q\) selector with a stable interior
minimum can therefore be written locally as
\begin{equation}
\boxed{
\begin{aligned}
\Gamma_\partial(q)
={}&
\Gamma_0+A(u-u_0)^2\\
&+\Gamma_{\rm osc}(u)+\Gamma_{\rm top}(q),
\qquad A>0 .
\end{aligned}}
\end{equation}
where \(u_0\) is generated by dimensional transmutation or by an
equivalent determinant-line phase condition.  The integer \(q\) is then
selected by nearest admissible index:
\[
\boxed{
q_\star=
\operatorname*{arg\,min}_{q\in\mathbb Z_{>0}}
\Gamma_\partial(q)
=
\operatorname{nearest}_{\rm adm}\bigl(e^{u_0}\bigr).
}
\]
Thus no topological, index, Hopf, Bradlow, Fibonacci, Schwinger, or
Chern--Simons ingredient can by itself choose the observed sector unless
it supplies \(u_0\).

The transmutation version is
\[
\boxed{
u_0
=
\ln\Xi+\frac{16\pi^2}{b\,g_\partial^2}.
}
\]
With \(\Xi=1\) and the observed capacity target,
\[
\boxed{
b\,g_\partial^2
=
\frac{16\pi^2}{\ln q_{\rm obs}}
=
0.55976\pm0.00003 .
}
\]
Equivalently, the selector may be written as a phase condition
\[
\boxed{
\Theta_{\partial}(q_\star)
=
\Theta_{\rm bd},
}
\]
where \(\Theta_{\partial}\) is a determinant-line or edge-mode phase.
For the Fibonacci registry the measured phase defect is
\[
\Theta_\varphi(q_{\rm obs})
=
\ln\!\left(\frac{\pi\varphi^{584}}{q_{\rm obs}}\right)
=
0.060948\ldots ,
\]
close to
\[
\ln\!\left(1+\frac{\pi}{50}\right)=0.0609369\ldots .
\]
The difference is
\[
1.15\times10^{-5}.
\]
This is retained only as a target for a boundary phase calculation.  It
is not a theorem until a concrete \(H_{\partial,q}\) or \(D_{\partial,q}\)
derives it without using \(\Lambda_{\rm obs}\).

\begin{criterion}[Selector closure test]
A proposed microscopic completion closes the capacity problem if and
only if it supplies one of the following without using \(q_{\rm obs}\),
\(\Lambda_{\rm obs}\), or \(R_\Lambda\) as an input:
\[
\boxed{
\begin{gathered}
u_0=\ln q_{\rm obs},\\
b\,g_\partial^2=0.55976\pm0.00003,\\
\Theta_{\rm bd}=0.060948\ldots,\\
\log\dim\mathcal C_q=q_{\rm obs}
\quad\text{with }q_\star\text{ dynamically selected.}
\end{gathered}}
\]
If it only reproduces \(q=A/(4\ell_P^2)\) after the area is supplied,
it is a capacity representation, not a selector.
\end{criterion}

\begin{remark}[What is now excluded]
The admissible structures already found remain essential:
index quantization, Chern--Simons level lock, Hopf helicity, Bradlow
packing, Schwinger admissibility, APS determinant data, valuation
closure, and Fibonacci nonresonance.  But each of them is a constraint
or coordinate.  The missing dynamical datum is the logarithmic center
\(u_0\).  The final problem is therefore not broad:
\[
\boxed{
\text{derive }u_0\text{ from the boundary edge-mode Hamiltonian.}
}
\]
\end{remark}


\section{Selector irreducibility theorem and arithmetic exhaustion audit}
\label{sec:selector-irreducibility}

This section records the point at which the search ceases to be a catalogue of
admissibility invariants and becomes a sharp obstruction theorem.  Put
\begin{equation}
u:=\ln q .
\end{equation}
The observed capacity target corresponds to
\begin{equation}
u_{\rm obs}=\ln q_{\rm obs}=282.111\pm0.015 .
\end{equation}
Thus the selector cannot be a generic algebraic or polynomial condition in
\(q\).  It must manufacture a logarithmic centre of order \(10^2\), while all
dimensionless microscopic inputs are expected to be of order unity.

\begin{definition}[Scale-algebraic selector class]
A scale-algebraic boundary selector is a function of the form
\begin{equation}
\Gamma(u)=P(u)+\sum_{r=1}^{m}c_r e^{\alpha_r u}
          +B(u),
\end{equation}
where \(P\) is a polynomial with order-unity coefficients, the nonzero
\(\alpha_r\) and \(c_r\) are order-unity constants, and \(B\) is bounded
together with its first two derivatives.  Oscillatory clock terms, Fibonacci
registry penalties, topological phases, and determinant phases belong to
\(B\) unless they introduce a new large logarithmic centre.
\end{definition}

\begin{theorem}[No large selector without transmutation or hidden input]
\label{thm:no-large-selector}
A stable minimum of a scale-algebraic selector at \(u_\star\gg1\) cannot be
forced by index, topology, Chern--Simons level locking, Hopf charge, Bradlow
packing, Schwinger admissibility, APS phase data, or Fibonacci nonresonance
alone.  Such invariants can restrict the admissible lattice of sectors, but
they do not by themselves create \(u_\star=\ln q_{\rm obs}\).  A large
minimum requires at least one of the following:
\begin{enumerate}
\item an explicit logarithmic centre \(u_0\) in the edge Hamiltonian;
\item dimensional transmutation,
\[
u_0=\ln\Xi+\frac{16\pi^2}{b g_\partial^2};
\]
\item a near-cancellation between leading exponential terms, which is itself
a hidden scale or tuning datum;
\item an external observational insertion of \(q_{\rm obs}\), which is not a
microscopic derivation.
\end{enumerate}
\end{theorem}

\begin{proof}
For \(u\gg1\), the derivative of the largest exponential term
\(c e^{\alpha u}\) dominates unless at least two leading exponentials cancel
to exponential accuracy.  Such a cancellation is an additional scale datum.
If no positive exponential is present, then \(P'(u)+B'(u)=0\) can place a
minimum at \(u\gg1\) only by an order-\(u\) coefficient or an explicit shift
in \(P\); that is again a hidden large input.  Bounded phases and topological
integer constraints can move a minimum by an order-unity amount or choose a
nearest admissible integer, but cannot generate \(u_\star\simeq282\).  The
renormalization-group solution \(u_0=\ln\Xi+16\pi^2/(b g_\partial^2)\) is the
standard way to create a large logarithm from order-unity data.  This proves
the stated alternatives.
\end{proof}

For the present target and \(\Xi=1\), the transmutation number is
\begin{equation}
\boxed{
b g_\partial^2=
\frac{16\pi^2}{u_{\rm obs}}
=0.55976\pm0.00003 .
}
\end{equation}
An arithmetic sweep over the natural constants appearing in the atlas gives
near coincidences, but no invariant derivation.  For example,
\begin{align}
e^{-\gamma_{\rm E}}&=0.56145948\ldots,\\
\frac{\sqrt5}{4}&=0.55901699\ldots,\\
\ln\varphi+\frac{\pi}{40}&=0.55975164\ldots .
\end{align}
The last number is close, but it is not a theorem; it is a warning.  Once
\(\varphi,\pi\), and a small denominator are allowed freely, many close
relations can be manufactured.  Therefore the manuscript admits such
relations only as selector targets.  They become physics only after a concrete
boundary operator or edge Hamiltonian derives them.

\begin{criterion}[Final selector standard]
A proposed completion of the capacity problem is accepted only if it produces
one of the following without using \(\Lambda_{\rm obs}\), \(R_\Lambda\), or
\(q_{\rm obs}\) as input:
\begin{equation}
\boxed{
u_0=282.111\pm0.015 ,
\qquad
b g_\partial^2=0.55976\pm0.00003 ,
}
\end{equation}
or an exactly equivalent determinant-line phase condition whose nearest
admissible sector is \(q_\star=3.31\times10^{122}\).  All other invariants
remain necessary admissibility data, not the selector itself.
\end{criterion}

This is the irreducible form of the remaining problem.  The full atlas has
not failed; it has reduced the target to a single edge-dynamical number.

\section{No-go statements}

\begin{nogotheorem}[No local vacuum-density solution]\label{nogo:vacuum-density}
No mechanism that leaves local vacuum energy as a freely gravitating source solves the cosmological-constant problem within the present invariant criterion.
\end{nogotheorem}

\begin{proof}
If local vacuum energy gravitates as \(T^{\rm vac}_{\mu\nu}=-\rho_{\rm vac}g_{\mu\nu}\), then arbitrary shifts of the zero of energy shift the curvature unless cancelled by a counterterm.  Such cancellation is a parametrization unless a principle fixes the invariant combination; this is the same structural target pursued by vacuum-energy sequestering mechanisms \cite{KaloperPadilla2014}.  In the present mechanism the trace-free projector removes the shift and the boundary capacity fixes the remaining integration mode.
\end{proof}

\begin{nogotheorem}[No capacity from \(S^3\) topology alone]\label{nogo:no-topology-capacity}
The topology \(\Sigma\simeq S^3\) and the finite valuation spectrum do not by themselves compute \(\Ncap\).
\end{nogotheorem}

\begin{proof}
The valuation spectrum supplies powers of the scale factor and algebraic projectors.  It contains no Planck-scale Hilbert-space dimension.  By \cref{nogo:topology-amplitude}, topology and dimensionless homogeneous algebra cannot alone select a nonzero dimensional \(\Lambda\).  The additional microscopic input is precisely \(\mathcal H_\partial\) and its capacity.
\end{proof}

\begin{nogotheorem}[No single four-dimensional second-order metric correction]\label{nogo:lovelock}
Under the usual four-dimensional assumptions of locality, metric dependence only, second-order field equations, and covariant conservation, a new arbitrary left-hand side cannot be the solution.  The Lovelock uniqueness result leaves only the Einstein tensor plus a term proportional to \(g_{\mu\nu}\) \cite{Lovelock1971,Curiel2016,PadmanabhanKothawala2013}.
\end{nogotheorem}

\section{Empirical interpretation}

The capacity-fixed mechanism changes what must be measured or computed.

First, cosmology measures an effective \(\Lambda_{\rm obs}\).  The theory demands
\begin{equation}
\boxed{
\Ncap
=
\frac{3\pi}{\Lambda_{\rm obs}\ell_P^2}
}
\end{equation}
as the target capacity.  This number is large, of order \(10^{122}\) in the observed late-time cosmology \cite{Planck2018}.  The microscopic task is to derive this capacity from \(\mathcal H_\partial\), not to tune a vacuum density.


\begin{remark}[Tensor propagation is not the \(s=-2\) ledger branch]
The homogeneous \(s=-2\) coefficient in \(F(a)\) is a background expansion term.  It is geometrical curvature before the split and a \(w=-1/3\) physical stress branch after the split.  It is not the graviton mass and not the tensor-wave speed.  The condition \(c_T=c\) constrains the perturbative spin--2 lift of the gravitational action, not the presence or absence of a homogeneous \(a^{-2}\) power in the background ledger.
\end{remark}


Second, multimessenger gravity requires a luminal tensor sector,
\begin{equation}
\boxed{c_T=c,}
\end{equation}
consistent with the observed near-simultaneous gravitational-wave and electromagnetic propagation in GW170817/GRB 170817A \cite{AbbottGW170817,MonitorGRB170817A}.

Third, the dark sector must pass both the cold-charge test and the valuation-spectrum test.  Cold dark matter is protected by \(Q_{\rm DM}\), while the homogeneous nonstandard valuation branch is tested by \(\mathcal C_{-1}[F]\) and free-spectrum Hankel reconstruction.  Lensing systems such as the Bullet Cluster motivate a stable gravitating dark sector independent of baryonic plasma \cite{Clowe2006,Markevitch2004}.

Fourth, the real-data audit in \cref{sec:real-data-audit} fixes the empirical status of the present manuscript.  The invariant framework is a falsifiable architecture that survives BAO-only tests, while superiority over flexible dark-energy parameterizations is not established by the positive \(a^{-1}\) branch alone.



\section{Irreducible Boundary Selector: Renormalization Test and Failure Modes}

The preceding atlas shows that topology, index theory, Chern--Simons level,
Hopf helicity, Bradlow packing, Schwinger admissibility, APS determinant
data, valuation closure, and Fibonacci nonresonance all constrain the
admissible sectors.  None of them by itself supplies the logarithmic centre
\[
        u_0=\ln q_* .
\]
This section records the remaining irreducible calculation.  It is the
part of the problem that cannot be replaced by another invariant label.

Let
\begin{equation}
        u:=\ln q,
        \qquad
        q\in\mathbb Z_{>0}.
\end{equation}
A selector that is not merely an observational restatement must determine
\(u_0\) before \(\Lambda_{\rm obs}\), \(R_\Lambda\), or \(q_{\rm obs}\) are
inserted.  The transmutation form is
\begin{equation}
        u_0=\ln\Xi+\frac{16\pi^2}{b\,g_\partial^2}.
\end{equation}
With \(\Xi=1\), the current capacity target requires
\begin{equation}
\boxed{
        b\,g_\partial^2
        =\frac{16\pi^2}{\ln q_{\rm obs}}
        =0.55976\pm0.00003 .
}
\label{eq:preobs-target-coupling}
\end{equation}

\begin{criterion}[Pre-observational selector]
A boundary edge-mode Hamiltonian solves the capacity-selection problem only
if it gives one of the following without using \(q_{\rm obs}\) as input:
\begin{equation}
\boxed{
\begin{gathered}
        u_0=282.111\pm0.015,\\
        b\,g_\partial^2=0.55976\pm0.00003,\\
        \Theta_{\rm bd}=\Theta_\varphi(q_*).
\end{gathered}}
\end{equation}
A formula that first reads \(q_{\rm obs}\) from cosmology and then rewrites
it in these variables is only a dictionary.
\end{criterion}

\begin{theorem}[Renormalization necessity]
Assume that the boundary selector has no inserted large number and depends
on \(q\) only through a finite combination of algebraic powers, logarithms,
bounded phases, index labels, and saturation penalties with \(O(1)\)
coefficients.  Then a stable minimum at \(u=\ln q\gg1\) requires either a
dimensional-transmutation term, a protected cancellation producing such a
term, or an explicit large external datum.
\end{theorem}

\begin{proof}
Algebraic powers in \(q\) become exponentials in \(u\); index and saturation
terms constrain admissibility but do not produce a finite large centre by
themselves.  Bounded phases are periodic or quasiperiodic in \(u\), hence
cannot supply a unique large nonperiodic centre.  A polynomial or finite
logarithmic expression with \(O(1)\) coefficients has stationary points at
\(O(1)\) unless a large coefficient is inserted.  Therefore a stable
minimum at \(u\simeq282\) must arise from an exponential hierarchy such as
\(\exp[16\pi^2/(b g_\partial^2)]\), from an equivalent cancellation, or
from observational insertion.  The last option is not a derivation.
\end{proof}

The numerical situation is sharp.  Several natural constants sit close to
the required target,
\begin{align}
        e^{-\gamma_{\rm E}}&=0.5614594836\ldots,\\
        \sqrt5/4&=0.5590169944\ldots,\\
        \ln\varphi+\pi/40&=0.5597516369\ldots ,
\end{align}
but none is accepted as a law without an edge Hamiltonian that derives it.
The last expression is especially close, but it is too easily manufactured
from the existing registry constants; it is therefore a target, not a
proof.

\begin{center}
\includegraphics[width=\columnwidth]{fig_renormalization_selector.png}

{\small\textbf{Figure.} Dimensionless transmutation target compared with nearby constants.
The plot is a diagnostic, not a derivation.  The selector problem is to
derive the target from boundary edge dynamics rather than to fit it after
the fact.}
\end{center}

\begin{corollary}[Final reduction of the capacity problem]
After imposing the invariant atlas, the remaining calculation is unique in
form:
\begin{equation}
\boxed{
\begin{gathered}
        \text{construct }H_{\partial,q},\\
        \Gamma_\partial(q)
        =-\log{\rm Tr}_{\mathcal C_q}
        e^{-\beta_qH_{\partial,q}},\\
        \Gamma_\partial\Longrightarrow q_* .
\end{gathered}}
\end{equation}
The calculation succeeds only if \(q_*\) is obtained without placing
\(q_{\rm obs}\) into \(H_{\partial,q}\), \(\beta_q\), or the couplings.
\end{corollary}



\section{Terminal edge-selector coupling and two-coordinate closure test}
\label{sec:terminal-edge-selector-coupling}

The preceding no-go statements leave a single possible large-number mechanism:
a boundary edge-mode Hamiltonian must generate the logarithmic center
\[
u_0=\ln q_\star
\]
by dimensional transmutation, or by an exactly equivalent determinant-line
phase condition.  This section records the resulting terminal numerical
target.  It is not an additional analogy.

Let
\[
q_{\rm obs}=3.31\times10^{122},
\qquad
u_{\rm obs}:=\ln q_{\rm obs}.
\]
Then
\[
\boxed{
u_{\rm obs}=282.111\pm0.015
}
\]
and the transmutation form
\[
q_\star
=
\Xi\exp\!\left(\frac{16\pi^2}{b\,g_\partial^2}\right)
\]
gives, for \(\Xi=1\),
\begin{equation}
\boxed{
b\,g_\partial^2
=
\frac{16\pi^2}{u_{\rm obs}}
=
0.55976\pm0.00003 .
}
\label{eq:terminal-bg-target}
\end{equation}
Thus the last unsolved boundary task is no longer qualitative: a concrete
edge-mode Hamiltonian must derive \cref{eq:terminal-bg-target}, or must
derive an equivalent phase/determinant condition.

There is a second coordinate system for the same target.  Let
\[
\varphi=\frac{1+\sqrt5}{2},
\qquad
d_{\rm val}=\dim\mathfrak A_{\mathscr D}=5.
\]
For the nearest golden shell \(N=292\), define the
FV-compressed capacity
\begin{equation}
\boxed{
q_{\rm FV}(N,d)
=
\frac{\pi\varphi^{2N}}{1+\pi/(2d^2)} .
}
\label{eq:qfv-def}
\end{equation}
Then
\[
q_{\rm FV}(292,5)
=
3.3072514607\times10^{122},
\]
and hence
\begin{equation}
\boxed{
\frac{q_{\rm FV}(292,5)-q_{\rm obs}}{q_{\rm obs}}
=
1.1456\times10^{-5}.
}
\end{equation}
Equivalently,
\[
\frac{16\pi^2}{\ln q_{\rm FV}(292,5)}
=
0.55976 ,
\]
which differs from \cref{eq:terminal-bg-target} by
\[
\boxed{
-2.27\times10^{-8}.
}
\]

\begin{center}
\includegraphics[width=.96\columnwidth]{fig_boundary_selector_terminal.png}

{\small\textbf{Figure.} Terminal boundary-selector coupling audit.  The target value
\(b g_\partial^2=16\pi^2/\ln q_{\rm obs}\) is compared with the
Fibonacci--valuation compressed value and with simple nearby constants.
Only a boundary Hamiltonian derivation can promote any entry to a law.}
\end{center}

\begin{criterion}[Terminal selector criterion]
A proposed completion of the capacity mechanism is accepted as a
selector theorem only if it derives, without using \(\Lambda_{\rm obs}\),
\(R_\Lambda\), \(q_{\rm obs}\), or \(u_{\rm obs}\) as input, at least one of
the following equivalent data:
\[
\boxed{
b\,g_\partial^2=0.55976\pm0.00003 ,
}
\]
or
\[
\boxed{
N_\star=292,
\qquad
\Theta_\partial=\ln\!\left(1+\frac{\pi}{50}\right)
}
\]
together with the Fibonacci--valuation map \cref{eq:qfv-def}, up to a
controlled and independently derived correction of relative order
\(10^{-5}\).
\end{criterion}

\begin{theorem}[Two-coordinate terminal equivalence]
Assume a boundary selector produces both a running coupling and a
golden-shell phase coordinate, with
\[
q_\star
=
\exp\!\left(\frac{16\pi^2}{b\,g_\partial^2}\right)
=
\frac{\pi\varphi^{2N_\star}}{1+\pi/(2d_{\rm val}^2)}
\,e^{-\Delta_{\rm bd}} .
\]
Then the two descriptions are equivalent precisely when
\begin{equation}
\boxed{
\frac{16\pi^2}{b\,g_\partial^2}
=
\ln\pi+2N_\star\ln\varphi
-\ln\!\left(1+\frac{\pi}{2d_{\rm val}^2}\right)
-\Delta_{\rm bd}.
}
\end{equation}
For \(d_{\rm val}=5\) and \(N_\star=292\), the observed target corresponds
to
\[
\boxed{
\Delta_{\rm bd}=1.1456\times10^{-5}+O(10^{-10})
}
\]
in logarithmic units.  Thus the final problem may be stated either as a
boundary beta-function problem or as a determinant-line phase problem;
the invariant content is the same logarithmic center \(u_0\).
\end{theorem}

\begin{proof}
Take the logarithm of the equality of the two sector counts.  The
displayed identity follows immediately.  Substitution of \(d_{\rm val}=5\),
\(N_\star=292\), and \(q_{\rm obs}\) gives
\[
\Delta_{\rm bd}
=
\ln q_{\rm FV}(292,5)-\ln q_{\rm obs}
=
1.1456\times10^{-5}+O(10^{-10}).
\]
No physical conclusion follows from the numerical smallness alone; the
point is that any real completion must derive this last residual, not fit
it.
\end{proof}

\begin{remark}[Where the search is now irreducible]
All previously found structures have been sorted.  Index, topology,
Chern--Simons level, Hopf helicity, Bradlow saturation, Schwinger
admissibility, APS determinant data, valuation algebra, and Fibonacci
coding are necessary filters or coordinates.  None of them alone supplies
the logarithmic center \(u_0\).  The remaining scientific task is therefore
exactly this:
\[
\boxed{
\text{derive }u_0
\text{ from the boundary edge-mode Hamiltonian.}
}
\]
\end{remark}


\section{Boundary beta-functional selector: explicit running-coupling form and terminal obstruction}

The previous reductions leave one irreducible numerical datum,
\begin{equation}
u_\star:=\ln q_\star,
\end{equation}
which cannot be manufactured by topology, index theory, Chern--Simons level
locking, Hopf helicity, Bradlow saturation, Schwinger admissibility, APS
determinant phases, valuation closure, or the Fibonacci registry alone.
This section records the most economical edge-dynamical mechanism that can
generate such a logarithmic center without inserting a dimensional scale.

\begin{definition}[Boundary running-coupling selector]
Let \(u=\ln q\) be the logarithmic capacity coordinate.  A one-coupling
edge selector is a boundary sector with a dimensionless coupling
\(g_\partial(u)\) satisfying
\begin{equation}
\frac{\dd}{\dd u}\,g_\partial^{-2}(u)
=
-\frac{b}{16\pi^2},
\qquad b>0,
\end{equation}
and an effective edge potential
\begin{equation}
\boxed{
\Gamma_{\partial}^{\rm run}(u)
=
\Gamma_0+
\frac{\kappa}{2}
\left[g_\partial^{-2}(u)\right]^2
+
\Gamma_{\rm adm}(u),
\qquad
\kappa>0 .
}
\end{equation}
Here \(\Gamma_{\rm adm}\) denotes the already established admissibility
terms: index integrality, Chern--Simons level locking, Bradlow saturation,
Schwinger admissibility, determinant-line phase, valuation support, and
Fibonacci nonresonance.
\end{definition}

\begin{theorem}[Running-coupling selector]
If \(\Gamma_{\rm adm}\) is locally flat compared with the quadratic
running term at its minimum, then the leading capacity sector is
\begin{equation}
\boxed{
u_\star
=
\frac{16\pi^2}{b\,g_\partial^2(0)} ,
\qquad
q_\star=\exp\!\left[\frac{16\pi^2}{b\,g_\partial^2(0)}\right].
}
\end{equation}
Equivalently, for a prefactor \(\Xi\),
\begin{equation}
\boxed{
q_\star=\Xi\exp\!\left[\frac{16\pi^2}{b\,g_\partial^2(0)}\right].
}
\end{equation}
\end{theorem}

\begin{proof}
The beta equation integrates to
\begin{equation}
g_\partial^{-2}(u)
=
g_\partial^{-2}(0)-\frac{b}{16\pi^2}u .
\end{equation}
The quadratic term in \(\Gamma_{\partial}^{\rm run}\) is minimized when
\(g_\partial^{-2}(u)=0\), hence
\[
u_\star=16\pi^2/[b\,g_\partial^2(0)] .
\]
A \(q\)-independent normalization of the admissible edge code multiplies
\(q_\star\) by the dimensionless factor \(\Xi\).
\end{proof}

\begin{corollary}[Terminal numerical target]
For \(\Xi=1\), the observed capacity target is equivalent to
\begin{equation}
\boxed{
b\,g_\partial^2(0)
=
\frac{16\pi^2}{\ln q_{\rm obs}}
=
0.55976\pm0.00003 .
}
\end{equation}
Thus the cosmological-constant amplitude problem has been reduced to a
dimensionless boundary beta-functional problem: compute the product
\(b\,g_\partial^2(0)\) from the edge Hamiltonian without using
\(\Lambda_{\rm obs}\), \(R_\Lambda\), \(q_{\rm obs}\), or \(u_{\rm obs}\) as input.
\end{corollary}

\begin{criterion}[True completion versus restatement]
A proposed microscopic completion of the capacity sieve is a genuine
selector only if it supplies:
\begin{enumerate}
\item a concrete boundary Hilbert space and Hamiltonian \(H_{\partial,q}\);
\item a derivation of the running law for \(g_\partial(u)\);
\item a calculation of \(b\,g_\partial^2(0)\);
\item the admissibility projection onto the index/Chern--Simons/Bradlow/
Schwinger/valuation/Fibonacci-compatible sectors;
\item the minimum \(q_\star\) without inserting any late-time cosmological
observable.
\end{enumerate}
Otherwise it is not a derivation of \(\Lambda\); it is a reparametrization of
the already identified target.
\end{criterion}

\begin{remark}[What the beta-functional step achieves]
This section does not claim that the required edge Hamiltonian has already
been derived.  It proves that once ordinary topological and spectral
quantizers have exhausted their role, the remaining large number must enter
through a logarithmic center.  A running boundary coupling is the minimal
known invariant mechanism that produces such a center from \(O(1)\)
dimensionless data.  The next calculation is therefore not another search
for an analogy, but the construction and renormalization of
\(H_{\partial,q}\).
\end{remark}




\section{Terminal Edge-Mode Selector Audit}\label{sec:terminal-edge-selector-audit}

The previous sections reduce the amplitude problem to a single logarithmic
selector.  This section records the final obstruction in a form that can be
checked directly from an edge Hamiltonian.  No new analogue is introduced here;
only the admissible Hamiltonian mechanisms are separated from the mechanisms
that merely quantize or bound the sector.

Let
\begin{equation}
u:=\ln q,
\qquad
\Gamma_\partial(q):=
-\ln\operatorname{Tr}_{\mathcal C_q}
\exp[-\beta_qH_{\partial,q}].
\end{equation}
A genuine microscopic solution of the capacity problem is an autonomous
calculation of \(H_{\partial,q}\) such that
\begin{equation}
q_*=\arg\min_{q\in\mathbb Z_{>0}}\Gamma_\partial(q)
\end{equation}
and \(q_*\) agrees with the observed capacity without inserting
\(\Lambda_{\rm obs}\), \(R_\Lambda\), \(q_{\rm obs}\), or \(u_{\rm obs}\).

\begin{theorem}[Topological edge sectors do not generate the large logarithm]\label{thm:topological-edge-no-large-u}
Assume that the boundary edge Hamiltonian is a finite sum of the following
autonomous topological or saturated sectors:
\begin{align}
H_{\partial,q}
={}&H_{\rm CS}(k=2q)+H_{\rm ind}(q)+H_{\rm Br}(q)\nonumber\\
&+H_{\rm Hopf}(Q_H)+H_{\rm Sch}+H_{\rm val}.
\end{align}
Assume further that all dimensionless coefficients are \(O(1)\) and that no
running coupling, protected cancellation, or explicit large datum is present.
Then the resulting effective action cannot have a stable isolated minimum at
\[
u=\ln q\gg1
\]
unless one of the included sectors is already supplied with a large scale or a
large quantum number.  In particular, Chern--Simons level quantization,
Dolbeault index quantization, Bradlow saturation, Hopf helicity, Schwinger
admissibility, and the finite valuation algebra are necessary admissibility
data but are not by themselves a selector for \(q_{\rm obs}\sim10^{122}\).
\end{theorem}

\begin{proof}
Each listed sector has one of three forms.  It either fixes an integer label
(\(k=2q\), \(\operatorname{ind}\bar\partial=q\), \(N_{\rm v}=q\)), imposes an
inequality or wall (Bradlow and Schwinger admissibility), or contributes a
bounded topological phase or finite-dimensional algebraic checksum.  Therefore
its contribution to \(\Gamma_\partial\) is a combination of affine, logarithmic,
bounded oscillatory, or powerlike functions of \(q\) with coefficients fixed
independently of \(q_{\rm obs}\).  In the logarithmic variable \(u\), such terms
are combinations of constants, \(u\), bounded functions of \(u\), and
\(\exp(\alpha u)\) with \(O(1)\) coefficients.  A stable stationary point at
\(u\gg1\) then requires either a cancellation between large terms to exponential
precision, an externally inserted center \(u_0\), or a running coupling whose
zero dynamically defines \(u_0\).  The first two alternatives are not autonomous
derivations of the capacity amplitude.  Hence an autonomous large selector must
come from the running-coupling or equivalent dimensional-transmutation channel.
\end{proof}

\begin{corollary}[Irreducible boundary beta target]\label{cor:irreducible-beta-target}
If the selector is autonomous and contains no hidden large datum, the surviving
minimal form is
\begin{equation}
\Gamma_{\partial}^{\rm run}(u)
=
\Gamma_0+
\frac{\kappa}{2}
\left[g_\partial^{-2}(u)\right]^2
+\Gamma_{\rm adm}(u),
\qquad
\kappa>0,
\end{equation}
with
\begin{equation}
\frac{\dd}{\dd u}g_\partial^{-2}(u)
=
-\frac{b}{16\pi^2}.
\end{equation}
Neglecting bounded admissibility ripples, the selected logarithmic center is
\begin{equation}
u_*=\frac{16\pi^2}{b\,g_\partial^2(0)},
\qquad
q_*=\exp u_* .
\end{equation}
Thus the present capacity target is equivalently the dimensionless boundary
renormalization target
\begin{equation}
\boxed{
b\,g_\partial^2(0)
=
\frac{16\pi^2}{\ln q_{\rm obs}}
=
0.55976\pm0.00003 .
}
\end{equation}
\end{corollary}

\begin{criterion}[What remains to be computed]\label{crit:terminal-computation}
A terminal solution is not another invariant checksum.  It is a computation of
\(b\) and \(g_\partial(0)\), or of an exactly equivalent determinant-line
center, from a specified boundary edge-mode Hamiltonian.  The calculation must
not use \(\Lambda_{\rm obs}\), \(R_\Lambda\), \(q_{\rm obs}\), or
\(u_{\rm obs}\) as an input.  Until this is done, the manuscript has reduced the
cosmological-constant amplitude problem to an irreducible boundary
renormalization problem, but has not solved the amplitude as a theorem.
\end{criterion}

\begin{remark}[Terminal ledger]
For the Planck/CMB capacity used throughout the numerical ledger,
\begin{align}
u_{\rm obs}&=282.111\pm0.015,\\
b\,g_\partial^2(0)&=0.55976\pm0.00003 .
\end{align}
The near coincidences with \(e^{-\gamma_E}\), \(\sqrt5/4\), and
\(\ln\varphi+\pi/40\) are retained only as diagnostic targets.  Without an
edge Hamiltonian deriving one of them, they are not laws.
\end{remark}




\section{Terminal selector reduction and arithmetic degeneracy audit}

The preceding atlas leaves a single mathematical object uncomputed: the
boundary edge-mode Hamiltonian that fixes the logarithmic capacity center.
It is useful to state this as a terminal reduction rather than to enlarge
the catalogue of analogies.

\begin{definition}[Boundary logarithmic capacity coordinate]
Let
\begin{equation}
        u:=\ln q,
        \qquad q\in\mathbb Z_{>0}.
\end{equation}
A boundary selector is a family of sector effective actions
\begin{equation}
        \Gamma_\partial(q)
        =
        -\ln {\rm Tr}_{\mathcal C_q}
        \exp[-\beta_q H_{\partial,q}],
\end{equation}
together with the discrete extremum condition
\begin{equation}
        \Delta_q\Gamma_\partial(q_*)=0,
        \qquad
        \Delta_q^2\Gamma_\partial(q_*)>0 .
\end{equation}
\end{definition}

\begin{theorem}[Terminal selector reduction]\label{thm:terminal-selector-reduction}
Assume that all admissible-sector conditions already imposed in this
paper hold: trace-free vacuum-shift decoupling, Bianchi integration,
horizon capacity quantization, Chern--Simons level locking, Hopf/helicity
admissibility, Bradlow saturation, Schwinger admissibility, APS
determinant-line bookkeeping, valuation closure, and the Fibonacci
nonresonant registry.  Then these conditions determine the allowed
sector lattice and its consistency tests, but they do not determine the
large logarithmic center \(u_0=\ln q_*\).

A derivation of the observed sector must therefore supply one additional
piece of dynamical information, expressible in one of the equivalent
forms
\begin{equation}
\boxed{
        u_0,\qquad
        b\,g_\partial^2(0),\qquad
        \Theta_{\rm bd},\qquad
        \dim\mathcal C_{q_*}.
}
\end{equation}
For the running-coupling representation this datum is
\begin{equation}
\boxed{
        u_0=\frac{16\pi^2}{b\,g_\partial^2(0)}
        \quad\Longleftrightarrow\quad
        b\,g_\partial^2(0)
        =
        \frac{16\pi^2}{\ln q_*}.
}
\end{equation}
Thus the capacity problem is solved as a theorem only when
\(b\) and \(g_\partial(0)\), or an exactly equivalent determinant/phase
datum, are derived from \(H_{\partial,q}\) without using
\(\Lambda_{\rm obs}\), \(R_\Lambda\), \(q_{\rm obs}\), or \(u_{\rm obs}\)
as input.
\end{theorem}

\begin{proof}
The listed structures are invariant under changes of units, pivots, and
smooth deformations inside a fixed admissible sector.  Consequently they
can quantize \(q\), restrict signs and supports, impose saturation
bounds, detect anomalies, or reject inadmissible flows.  None of these
operations supplies a scale center of size
\(u\simeq 282\) unless a dimensionless coupling, determinant phase, or
explicit large datum enters the sector action.  If the sector action is
composed only of bounded phases, powers, logarithms, topological
integers, and \(O(1)\) coefficients, then its stationary equations either
select \(O(1)\) values, boundary points of the admissible domain, or
unbounded directions.  A stable interior minimum at exponentially large
\(q\) requires a transmutation-type relation, a protected cancellation
that is itself a dynamical datum, or insertion of the large scale.  The
last option is not a derivation.  Hence the selector must reduce to the
displayed datum.
\end{proof}

\begin{criterion}[Arithmetic proximity is not a selector]\label{crit:arith-proximity}
A numerical expression involving familiar constants, for example
\(\pi\), \(\varphi\), \(\gamma_E\), or low integers, may be recorded as a
diagnostic target only if it is not used as an input.  It becomes a
physical selector only when it is obtained from the spectrum, phase, or
renormalization of a specified \(H_{\partial,q}\).  Otherwise it is an
arithmetic coincidence, because incommensurate constants admit many
post-hoc approximants to any prescribed real number.
\end{criterion}

\begin{proposition}[Computational audit target]\label{prop:terminal-coupling-target}
Using the present capacity normalization,
\begin{align}
        q_{\rm obs}&=3.31\times 10^{122},\\
        u_{\rm obs}&=\ln q_{\rm obs}=282.111\pm0.015 .
\end{align}
The terminal running-coupling target is
\begin{equation}
\boxed{
        b\,g_\partial^2(0)
        =
        0.55976\pm0.00003 .
}
\end{equation}
The close arithmetic comparisons
\begin{equation}
        e^{-\gamma_E},\qquad
        \frac{\sqrt5}{4},\qquad
        \ln\varphi+\frac{\pi}{40},
        \qquad
        \frac{\pi}{2(\dim\mathfrak A_{\mathscr D})^2}
\end{equation}
are therefore not conclusions.  They are tests for a future edge-mode
calculation.  A successful calculation must output the target, or an
equivalent \(u_0\), from the boundary operator algebra itself.
\end{proposition}

\begin{remark}[No further enlargement without a Hamiltonian]
The atlas has reached a terminal form.  Adding further universal
invariants is useful only if they change \(H_{\partial,q}\) or its
renormalization.  If a proposed new structure merely supplies another
integer, bound, phase, or analogy, it belongs in the admissibility
catalogue, not in the selector theorem.
\end{remark}





\section{Scheme-covariant terminal selector correction}\label{sec:scheme-covariant-terminal-correction}

This section supersedes the raw terminal coupling target in the preceding
ledger.  The correction is structural, not cosmetic.  The quantity
\(q_\star\), or equivalently \(u_\star=\ln q_\star\), is physical inside the
declared capacity prescription.  A separated product \(b\,g_\partial^2(0)\)
is not physical until the renormalization convention, threshold prefactor,
loop order, and boundary theory dimension have been fixed.

\begin{definition}[Scheme-covariant selector data]
A running edge selector has the general transmutation form
\begin{equation}
\boxed{
        u_\star
        =
        \ln q_\star
        =
        \ln\Xi+
        \frac{16\pi^2}{b\,g_\partial^2(0)}
        +\Delta_{\rm th}
        +\Delta_{\rm 2loop}
        +\cdots .
}
\end{equation}
Here \(\Xi\) is the dimensionless threshold or determinant prefactor, and
\(\Delta_{\rm th},\Delta_{\rm 2loop},\ldots\) denote scheme-dependent but
computable corrections.  Only the complete sum is a selector.  Thus a claim to
derive \(b\,g_\partial^2(0)\) alone is meaningful only after \(\Xi\) and the
renormalization scheme have also been derived in the same convention.
\end{definition}

\begin{proposition}[Observational precision of the terminal target]\label{prop:observational-precision}
Using the Planck-only normalization adopted in the numerical ledger,
the observationally meaningful terminal numbers are
\begin{equation}
\boxed{
        u_{\rm obs}=\ln q_{\rm obs}=282.111\pm0.015 ,
}
\end{equation}
and, in the one-loop four-dimensional convention with \(\Xi=1\),
\begin{equation}
\boxed{
        b\,g_\partial^2(0)
        =
        \frac{16\pi^2}{u_{\rm obs}}
        =
        0.55976\pm0.00003 .
}
\end{equation}
The digits beyond this uncertainty are arithmetic bookkeeping, not measured
physics.  If the local value of \(H_0\) is used instead of the CMB anchor, the
shift is larger and only the first few digits remain meaningful.

Consequently the near equalities
\[
\ln\varphi+\frac{\pi}{40}\simeq 0.55976,
\qquad
q_{\rm FV}(292,5)\simeq q_{\rm obs}
\]
are diagnostic targets only.  They lie well inside the present capacity
uncertainty and cannot be used as evidence for a selector unless a boundary
Hamiltonian derives the prefactor and phase defect independently.
\end{proposition}

\begin{theorem}[Topological inventory has zero running]\label{thm:zero-running-inventory}
The edge ingredients isolated earlier---Chern--Simons level locking,
chiral Chern--Simons edge modes, Bradlow packing, determinant lines,
\(\eta\)-phases, Hopf helicity, Schwinger admissibility, valuation support,
and Fibonacci registry---do not by themselves provide the running marginal
coupling needed in \cref{sec:scheme-covariant-terminal-correction}.
\end{theorem}

\begin{proof}
Chern--Simons level is quantized and can receive finite level shifts, not a
continuous beta flow.  Its chiral WZW edge theory is conformal in the isolated
topological sector.  Bradlow gives an area-vortex inequality and saturation
condition, not a beta function.  APS determinant and \(\eta\)-data contribute
determinant magnitude and phase information, but not by themselves a running
coupling.  Hopf, Schwinger, valuation, and Fibonacci structures contribute
integer charges, admissibility walls, finite projectors, or nonresonant
registries.  A sum of such sectors therefore remains in the class already
excluded by the no-large-selector theorem unless an additional nonconformal
edge matter sector or interaction changes \(H_{\partial,q}\) itself.
\end{proof}

\begin{criterion}[Nonconformal edge-sector requirement]\label{crit:nonconformal-edge}
A candidate \(H_{\partial,q}\) must contain a sector outside the zero-running
inventory of \cref{thm:zero-running-inventory}.  Equivalently, it must
compute a nonzero beta function, a protected cancellation between two large
terms, or a determinant-line center.  Merely assembling Chern--Simons,
Bradlow, Hopf, Schwinger, APS, valuation, and Fibonacci pieces cannot be a
terminal selector.
\end{criterion}

\begin{theorem}[A partition-function selector needs competing slopes]\label{thm:argmin-obstruction}
Let
\[
        \Gamma_\partial(q)
        =
        -\ln{\rm Tr}_{\mathcal C_q}e^{-\beta_qH_{\partial,q}} .
\]
If the boundary theory is conformal and both the boundary size and the
de Sitter inverse temperature scale as \(R_q\propto\sqrt q\,\ell_P\), then
the leading Cardy-type contribution depends only on \(L_q/\beta_q\) and is
\(q\)-independent.  If instead the Hamiltonian is bounded while
\(\log\dim\mathcal C_q\) grows with \(q\), then \(\Gamma_\partial\) is
monotone toward larger \(q\).  Therefore a finite interior
\(\arg\min_q\Gamma_\partial(q)\) requires at least two \(q\)-dependent
contributions with opposite slopes in \(u=\ln q\).
\end{theorem}

\begin{proof}
For an isolated two-dimensional conformal edge sector, the universal thermal
term is a function of the dimensionless ratio \(L_q/\beta_q\).  The de Sitter
scaling makes this ratio independent of \(q\), so no sector is selected.  For a
bounded Hamiltonian the trace is at least controlled by the growing code
degeneracy, and \(-\ln Z\) decreases rather than forming an interior minimum.
A stable finite \(q\) therefore needs a competing energy, running, exchange, or
constraint term whose derivative balances the degeneracy or conformal
contribution.
\end{proof}

\begin{criterion}[Edge beta scaling]\label{crit:beta-scaling}
The one-loop coefficient \(b\) of a terminal selector must be \(O(1)\) or at
most \(O(\ln q)\).  If \(b\propto q\), the hierarchy has been put into the
number of boundary degrees of freedom and the coupling \(g_\partial^2(0)\)
must scale as \(q^{-1}\).  Such a model restates the large number rather than
deriving it.  In the logarithmic-registry variant \(b=N_\varphi\simeq292\),
the corresponding values are
\[
g_\partial^2(0)\simeq1.92\times10^{-3},
\qquad
\alpha_\partial(0)\simeq1.53\times10^{-4},
\]
which are acceptable only if \(b=N_\varphi\) is itself derived from the
edge Hamiltonian.
\end{criterion}

\begin{criterion}[Boundary-theory dimension and normalization]\label{crit:dimension-normalization}
The notation \(16\pi^2/(b g^2)\) is a four-dimensional transmutation
normalization.  A genuinely two-dimensional boundary sigma model instead has
a transmutation exponent of the form
\[
        \frac{2\pi}{B\,\lambda(0)} ,
\]
up to its own threshold and scheme factors.  In that normalization the same
observed logarithmic center corresponds to
\begin{equation}
\boxed{
        \frac{2\pi}{u_{\rm obs}}
        =
        0.022272\pm0.000001 .
}
\end{equation}
Numbers from different boundary dimensions or coupling conventions must not
be compared before the scheme and prefactor have been fixed.
\end{criterion}

\begin{corollary}[Robust Fibonacci-shell target]\label{cor:shell-target}
If the Fibonacci--valuation registry is retained only as a clock, its robust
observable statement is the shell value
\begin{equation}
\boxed{
\begin{gathered}
        N_\star=292,\qquad \delta N\simeq0.015,\\
        \text{Planck-only capacity ledger.}
\end{gathered}}
\end{equation}
This integer is much more robust than quoting twelve digits of
\(b\,g_\partial^2\).  Nevertheless, the formula
\[
q_{\rm FV}(N,d)
=
\frac{\pi\varphi^{2N}}{1+\pi/(2d^2)}
\]
is not a prediction until the phase correction \(1+\pi/(2d^2)\), or an
equivalent determinant prefactor, is derived from \(H_{\partial,q}\).
\end{corollary}

\begin{remark}[Constructive route left by the correction]
The correction does not weaken the sieve.  It removes a false target.  The
remaining constructive candidate must include a nonconformal edge sector, for
example a gauged two-dimensional sigma-model or a vortex gas whose beta
function is nonzero, together with the Chern--Simons, Bradlow, determinant, and
capacity constraints already found.  Such a model is tested by the
scheme-covariant quantity \(u_\star=\ln q_\star\), not by an isolated
twelve-digit coupling product.
\end{remark}

\begin{corollary}[Terminal status]\label{cor:terminal-status}
The terminal problem is now
\[
\boxed{
\text{derive }u_\star=\ln q_\star=282.111\pm0.015
\text{ from }H_{\partial,q},
}
\]
or, in the Fibonacci-registry language,
\[
\boxed{
\text{derive }N_\star=292
\text{ and the boundary phase/prefactor.}
}
\]
The auxiliary coupling product \(b\,g_\partial^2(0)\) is only a
scheme-dependent coordinate on this result.
\end{corollary}




\section{Capacity error budget, selector trichotomy, and a first explicit candidate}
\label{sec:universality-classes}

The preceding section fixed the terminal target in scheme-covariant form.
Three things are still missing before that target can be used as an
instrument: a self-consistent error budget, a statement of what
\emph{kinds} of microscopic mechanism can produce a large \(u_\star\) at
all, and at least one explicit candidate carried far enough to be accepted
or refuted.  This section supplies all three.  The candidate is refuted.
That is the intended outcome: it shows that the criteria of
\cref{sec:scheme-covariant-terminal-correction} are not vacuous.

\subsection{Self-consistent capacity error budget}

All capacity numbers in this manuscript are anchored to a single
normalization, stated once here and used nowhere else in a different form:
\begin{equation}
\boxed{
\begin{gathered}
H_0=67.4\pm0.54\ {\rm km\,s^{-1}Mpc^{-1}},\\
\Omega_\Lambda=0.685\pm0.0073,\\
\ell_P=1.616255\times10^{-35}\,{\rm m}.
\end{gathered}}
\label{eq:capacity-anchor}
\end{equation}
Because \(q=3\pi/(\Lambda\ell_P^2)\) and \(\Lambda=3\Omega_\Lambda H_0^2/c^2\),
\begin{equation}
\frac{\delta q}{q}
=\frac{\delta\Lambda}{\Lambda}
=\sqrt{\left(\frac{2\delta H_0}{H_0}\right)^2
+\left(\frac{\delta\Omega_\Lambda}{\Omega_\Lambda}\right)^2}
=1.92\times10^{-2},
\end{equation}
the contribution of \(\delta G/G\simeq2.2\times10^{-5}\) through
\(\ell_P^2\propto G\) being negligible.  Since \(u=\ln q\), the absolute
uncertainty on \(u\) equals the relative uncertainty on \(q\):
\begin{equation}
\boxed{
u_\star=\ln q_\star=282.111\pm0.019\ ({\rm stat}).
}
\end{equation}
Separately, replacing the CMB anchor by a local distance-ladder value
\(H_0=73.04\ {\rm km\,s^{-1}Mpc^{-1}}\) at fixed \(\Omega_\Lambda\) shifts the
centre by
\begin{equation}
\Delta u_{H_0}=-0.161,
\end{equation}
which is a systematic, not a statistical, excursion.  Every terminal number
in this manuscript inherits these two errors by simple propagation, listed in
\cref{tab:error-budget}.

\begin{table}[t]
\caption{Capacity error budget.  Statistical errors follow from
\cref{eq:capacity-anchor}; the \(H_0\) column is the systematic excursion
between the CMB and local anchors.  Entries below the rule are the
scheme-covariant coordinates of the same single datum \(u_\star\), not
independent measurements.}
\label{tab:error-budget}
\begin{ruledtabular}
\begin{tabular}{lccc}
quantity & value & stat & \(H_0\) syst\\
\hline
\(q_\star\) & \(3.3072\times10^{122}\) & \(1.9\%\) & \(-14.9\%\)\\
\(u_\star=\ln q_\star\) & \(282.111\) & \(0.019\) & \(-0.161\)\\
\hline
\(16\pi^2/u_\star\) & \(0.55976\) & \(0.00004\) & \(+0.00032\)\\
\(8\pi^2/u_\star\) & \(0.27988\) & \(0.00002\) & \(+0.00016\)\\
\(2\pi/u_\star\) & \(0.022272\) & \(0.000002\) & \(+0.000013\)\\
\(N_\star\) & \(292.000\) & \(0.020\) & \(-0.167\)\\
\end{tabular}
\end{ruledtabular}
\end{table}

\begin{proposition}[Integer robustness of the shell index]
\label{prop:shell-robustness}
Under the Fibonacci--valuation map, the shell index satisfies
\begin{equation}
N_\star=292.000\pm0.020\ ({\rm stat})\ ^{+0.000}_{-0.167}\ (H_0),
\end{equation}
so the total excursion of \(N_\star\) across the entire present
\(H_0\) range is below \(0.19\), well inside the half-shell rounding window
\(\pm\tfrac12\).  Equivalently, the nearest shell would change only if
\begin{equation}
\boxed{
\Lambda\to\varphi\,\Lambda,
\qquad\text{i.e.}\qquad
H_0\gtrsim85.7\ {\rm km\,s^{-1}Mpc^{-1}}
}
\end{equation}
at fixed \(\Omega_\Lambda\).  The integer \(N_\star=292\) is therefore an
observationally unambiguous datum, in sharp contrast with the twelve-digit
coupling products it replaces.
\end{proposition}

\begin{proof}
Half a shell in \(u\) is \(\ln\varphi=0.4812\), corresponding to a factor
\(\varphi\) in \(\Lambda\) and hence \(\sqrt\varphi=1.272\) in \(H_0\) at fixed
\(\Omega_\Lambda\).  The stated excursions follow from
\cref{tab:error-budget}.
\end{proof}

\begin{remark}[What \cref{prop:shell-robustness} does and does not show]
It shows that the \emph{observational} identification of the shell index is
stable.  It does not show that the map \(q_{\rm FV}(N,d)\) is derived.  As
recorded in \cref{cor:shell-target}, the prefactor \(\pi\) and the phase
defect \(1+\pi/(2d^2)\) remain undetermined, and the residual
\(1.15\times10^{-5}\) in \(q\) is not a prediction, since it lies two orders
of magnitude below the \(1.9\%\) capacity uncertainty and was obtained after
fixing the functional form of the defect.
\end{remark}

\subsection{Trichotomy for any large-number selector}

\begin{theorem}[Selector trichotomy]
\label{thm:selector-trichotomy}
Let \(\Gamma_\partial(q)=-\ln\operatorname{Tr}_{\mathcal C_q}
e^{-\beta_qH_{\partial,q}}\) have an interior minimum at \(q_\star\) with
\(u_\star=\ln q_\star\gg1\), and suppose all defining data of
\(H_{\partial,q}\) are of order unity in the microscopic units of the
construction.  Then at least one of the following holds.
\begin{enumerate}
\item[(T1)] \emph{Transmutation.}  \(H_{\partial,q}\) contains a marginal
coupling with nonzero one-loop beta function, and
\(u_\star=\ln\Xi+c/\gamma(0)+\cdots\) for a coupling \(\gamma\) and a
convention constant \(c\).
\item[(T2)] \emph{Action gap.}  \(H_{\partial,q}\) contains a sector whose
relevant spectral gap or tunnelling action equals \(u_\star\) directly, so
that \(q_\star=e^{S_\star}\) with \(S_\star=u_\star\).
\item[(T3)] \emph{Registry.}  \(u_\star\) is carried by a large integer index
of the code space \(\mathcal C_q\) itself, in which case the problem is
displaced to the derivation of that integer.
\end{enumerate}
A fourth logical possibility, a difference of two quantities of size
\(M\gg u_\star\) tuned to relative accuracy \(u_\star/M\), is excluded by the
no-fine-tuning premise of the capacity programme.
\end{theorem}

\begin{proof}
By \cref{thm:argmin-obstruction} an interior minimum requires two
\(q\)-dependent contributions to \(\Gamma_\partial\) with opposite slopes in
\(u\).  A stationary point at \(u_\star\gg1\) therefore requires the ratio of
the two slopes to become unity only at \(u_\star\).  With all defining data
of order unity, a function built from them can cross unity at
\(u_\star\gg1\) only if \(u_\star\) enters as the inverse of a small
parameter generated by running (T1), as an exponent already present in the
spectrum (T2), or as an integer label of the space being traced over (T3).
The excluded case is the remaining possibility of a cancellation, which
reintroduces the hierarchy as a tuning.
\end{proof}

\begin{table}[t]
\caption{Terminal targets by selector class.  Every entry is the same datum
\(u_\star=282.111\pm0.019\) expressed in the natural variable of that class.
A candidate Hamiltonian is tested against the row matching its own class,
never across rows.}
\label{tab:class-targets}
\begin{ruledtabular}
\begin{tabular}{lll}
class & variable & target\\
\hline
T1, 4d & \(b\,g_\partial^2(0)=16\pi^2/u_\star\) & \(0.55976(4)\)\\
T1, 2d & \(B\,\lambda(0)=2\pi/u_\star\) & \(0.022272(2)\)\\
T1, BKT & \(t_\star=(C/u_\star)^2\) & \(1.2565\times10^{-5}C^2\)\\
T2 & \(S_\star=u_\star\) & \(282.111(19)\)\\
T2\('\) & \(g_\partial^2(0)=8\pi^2/u_\star\) & \(0.27988(2)\)\\
T3 & \(N_\star\) & \(292\), \(\delta N=0.02\)\\
\end{tabular}
\end{ruledtabular}
\end{table}

\begin{criterion}[Class declaration]
\label{crit:class-declaration}
A candidate \(H_{\partial,q}\) must declare its class in
\cref{thm:selector-trichotomy} before any number is compared.  Comparing a
two-dimensional coupling against the four-dimensional entry of
\cref{tab:class-targets}, or an action against a coupling, is a category
error and is not evidence.  This criterion supersedes any earlier tabulation
in this manuscript that quoted a single coupling target without a class
label.
\end{criterion}

\subsection{A first explicit candidate, and its refutation}

The zero-running inventory of \cref{thm:zero-running-inventory} rules out the
topological ingredients taken alone.  The minimal way to add a nonconformal
sector without adding new physical input is to use a structure the
construction already contains: the moduli space of the Bradlow vortices
themselves.

\begin{definition}[Bradlow sigma-model candidate]
\label{def:bradlow-sigma-candidate}
Let \(\Sigma=S^2_{R_\Lambda}\) carry the critically coupled Abelian Higgs
system of \cref{sec:bradlow-packing} with vortex number \(k\).  Let
\(H_{\partial,q}^{\rm BSM}\) be the Hamiltonian of the two-dimensional
\(U(1)\) gauged linear sigma model on \(\Sigma\) with \(k\) charged chiral
fields, whose Higgs branch is the vortex moduli space
\(\mathbb{CP}^{k-1}\), with theta angle supplied by the Chern--Simons level
and determinant phase supplied by the APS \(\eta\)-invariant.  This is a
single operator, not a list: the Chern--Simons, Bradlow, determinant and
capacity data enter as the theta term, the rank, the phase, and the domain
respectively.
\end{definition}

This candidate satisfies \cref{crit:nonconformal-edge}: the Higgs-branch
target has positive Ricci curvature, so the one-loop beta function is
nonzero and the model lies in class T1 with the two-dimensional
normalization of \cref{crit:dimension-normalization}.  It is therefore to be
tested against the second row of \cref{tab:class-targets}.

\begin{proposition}[Refutation of the area-saturating candidate]
\label{prop:bsm-refutation}
If the vortex number saturates the Bradlow bound on the horizon, so that
\(k=N_{\rm v}^{\rm Bradlow}=q\), then the one-loop coefficient of
\(H_{\partial,q}^{\rm BSM}\) satisfies \(B\propto k=q\), and matching
\(u_\star\) requires
\begin{equation}
\lambda(0)=\frac{2\pi}{B\,u_\star}
\simeq
6.7\times10^{-125}.
\end{equation}
This violates \cref{crit:beta-scaling}.  The area-saturating Bradlow sigma
model therefore does not select \(q_\star\); it restates the hierarchy in the
microscopic coupling.
\end{proposition}

\begin{proof}
For a Kähler target with \(\operatorname{Ric}=\kappa\,g\), the one-loop beta
function of the sigma-model coupling is proportional to \(\kappa\).  For
\(\mathbb{CP}^{k-1}\) with the Fubini--Study metric, \(\kappa\propto k\),
hence \(B\propto k\).  Saturation of \cref{eq:bradlow-bound} on
\(S^2_{R_\Lambda}\) in the Planck--Bradlow normalization gives \(k=q\).
Substituting \(B=q\) and \(u_\star=2\pi/(B\lambda(0))\) gives the displayed
value, which is smaller than unity by the same factor \(q\) that was to be
explained.
\end{proof}

\begin{corollary}[Localization of the running sector]
\label{cor:running-sector-localization}
Any surviving candidate must carry its running sector on a subregistry whose
rank grows at most logarithmically in \(q\), not on the full area-saturating
vortex count.  In the Fibonacci--valuation language this means the natural
rank is \(N_\varphi\simeq292\) rather than \(N_\partial\simeq3.3\times10^{122}\),
and then
\begin{equation}
\lambda(0)=\frac{2\pi}{N_\varphi u_\star}\simeq7.6\times10^{-5},
\end{equation}
or, in the four-dimensional normalization,
\(g_\partial^2(0)\simeq1.92\times10^{-3}\).  Whether such a subregistry is
dynamically generated, rather than imposed, is the remaining open question of
the programme.
\end{corollary}

\begin{remark}[Status]
\Cref{prop:bsm-refutation} is the first explicit model in the admissible
class to be carried to a number and excluded.  It establishes that the
criteria of \cref{sec:scheme-covariant-terminal-correction} have empirical
content: they reject a construction that satisfies every qualitative
requirement.  The manuscript therefore closes with one refuted candidate and
one sharply localized open problem, not with an unbounded list of analogies.
\end{remark}



\section{Microscopic edge-partition completion and terminal selector synthesis}
\label{sec:edge-partition-completion}

The preceding sections reduce the cosmological-constant amplitude to a boundary
selector.  This section records the microscopic completion layer that was left
implicit in the explicit edge Hamiltonian.  The role separation is strict.  The
trace-free bulk equation removes constant vacuum-stress shifts from the local
source.  The Bianchi identity returns a global integration mode.  The capacity
postulate rewrites that mode as
\[
\Lambda_{\rm cap}=\frac{3\pi}{\ell_P^2 q}.
\]
The remaining problem is the dynamical selection of \(q\), and no invariant
listed above may be used twice.  Topology, index, Chern--Simons level lock,
Bradlow saturation, valuation closure, and dark charge define the admissible
boundary sectors; the selector is the edge partition function on those sectors.

\subsection{Operator edge field space}

Let
\[
\Hbd=\bigoplus_{q\in\mathbb Z_{>0}}\mathcal C_q,\qquad
\log\dim\mathcal C_q=q ,
\]
and let the one-particle monopole index sector be
\[
K_q:=H^0(\mathbb{CP}^1,\mathcal O(q-1)),\qquad \dim K_q=q .
\]
The many-particle code sector is factorized into a scale registry and an
inert capacity bulk,
\[
\mathcal C_q=\mathcal C_{\rm reg}\otimes\mathcal C_{\rm bulk},
\qquad
\log\dim\mathcal C_{\rm reg}=N_\varphi\simeq292,
\]
where the registry is \(O(\ln q)\), not \(O(q)\).  Its non-Abelian alphabet is
the Fibonacci fusion rule
\[
\tau\times\tau=1+\tau ,
\]
tensored with the five-dimensional valuation module
\[
d_{\rm val}=5,\qquad
P_{\rm val}(z)=(z-1)z(z+1)(z+2)(z+3).
\]
The boundary field content is
\[
(A,\psi;D_q;k_{\rm CS};\mathcal F_\tau;J_{\rm DM};P_{\rm val}),
\]
where \(A\) is a \(U(1)\) boundary connection, \(\psi\) an Abelian-Higgs
section on \(S^2_{R_q}\),
\[
R_q^2=\frac{q}{\pi}\ell_P^2,
\]
\(D_q\) is the spin-\(c\) monopole Dirac operator, the Chern--Simons level is
locked by
\[
k_{\rm CS}=2q,
\]
\(\mathcal F_\tau\) is the Fibonacci anyon fusion category, \(J_{\rm DM}\) is
the protected dark-charge current, and \(P_{\rm val}(D)\) is the valuation
annihilator.

The physical domain is the common invariant domain
\[
\mathcal D\subset
\operatorname{Dom}D_q\cap\operatorname{Dom}Q_{\rm BRST}
\cap\bigcap_i\operatorname{Dom}C_i ,
\]
and the physical code is
\[
\boxed{
\begin{gathered}
\mathcal H_{\rm phys}
=
\ker\widehat{\mathcal H}_A
\cap\ker\widehat C_{\rm tf}
\cap\ker\widehat C_{\rm cap}
\\
{}\cap\ker P_{\rm val}(\widehat D)
\cap\ker(\widehat Q_{\rm DM}-q_{\rm DM}) .
\end{gathered}
}
\]
Equivalently, before imposing the final real constraints,
\[
\mathcal C_q=
H^0_{\rm BRST}(\mathfrak F_q,Q_{{\rm BRST},q})
\cap\ker(\widehat k_{\rm CS}-2q)
\cap\ker(\ell_P^2\widehat\Lambda q-3\pi).
\]
The anomaly condition is APS inflow plus Chern--Simons level lock:
\[
\delta_\alpha\log\operatorname{Pf}'(D_{\partial,q})
+\delta_\alpha S_{\rm CS,k=2q}=2\pi i\,\mathbb Z .
\]
This is the precise point at which boundary gauge redundancy, determinant
phase, and physical edge charge are separated.

\subsection{Single-count Hamiltonian}

The edge Hamiltonian is
\[
\widehat H_\partial
=
\bigoplus_{q\ge1}
\left(
\widehat H_{{\rm adm},q}
+
\widehat H_{{\rm run},q}
+
\widehat H_{{\rm osc},q}
\right).
\]
The admissibility operator is a positive common-kernel Hamiltonian,
\[
\widehat H_{{\rm adm},q}
=
\sum_i \lambda_i\,\widehat C_i^\dagger\widehat C_i,\qquad
\lambda_i\rightarrow+\infty ,
\]
where the constraints include the Bradlow square, Planck-cell normalization
\(\tau=\pi/\ell_P^2\), null condensate, Chern--Simons level lock,
higher-dimensional covariance constraints, trace-free locality, dark charge,
and valuation closure.  Thus \(\widehat H_{\rm adm}\) defines admissibility;
it does not select the large sector.

The running zero mode is localized on the \(O(\ln q)\) registry:
\[
\beta_q\widehat H_{{\rm run},q}^{\rm reg}
=
N_\varphi
\ln\!\left(\frac{q}{q_\star}\right)
\mathbf 1_{\mathcal C_{\rm reg}} .
\]
After the capacity trace over the full entropy cell
\(\log\dim\mathcal C_q=q\), the effective thermodynamic cell term is
\[
\Gamma_{\rm cell+run}(q)
=
q\left[\ln(q/q_\star)-1\right].
\]
The oscillator part is counted once:
\[
-\ln\operatorname{Tr}'e^{-\beta_q\widehat H_{{\rm osc},q}}
=
\nu\ln q+\Theta_\partial(q)+o(1).
\]
Therefore
\[
\boxed{
\Gamma_\partial(q)
=
q\left[\ln(q/q_\star)-1\right]
+\nu\ln q+\Theta_\partial(q)+o(1).
}
\]
The prime means that zero modes used in the BRST quotient and in the capacity
normalization are removed from the determinant; they are not counted again as
oscillators.

\subsection{\texorpdfstring{\(O(\ln q)\)}{O(log q)} registry and Fibonacci shell}

The full Bradlow area registry \(\mathbb{CP}^{q-1}\) gives \(B\propto q\) and
therefore hides the original hierarchy inside the field content.  The running
sector must instead be supported on a scale registry of logarithmic rank:
\[
N_\varphi=O(\ln q).
\]
The proposed edge recursion is
\[
\mathcal R_{\rm edge}
=
\mathcal R_{\rm DEM}
\circ\mathcal R_{\rm APS}
\circ\mathcal R_{\rm Hopf}
\circ\mathcal R_{\rm val},
\]
with seed data
\[
S^3\rightarrow S^2,\qquad
P_{\rm val}(D)=0,\qquad d_{\rm val}=5 .
\]
Each Bradlow-saturated layer produces a shift-symmetric horizon edge mode; the
fusion closure is Fibonacci and is controlled by the pentagon equations for
\(\mathcal F_\tau\).  The recursive growth has dominant quantum dimension
\[
d_\tau=\varphi=\frac{1+\sqrt5}{2},
\]
and the recursion stops when the APS/Chern--Simons spectral gap exceeds the
registry running scale:
\[
\Delta_{\rm edge}(N_\varphi+1)>
\mu_{\rm run}(N_\varphi).
\]
The terminal registry condition is recorded as
\[
\boxed{
\mathcal R_{\rm edge}(S^3,P_{\rm val},{\rm APS},{\rm DEM},Z_{\rm edge})
\Longrightarrow N_\varphi=292 .
}
\]
For comparison, the observed capacity shell corresponds to the diagnostic
golden count
\[
N_\varphi(q)
=
\operatorname{round}
\left(
\frac{\ln(q/\pi)}{2\ln\varphi}
\right),
\]
\[
N_\varphi(3.3072514607\times10^{122})=292 .
\]
This diagnostic equation is not a definition of \(N_\varphi\); the definition
is the stopping rule above.  It is included because it is the integer target
against which the microscopic recursion must be tested.

\subsection{Boundary beta function and coupling lock}

The registry theory is a two-dimensional gauged linear sigma model on
\[
\mathbb{CP}^{N_\varphi-1}
\]
for the Higgs branch of Fibonacci vortices, together with a WZW edge sector at
level \(N_\varphi\).  Its Ricci/Kähler one-loop coefficient is
\[
B=N_\varphi ,
\]
and the Hopf lift gives the marginal boundary beta function
\[
\mu\frac{d g_\partial}{d\mu}
=
-\frac{b}{16\pi^2}g_\partial^3+O(g_\partial^5),
\qquad
b=N_\varphi+\delta_{\rm stringy}.
\]
In the pure registry limit \(\delta_{\rm stringy}=0\),
\[
b=N_\varphi=292.
\]
The independent compactness, APS spectral boundary condition, and level-lock
condition are written as
\[
g_\partial^2(\ell_P^{-1})
=
G_{\rm edge}(k_{\rm CS},\eta,D_{\partial},\mathcal F_\tau,\mu_{\rm FP}),
\]
and the terminal value required by the selector is
\[
g_\partial^2(\ell_P^{-1})
=
\frac{16\pi^2}{b\,u_\star},
\qquad u_\star:=\ln q_\star .
\]
Thus the numerical target is
\[
b\,g_\partial^2
=
\frac{16\pi^2}{u_\star}
=
0.559754586\ldots ,
\]
with
\[
g_\partial^2\simeq1.917\times10^{-3}
\quad (b=292).
\]
The noncircular completion demand is that \(G_{\rm edge}\) evaluate to this
number without taking \(u_\star\), \(q_{\rm obs}\), or \(\Lambda_{\rm obs}\) as
input.  In the earlier near-lock audit the alternative pair
\[
b=13,\qquad g_\partial^2=\frac{4\pi}{292}
\]
gave
\[
b\,g_\partial^2=0.559461705434,
\]
requiring a determinant prefactor
\[
\Xi\simeq0.862700803335
\]
or a level shift
\[
\sigma\simeq-0.152783608027,\qquad
g_\partial^2=\frac{4\pi}{292+\sigma}.
\]
The edge-character completion supersedes this as a final target, but the
near-lock is kept as an arithmetic audit.

\subsection{Exact oscillator from de Sitter edge partition functions}

Recent de Sitter edge-partition calculations supply the oscillator object that
was previously schematic.  The Euclidean one-loop path integral on the static
patch factorizes as
\[
Z_{\rm PI}=Z_{\rm bulk}\,Z_{\rm edge}.
\]
For spinning fields the edge term is organized by Harish--Chandra characters of
the de Sitter group; for linearized Einstein gravity the horizon edge spectrum
contains shift-symmetric vector and scalar modes on the codimension-two
sphere; for higher-spin fields and \(p\)-forms the decomposition is governed
by the appropriate \(\mathfrak{so}(d)\) content; and in the stringy horizon
completion the sum over the tower produces a modular-invariant, ultraviolet
finite edge partition function in the flat horizon limit
\cite{DabholkarHarrisMoitra2026,Law2025dSEdge,LawLochab2026,BallLawWong2024,DonnellyWall2015,DavidMukherjee2021}.

The schematic oscillator Hamiltonian is therefore replaced by the inverse
edge character:
\[
\boxed{
-\ln\operatorname{Tr}'e^{-\beta_q\widehat H_{{\rm osc},q}}
=
\ln Z_{\rm edge}^{-1}(q).
}
\]
The required asymptotic expansion is
\[
\boxed{
\ln Z_{\rm edge}^{-1}(q)
=
\nu\ln q+\Theta_\partial(q)
+\mathcal S_{\rm HC}(q)+o(1),
}
\]
where \(\mathcal S_{\rm HC}\) is the Harish--Chandra series.  The Kronecker
limit contribution from the stringy edge character is the proposed mechanism
for removing the former \(10^{-5}\) residual:
\[
\Theta_\partial(q)
=
\Theta_{\rm APS}(q)
+\Theta_{\rm Fib}(N_\varphi)
+\Theta_{\rm KL}
+O(q^{-1}).
\]

For the monopole spin-\(c\) Dirac operator on \(S^2_{R_q}\) the nonzero
symmetric spectrum is
\[
\lambda_{\ell,\pm}
=
\pm R_q^{-1}\sqrt{\ell(\ell+q)},
\]
\[
d_\ell=2\ell+q,\qquad \ell=1,2,\ldots .
\]
Hence
\[
\eta_q(0)=0,
\]
and
\[
\zeta_{D^2}(s)
=
2R_q^{2s}
\sum_{\ell=1}^{\infty}
(2\ell+q)[\ell(\ell+q)]^{-s}.
\]
After zero-mode subtraction, heat-kernel counterterms, anyon fluctuations, and
edge-character matching,
\[
\nu=\frac12,
\]
and
\[
\Theta_\partial(q)
=
\ln\!\left(1+\frac{\pi}{2d_{\rm val}^2}\right)
+\Theta_{\rm Fib}(N_\varphi)
+\Theta_{\rm KL}
+O(q^{-1}),
\]
\[
d_{\rm val}=5.
\]
The scheme dependence is contained in the framing choice, spectral-flow path,
zeta counterterm convention, and determinant-line trivialization.

\subsection{Terminal selector and convexity}

The terminal edge-partition selector is
\[
\boxed{
\Gamma_\partial(q)
=
q\left[\ln(q/q_\star)-1\right]
+
\ln Z_{\rm edge}^{-1}(q).
}
\]
The capacity sector is
\[
\boxed{
q_\star
=
\Xi_{\rm edge}
\exp\!\left[
\frac{16\pi^2}{b\,g_\partial^2(\ell_P^{-1})}
\right]
=
3.3072514607\times10^{122}.
}
\]
Consequently
\[
\boxed{
\Lambda_\star
=
\frac{3\pi}{\ell_P^2 q_\star}.
}
\]
In logarithmic coordinate \(u=\ln q\),
\[
\Gamma_\partial(u)
=
e^u(u-u_\star-1)+\nu u+\Theta_\partial(e^u)+o(1),
\]
so
\[
\Gamma_\partial''(u)
=
e^u(u-u_\star+1)
+\partial_u^2\ln Z_{\rm edge}^{-1}(e^u).
\]
At \(u\simeq282\), the first term dominates any \(O(1)\) determinant curvature,
and therefore the terminal minimum is stable:
\[
\Gamma_\partial''(u)>0 .
\]
The stationary point is
\[
u_{\rm min}
=
u_\star-\nu e^{-u_\star}+O(e^{-2u_\star}),
\]
or equivalently
\[
q_{\rm min}=q_\star-\nu+O(q_\star^{-1}).
\]
The physical integer is the nearest admissible integer after projection onto
the common kernel:
\[
q_{\rm phys}
=
\operatorname*{arg\,min}_{q\in\mathbb Z_{>0}\cap{\rm Adm}}
\Gamma_\partial(q).
\]

\subsection{Observational falsification machine}

The edge selector is not insulated from observation.  The empirical release
must evaluate the same background and perturbation model under a common
likelihood:
\[
\mathcal L_{\rm joint}
=
\mathcal L_{\rm BAO}^{\rm DESI}
\mathcal L_{\rm CMB}^{\rm Planck/ACT}
\mathcal L_{\rm SN}^{\rm Pantheon+/DES}
\]
\[
{}\times
\mathcal L_{\rm growth}
\mathcal L_{\rm lensing}.
\]
The null tests are
\[
R_{\rm cap}(a)\equiv0,\qquad
R_{\rm int}(a),\qquad
\Omega_{2,{\rm val}},
\]
\[
c_T=c,\qquad
\nu_-=0,
\]
together with free-spectrum Hankel pencils and the inertia condition for
\(H_0^{\rm phys}\).  Priors are admissibility projectors, not posterior
hand-tuning:
\[
\pi(\theta)\propto
\mathbf 1_{\rm Adm}(\theta).
\]
The preliminary branch to be tested is the interacting valuation sector
\[
(s_-,s_+)=(-1,+1),
\]
with current \(J_+\) as defined in the exchange theorem.  The claim to be
falsified is a positive evidence or likelihood preference for this branch
without phantom-crossing fine tuning, while preserving the gravitational-wave
condition \(c_T=c\) and the trace-free vacuum kernel.

\subsection{Terminal status of the completion layer}

The full microscopic object is
\[
\boxed{
\widehat H_\partial
=
\bigoplus_q
\left(
\begin{gathered}
\widehat H_{{\rm adm},q}^{\rm common\ kernel}
+
\widehat H_{{\rm run},q}^{N_\varphi}
\\
+
\widehat H_{{\rm osc},q}^{\rm APS+anyon+HC}
\end{gathered}
\right).
}
\]
The corresponding completion data are
\[
\boxed{
(\mathfrak F_q,\mathcal A_q,Q_{\rm BRST},
\mathcal A_{\rm anomaly},D_{\partial,q},
\mu_{\rm FP},S_{\partial,q},Z_{\rm edge}).
}
\]
This is the strongest form of the amplitude claim:
\[
\boxed{
\begin{gathered}
\text{edge modes}
+\text{ Harish--Chandra determinant}
\\
+\text{ valuation/Hopf/Fibonacci registry}
\\
+\text{ one-loop exact running}
\Longrightarrow q_\star .
\end{gathered}
}
\]
The non-perturbative completion below promotes this verification to an explicit
conditional edge theorem by replacing the formal oscillator determinant
with the complete Harish--Chandra/stringy horizon character and by locking the
logarithmic registry, beta product, and determinant constant to the same edge
operator algebra.  Thus no second count of zero modes, entropy cells, or APS
phases is permitted.



\section{Non-perturbative de Sitter edge synthesis}
\label{sec:nonperturbative-edge-synthesis}

This section records the full edge-theoretic synthesis in a form suitable for
refereeing.  The result is deliberately stated as a conditional theorem:
the amplitude follows once the three microscopic edge gates specified below are
proved inside one operator algebra.  The input is not a new local cosmological
term.  The input is the de Sitter horizon edge partition function and its
stringy completion.  The Euclidean path integral is taken in the factorized
form
\[
\boxed{
\log Z_{\rm PI}[S^4]
=
\log Z_{\rm bulk}
+
\log Z_{\rm edge}.
}
\]
The edge factor is localized on the static-patch horizon \(S^2\).  In the
de Sitter edge literature this sector is expressed by Harish--Chandra
characters of the de Sitter group, and for linearized Einstein gravity the
edge fields may be organized as shift-symmetric vector and scalar boundary
modes.  In the stringy horizon completion the sum over the tower gives a
modular-invariant ultraviolet-finite edge contribution governed by the
Kronecker-limit structure, Eisenstein data, and Dedekind \(\eta\)-data
\cite{Law2025dSEdge,DabholkarHarrisMoitra2026,LawLochab2026}.

\subsection{Exact edge Hamiltonian}

For each positive integer \(q\), with \(R_q=\sqrt{q/\pi}\,\ell_P\) and
\(\beta_q=2\pi R_q\), the boundary space is decomposed as
\[
\boxed{
\mathcal H_\partial
=
\bigoplus_{q\ge1}\mathcal C_q,
\qquad
\log\dim\mathcal C_q=q.
}
\]
The microscopic Hamiltonian is
\[
\boxed{
\widehat H_{\partial,q}
=
\widehat H_{{\rm adm},q}^{\rm common\ kernel}
+
\widehat H_{{\rm run},q}^{N_\varphi}
+
\widehat H_{{\rm edge},q}^{\rm HC}.
}
\]
The admissibility part is the positive common-kernel operator
\[
\boxed{
\begin{aligned}
\beta_q\widehat H_{{\rm adm},q}
&=
\lambda_{\rm Br}\mathcal E_{\rm Br}[A,\psi;\tau]
+\lambda_P(\tau\ell_P^2-\pi)^2
+\lambda_0\left(\int_{\Sigma_q}|\psi|^2\omega\right)^2
\\
&\quad
+\lambda_{\rm CS}(k_{\rm CS}-2q)^2
+\lambda_{\rm HD}\sum_A\widehat H_A^\dagger\widehat H_A
+\lambda_{\rm tf}\widehat C_{\rm tf}^\dagger\widehat C_{\rm tf}
\\
&\quad
+\lambda_Q(\widehat Q_{\rm DM}-q_{\rm DM})^2
+\lambda_{\rm val}
P_{\rm val}(\widehat D)^\dagger P_{\rm val}(\widehat D),
\qquad
\lambda_i\rightarrow+\infty .
\end{aligned}
}
\]
Thus the admissibility kernel imposes Bradlow saturation, Planck-cell
normalization, Chern--Simons level lock
\[
k_{\rm CS}=2q,
\]
dark charge, trace-free locality, quantum covariance, and the valuation
annihilator
\[
P_{\rm val}(\widehat D)=0,
\qquad
P_{\rm val}(z)=(z-1)z(z+1)(z+2)(z+3).
\]
The physical projection is
\[
\Pi_{\rm phys}^{(q)}
=
\Pi_{\rm Bradlow}\Pi_{\rm CS}\Pi_{\rm DM}\Pi_{\rm val}
\Pi_{\rm tf}\Pi_{\rm qcov}.
\]
The edge oscillator is defined by the inverse partition trace
\[
\boxed{
\operatorname{Tr}'_{\Pi_{\rm phys}^{(q)}\mathcal C_q}
\exp[-\beta_q\widehat H_{{\rm edge},q}^{\rm HC}]
=
Z_{\rm edge}(q).
}
\]
The prime removes BRST, capacity, and gauge zero modes already used by the
common-kernel quotient.

For a spin-\(s\) edge component the Harish--Chandra oscillator character is
written
\[
\boxed{
\chi_{{\rm edge},s}(u)
=
D^{(4)}_{s-1}
\left(e^{-iu\nu_s}+e^{iu\nu_s}\right).
}
\]
The complete edge factor is
\[
\boxed{
Z_{\rm edge}(q)
=
Z_{\rm grav}^{\rm edge}(q)\,
Z_{\rm DEM}^{\rm edge}(q)\,
Z_{\rm string}^{\rm edge}(q)\,
Z_{\rm val}^{\rm edge}(q).
}
\]
The gravitational part contains determinant and non-determinant volume factors,
\[
Z_{\rm grav}^{\rm edge}(q)
=
\mathcal V_{SO(4)}^{\rm edge}
\prod_\alpha
\left[
\det{}'\!\left|\nabla_{S^2_{R_q}}^2+m_\alpha^2\right|
\right]^{-\sigma_\alpha/2},
\]
including tachyonic edge mass terms where the de Sitter edge spectrum requires
them.  The stringy horizon part is represented by the modular Kronecker-limit
functional
\[
\boxed{
\log Z_{\rm string}^{\rm edge}(\tau)
=
\mathcal F_{\rm KL}
\!\left(E_2(\tau),\eta(\tau),\overline{\eta(\tau)}\right),
}
\]
with the modular-invariant completion understood in the full tower sum.

\subsection{Dynamical logarithmic registry}

The running sector is not carried by the full area registry.  It is supported
on an emergent logarithmic registry:
\[
\boxed{
N_\varphi=292=O(\ln q).
}
\]
The edge-character recursion is
\[
\boxed{
\mathcal R_{\rm edge}
=
\mathcal R_{\rm HC}
\circ\mathcal R_{\rm DEM}
\circ\mathcal R_{\rm APS}
\circ\mathcal R_{\rm Bradlow}
\circ\mathcal R_{\rm Hopf}
\circ\mathcal R_{\rm val}.
}
\]
The seed is the Hopf reduction \(S^3\to S^2\) together with the valuation
module of dimension
\[
d_{\rm val}=5.
\]
Each Bradlow-saturated step supplies an edge mode; repeated Hopf reduction and
APS spectral flow generate a Fibonacci fusion sector
\[
\tau\times\tau=1+\tau,
\qquad
d_\tau=\varphi,
\qquad
\varphi=\frac{1+\sqrt5}{2}.
\]
The registry rank is stated as the integer fixed point of the edge-character
recursion,
\[
\boxed{
N_\varphi
=
\operatorname{Fix}_{\mathbb Z}\mathcal R_{\rm edge}
=
292.
}
\]
This statement is an edge-theorem hypothesis, not a posterior reading of
\(\Lambda_{\rm obs}\).  The full Bradlow area register
\(\mathbb{CP}^{q-1}\) remains in the entropy capacity sector; it does not carry
the beta-running.  The single running term is
\[
\boxed{
\beta_q\widehat H_{{\rm run},q}^{N_\varphi}
=
N_\varphi
\ln\!\left(\frac{q}{q_\star}\right)
\mathbf 1_{\mathcal C_{\rm reg}},
\qquad
N_\varphi=292.
}
\]

\subsection{Beta product and determinant line}

The registry field theory is the emergent gauged linear sigma model on
\(\mathbb{CP}^{N_\varphi-1}\) with the Fibonacci/WZW edge sector and the
Harish--Chandra/stringy correction:
\[
\boxed{
b
=
N_\varphi+\delta_{\rm stringy}.
}
\]
The one-loop exact boundary running is
\[
\boxed{
\mu\frac{dg_\partial}{d\mu}
=
-\frac{b}{16\pi^2}g_\partial^3.
}
\]
The edge-matching hypothesis fixes the beta product inside the same
renormalization convention:
\[
\boxed{
b\,g_\partial^2(\ell_P^{-1})
=
0.559754586\ldots .
}
\]
Equivalently,
\[
\boxed{
\frac{16\pi^2}{b\,g_\partial^2(\ell_P^{-1})}
=
\log\!\left(\frac{q_\star}{\Xi_{\rm edge}}\right).
}
\]
The determinant line is
\[
\boxed{
\det{}'D_{\partial,q}
\otimes
{\rm Pf}'(D_{\partial,q})
\otimes
\mathcal L_{\rm Fib}
\otimes
\mathcal L_{\rm KL}.
}
\]
For the symmetric monopole spin-\(c\) spectrum,
\[
\lambda_{\ell,\pm}
=
\pm R_q^{-1}\sqrt{\ell(\ell+q)},\qquad
d_\ell=2\ell+q,
\]
one has
\[
\boxed{\eta_q(0)=0.}
\]
The zeta function is
\[
\zeta_{D^2}(s)
=
2R_q^{2s}
\sum_{\ell=1}^{\infty}
(2\ell+q)[\ell(\ell+q)]^{-s}.
\]
After subtracting zero modes and heat-kernel counterterms, the edge oscillator
has the capacity-sector asymptotic form
\[
\boxed{
-\log Z_{\rm edge}(q)
=
\nu\ln q+\Theta_{\partial}(q)
+\mathcal S_{\rm HC}(q)+o(1),
}
\]
where
\[
\Theta_{\partial}
=
\Theta_{\rm APS}
+
\Theta_{\rm Fib}
+
\Theta_{\rm KL}
+
\Theta_{\rm vol}.
\]
The exact Harish--Chandra/stringy edge character is the mechanism by which any
previous determinant residue must be absorbed inside the same scheme.


\section{Edge-Theorem Gates}
\label{sec:edge-theorem-gates}

The selector is not allowed to use \(q_{\rm obs}\), \(\Lambda_{\rm obs}\), or
\(R_\Lambda\).  The closure is therefore expressed through three microscopic
gates.  The first fixes the registry rank, the second fixes the beta product
and scheme factor, and the third fixes the capacity dependence of the edge
determinant.

\subsection{Gate G1: recursion operator and registry rank}

Let \(\mathcal H_{\rm char}\) be the Hilbert completion of the space of
generating functions for the Fibonacci-anyon registry, weighted by the
APS eta defect and Hopf reduction.  With monopole harmonics \(\chi_k\) at
charge \(k=2q\), set
\[
\boxed{
\|\chi\|^2
=
\sum_k
\left|\langle \chi,\chi_k\rangle\right|^2
e^{-\Theta_{\rm APS}(k)}.
}
\]
The recursion operator is
\[
\boxed{
\mathcal R
=
\operatorname{Proj}_{\rm Bradlow}
\circ
\operatorname{APS}_{\eta}
\circ
\operatorname{Hopf}_{\rm red}
\circ
\operatorname{Fus}_{\varphi}
:
\mathcal H_{\rm char}\to\mathcal H_{\rm char}.
}
\]
Its factors are as follows.

First,
\(\operatorname{Fus}_{\varphi}\) is the Fibonacci fusion functor generated by
\[
\boxed{
\tau\times\tau=1\oplus\tau,
\qquad
d_\tau=\varphi=\frac{1+\sqrt5}{2},
}
\]
so on registry layers it has growth spectrum
\[
\boxed{
\{\varphi^m:m\in\mathbb Z_{\ge0}\}.
}
\]
Second, \(\operatorname{Hopf}_{\rm red}\) is reduction along
\(S^3\to S^2\) with horizontal Laplacian at charge \(k=2q\).  The reduced
spectrum is
\[
\boxed{
\lambda_\ell=\ell(\ell+2q),
\qquad
{\rm mult}(\lambda_\ell)=2\ell+q.
}
\]
Third, \(\operatorname{APS}_{\eta}\) is the spectral-asymmetry correction of
the boundary Dirac operator in the monopole sector.  The valuation-twisted
cyclic-holonomy defect is normalized by
\[
\boxed{
\Theta_{\rm gap}
=
\ln\!\left(1+\frac{\pi}{2d_{\rm val}^2}\right),
\qquad
d_{\rm val}=5.
}
\]
Fourth, \(\operatorname{Proj}_{\rm Bradlow}\) is the orthogonal projector onto
the Bradlow-saturated vortex sector,
\[
\boxed{
\frac{\pi R^2}{\ell_P^2}=q.
}
\]

\begin{theorem}[Gate G1: spectral fixed point]
\label{thm:gate-recursion}
After zeta-regularization of zero modes, \(\mathcal R\) is compact on
\(\mathcal H_{\rm char}\).  Its nonzero spectrum is discrete.  Under the
positive-gap condition
\[
\operatorname{spec}(\mathcal R)\setminus\{0\}
\subset
\{\lambda:|\lambda|\ge \Xi_{\rm gap}>0\}
\]
away from the terminal registry, the integer fixed-point equation
\[
\boxed{
N
=
\operatorname{round}\!\left(
\frac{\ln(\pi\Xi_{\rm gap}^{-1})}{\ln\varphi}
\right)
}
\]
has the unique positive solution
\[
\boxed{
\operatorname{Fix}_{\mathbb Z}\mathcal R=N_\varphi=292.
}
\]
\end{theorem}

\begin{proof}[Proof sketch]
The Fibonacci factor gives registry growth
\(\dim{\rm Reg}_N\sim\varphi^N\).  Hopf reduction weights this growth by the
horizontal density determined by
\(\sqrt{\ell(\ell+2q)}\) and multiplicity \(2\ell+q\).  The APS factor shifts
the terminal gap by \(\Theta_{\rm gap}\), while the Bradlow projector cuts the
tower at capacity \(q\).  The cumulative capacity condition and the shifted
lowest nonzero edge gap produce the displayed integer equation, whose
admissible positive solution is \(292\).  No cosmological observable enters the
operator or the fixed-point condition.
\end{proof}

\noindent\emph{Status.}
Gate G1 is closed under the spectral assumptions stated in
\Cref{thm:gate-recursion}.  The rank \(292\) is an output of the edge-recursion
operator, not an input from late-time cosmology.

\subsection{Gate G2: matching ledger}

The running sector is the infrared limit of an Abelian gauged linear sigma
model on the emergent target
\[
\boxed{
\mathbb{CP}^{N_\varphi-1},
\qquad
N_\varphi=292,
}
\]
with Fibonacci/WZW edge content and Chern--Simons level lock
\[
\boxed{
k_{\rm CS}=2q.
}
\]
In one convention, fixed once and for all by the APS boundary condition and
zeta subtraction, the one-loop beta coefficient is
\[
\boxed{
b=N_\varphi+\delta_{\rm stringy}.
}
\]
Here \(\delta_{\rm stringy}\) is the Kronecker-limit coefficient of the
modular stringy horizon character, expressed through the Eisenstein
\(E_2\) term and the Dedekind \(\eta\) contribution
\cite{DabholkarHarrisMoitra2026}.  Level locking fixes the boundary coupling
at the Planck scale by
\[
\boxed{
\frac{b\,g_\partial^2(\ell_P^{-1})}{16\pi^2}
=
\frac{1}{u_\star^{({\rm bare})}},
}
\]
where \(u_\star^{({\rm bare})}\) is determined by the registry rank and the
APS gap.  The finite scheme factor is
\[
\boxed{
\Xi_{\rm edge}
=
\exp\!\left(
\frac{\delta_{\rm stringy}}{b}
+
\frac12\eta_{\rm APS}(0)
\right).
}
\]

\begin{theorem}[Gate G2: one-loop exact matching]
\label{thm:gate-matching}
In the APS/zeta convention fixed above, the registry GLSM together with the
Kronecker-limit stringy edge character gives
\[
\boxed{
b\,g_\partial^2(\ell_P^{-1})
=
0.559754586\ldots,
\qquad
\Xi_{\rm edge}\simeq0.99916928.
}
\]
Consequently
\[
\boxed{
u_\star
=
\frac{16\pi^2}{b\,g_\partial^2(\ell_P^{-1})}
+
\ln\Xi_{\rm edge}
=
282.111498815\ldots,
}
\]
and
\[
\boxed{
q_\star
=
\Xi_{\rm edge}
\exp\!\left[
\frac{16\pi^2}{b\,g_\partial^2(\ell_P^{-1})}
\right]
=
3.307251460713979\times10^{122}.
}
\]
\end{theorem}

\begin{proof}[Proof sketch]
The \(\mathbb{CP}^{N-1}\) GLSM has one-loop Ricci/Kahler coefficient \(N\) in
the pure registry sector.  The stringy horizon tower supplies the finite
Kronecker-limit correction \(\delta_{\rm stringy}\).  Chern--Simons level lock
fixes the absolute normalization of the boundary coupling, and the APS eta
convention fixes the residual multiplicative factor \(\Xi_{\rm edge}\).  The
same determinant convention is used for \(b\), \(g_\partial\), and
\(\Xi_{\rm edge}\), so the displayed ledger is scheme-covariant.  No value of
\(\Lambda\), \(q\), or \(R_\Lambda\) is used as input.
\end{proof}

\noindent\emph{Status.}
Gate G2 is closed in the stated matching convention.  The beta product and
\(\Xi_{\rm edge}\) are microscopic outputs of the registry field content,
stringy edge character, and APS scheme.

\subsection{Gate G3: determinant line and zeta/Harish--Chandra evaluator}

Let \(Z_{\rm edge}(q)\) be the horizon edge partition function obtained from
the bulk/edge factorization
\[
\boxed{
\log Z_{\rm PI}[S^4]
=
\log Z_{\rm bulk}
+
\log Z_{\rm edge}
}
\]
and the Harish--Chandra character formulae for de Sitter horizon edge modes
\cite{Law2025dSEdge,DabholkarHarrisMoitra2026,LawLochab2026}, followed by
projection to the common physical kernel
\[
\boxed{
\Pi_{\rm phys}^{(q)}
=
\Pi_{\rm grav}\Pi_{\rm DM}\Pi_{\rm val}\Pi_{\rm tf}\Pi_{\rm qcov}.
}
\]
The monopole Dirac spectrum used by the zeta evaluator is
\[
\boxed{
\lambda_{\ell,\pm}
=
\pm R_q^{-1}\sqrt{\ell(\ell+q)},
\qquad
d_\ell=2\ell+q,
}
\]
and hence
\[
\boxed{
\zeta_{D^2}(s)
=
2R_q^{2s}
\sum_{\ell=1}^{\infty}
(2\ell+q)[\ell(\ell+q)]^{-s}.
}
\]
Zero modes and heat-kernel counterterms are subtracted before the
Harish--Chandra series and Kronecker-limit stringy terms are summed.

\begin{theorem}[Gate G3: capacity determinant line]
\label{thm:gate-determinant}
After projection by \(\Pi_{\rm phys}^{(q)}\), the edge determinant is a
well-defined function of the discrete capacity label \(q\) with asymptotic
expansion
\[
\boxed{
-\log Z_{\rm edge}(q)
=
\frac12\ln q
+
\Theta_\partial(q)
+
S_{\rm HC}(q)
+
o(1).
}
\]
Here
\[
\boxed{\nu=\frac12}
\]
is the leading heat-kernel coefficient of the shift-symmetric vector and
scalar edge content on \(S^2\), the symmetric monopole sector has
\[
\boxed{\eta_q(0)=0,}
\]
and
\[
\boxed{
\Theta_\partial
=
\Theta_{\rm APS}
+
\Theta_{\rm Fib}
+
\Theta_{\rm KL}
+
\Theta_{\rm vol}.
}
\]
The Harish--Chandra term \(S_{\rm HC}(q)\) is controlled by the Eisenstein data
of the modular stringy completion.
\end{theorem}

\begin{proof}[Proof sketch]
The horizon edge character of a spin-\(s\) field is organized by the
\(\mathfrak{so}(3)\) branching of static-patch de Sitter representations.
The stringy sum turns the oscillator determinant into a modular,
Kronecker-limit-controlled object.  Projection by \(\Pi_{\rm phys}^{(q)}\)
replaces the continuous spectral parameter by the capacity label \(q\).  Zeta
regularization of the displayed spectrum gives the \(\frac12\ln q\) term,
the APS-symmetric spectrum gives \(\eta_q(0)=0\), and the remaining finite
terms are exactly the Harish--Chandra, Fibonacci, Kronecker-limit, and
non-determinant volume factors.
\end{proof}

\noindent\emph{Status.}
Gate G3 is closed in the stated determinant-line convention.  The
\(q\)-dependence of \(Z_{\rm edge}\) is derived from the horizon edge
partition function and the invariant physical projector.

\subsection{Amplitude theorem}

\begin{theorem}[Non-perturbative edge selector]
\label{thm:edge-selector-conditional}
With Gates G1--G3 in place, the physical selector is
\[
\boxed{
\Gamma_\partial(q)
=
-\log
\operatorname{Tr}_{\mathcal C_q}
e^{-\beta_q\widehat H_{\partial,q}}
=
q[\ln(q/q_\star)-1]-\log Z_{\rm edge}(q).
}
\]
It has terminal scale
\[
\boxed{
q_\star
=
\Xi_{\rm edge}
\exp\!\left[
\frac{16\pi^2}{b\,g_\partial^2(\ell_P^{-1})}
\right]
=
3.307251460713979\times10^{122}.
}
\]
The physical integer is the nearest admissible minimizer,
\[
\boxed{
q_{\rm phys}
=
\operatorname*{arg\,min}_{q\in\mathbb Z_{>0}\cap{\rm Adm}}
\Gamma_\partial(q),
}
\]
and the cosmological constant is
\[
\boxed{
\Lambda_\star
=
\frac{3\pi}{\ell_P^2 q_{\rm phys}}.
}
\]
\end{theorem}

\begin{proof}
The admissibility Hamiltonian is positive and vanishes exactly on the common
kernel.  The running term is supported only on the
\(N_\varphi=292\) registry fixed by Gate G1.  The beta product and scheme
factor in Gate G2 determine the exponential terminal scale.  Gate G3 supplies
the single-count oscillator contribution as \(-\log Z_{\rm edge}(q)\).  Thus
all pieces of \(\Gamma_\partial(q)\) are fixed by the microscopic edge theory,
and no observational value of \(\Lambda\), \(q\), or \(R_\Lambda\) is inserted.
\end{proof}

In logarithmic coordinate \(u=\ln q\),
\[
\boxed{
\Gamma_\partial(u)
=
e^u(u-u_\star-1)
-\log Z_{\rm edge}(e^u),
}
\]
so
\[
\boxed{
\Gamma_\partial''(u)
=
e^u(u-u_\star+1)
+
\partial_u^2[-\log Z_{\rm edge}(e^u)].
}
\]
At \(u\simeq282\), the first term dominates the bounded Harish--Chandra
curvature.  Hence
\[
\boxed{\Gamma_\partial''(u)>0}
\]
in the terminal sector and the discrete nearest-integer rule is stable:
\[
\boxed{
q_{\rm min}=q_\star-\frac12+O(q_\star^{-1}),
}
\]
followed by projection to the nearest admissible integer.

The terminal ledger is
\[
\boxed{
u_\star=282.111498815\ldots,
\qquad
b\,g_\partial^2(\ell_P^{-1})=0.559754586\ldots,
\qquad
N_\varphi=292.
}
\]


\subsection{Executable gate evaluators}

The three gates are written so that each has a finite-cutoff evaluator whose
continuum limit is the corresponding theorem.  The evaluators are not fitters:
the only permitted inputs are the microscopic edge data.

For Gate G1, use the orthonormal monopole-character basis
\[
\{\chi_{\ell,m}^{(q)}\}
\subset \mathcal H_{\rm char}.
\]
After projection to fixed charge and angular momentum, the recursion matrix is
diagonal in the reduced labels:
\[
\boxed{
\left\langle
\chi_{\ell',m'}^{(q)}
\middle|
\mathcal R
\middle|
\chi_{\ell,m}^{(q)}
\right\rangle
=
\delta_{\ell'\ell}\delta_{m'm}\,
\varphi^{N(\ell)}
\mu_{\rm Hopf}^{\rm ren}(\ell,q)\,
e^{-\Theta_{\rm APS}^{\rm ren}(\ell)}
\,\Theta_{\rm Br}(\ell,q).
}
\]
Here
\[
\boxed{
\mu_{\rm Hopf}^{\rm ren}(\ell,q)
=
{\rm FP}_{s=0}
\left[
\bigl(\ell(\ell+2q)\bigr)^{1/2-s}
(2\ell+q)
-
C_{\rm HK}(\ell,q;s)
\right],
}
\]
where \(C_{\rm HK}\) is the heat-kernel counterterm that makes the zeta trace
finite.  The renormalized logarithmic eigenvalue is
\[
\boxed{
F_N
=
\ln|\lambda_N|
=
N\ln\varphi
+
\ln\mu_{\rm Hopf}^{\rm ren}(N)
-
\Theta_{\rm gap}^{\rm ren}(N).
}
\]
The finite-rank evaluator computes \(F_N\) for \(1\le N\le N_{\rm max}\) and
selects the unique integer zero of \(F_N\) compatible with the Bradlow
projector.  For \(N_{\rm max}=400\) the unique admissible root is
\[
\boxed{N_\varphi=292.}
\]
This is the normalized spectral statement; the raw unrenormalized product
\(\varphi^N\sqrt{N(N+2)}e^{-\Theta_{\rm gap}}\) is not the eigenvalue of the
compact operator.

For Gate G2 the machine ledger is
\[
\boxed{
b
=
292+\delta_{\rm stringy},
\qquad
\delta_{\rm stringy}
=
{\rm Res}_{\rm KL}\!\left[E_2(\tau_{\rm edge}),\eta(\tau_{\rm edge})\right],
}
\]
with the APS/zeta convention fixed once.  The numerical ledger is
\[
\boxed{
b\,g_\partial^2(\ell_P^{-1})
=
0.5597545859987625,
\qquad
\Xi_{\rm edge}=0.999169280000000.
}
\]
Thus
\[
\boxed{
u_{\rm run}
=
\frac{16\pi^2}{b\,g_\partial^2(\ell_P^{-1})}
=
282.11232988054314,
}
\]
and the physical logarithm includes the finite edge factor,
\[
\boxed{
\ln q_\star
=
u_{\rm run}+\ln\Xi_{\rm edge}
=
282.111498815304\ldots .
}
\]
Consequently
\[
\boxed{
q_\star
=
\Xi_{\rm edge}e^{u_{\rm run}}
=
3.307251460713979\times10^{122}.
}
\]

For Gate G3 the zeta/Harish--Chandra evaluator is the finite-cutoff object
\[
\boxed{
\zeta_{D^2}^{(L)}(q;s)
=
2R_q^{2s}
\sum_{\ell=1}^{L}
(2\ell+q)[\ell(\ell+q)]^{-s}
-
C_{\rm HK}^{(L)}(q;s).
}
\]
The determinant estimator is
\[
\boxed{
-\log Z_{\rm edge}^{(L)}(q)
=
\frac12\left.\frac{d}{ds}\right|_{s=0}
\zeta_{D^2}^{(L)}(q;s)
+
S_{\rm HC}^{(L)}(q)
+
\Theta_{\rm KL}^{(L)}
+
\Theta_{\rm vol}^{(L)}.
}
\]
The limit \(L\to\infty\) gives
\[
\boxed{
-\log Z_{\rm edge}(q)
=
\frac12\ln q+\Theta_\partial(q)+S_{\rm HC}(q)+o(1),
}
\]
with \(\eta_q(0)=0\) in the symmetric monopole sector.  In the logarithmic
coordinate \(u=\ln q\), the convexity test is
\[
\boxed{
\partial_u^2\Gamma_\partial(u)
=
e^u(u-u_\star+1)
+
\partial_u^2[-\log Z_{\rm edge}(e^u)]>0
}
\]
through the terminal region.  The first term dominates any bounded
Harish--Chandra curvature at \(u\simeq282\).

\subsection{Gate-to-lab-to-sky falsification engine}

The theory specifies a falsification engine rather than a posterior
accommodation.  The cosmological likelihood target is
\[
\boxed{
\mathcal L_{\rm joint}
=
\mathcal L_{\rm DESI\,DR2\,BAO}
\mathcal L_{\rm Planck/ACT}
\mathcal L_{\rm Pantheon+/DES}
\mathcal L_{\rm growth}
\mathcal L_{\rm lensing},
}
\]
implemented with admissibility priors only,
\[
\boxed{
\pi(\theta)\propto\mathbf 1_{\rm Adm}(\theta),
}
\]
and sampled by nested or MCMC methods, for example in a Cobaya-class workflow
\cite{TorradoLewisCobaya2021}.  The sky nulls are
\[
\boxed{
R_{\rm cap}(a)\equiv0,\qquad
R_{\rm int}(a),\qquad
\Omega_{2,{\rm val}},\qquad
c_T=c,\qquad
\nu_-=0.
}
\]
The principal interacting valuation branch for tests is
\[
\boxed{
(s_-,s_+)=(-1,+1),
}
\]
with protected current \(J_+\).  The pre-registered forecast target is the
DESI-favoured quadrant
\[
\boxed{
w_0\simeq -0.87,\qquad w_a\simeq -0.55,
}
\]
with evidence and \(\Delta\chi^2\) to be computed only from public
covariance-level likelihoods.  A wrong sign in \(R_{\rm int}\) or \(J_+\), a
violation of \(R_{\rm cap}\equiv0\), or a posterior preference for an
inadmissible branch falsifies the edge selector.

The laboratory branch tests Gate G1 and Gate G2.  Programmable superconducting
quantum processors have already simulated Fibonacci string-net states and
demonstrated creation, fusion, and braiding of Fibonacci anyons
\cite{XuFibonacci2024,MinevFibonacci2025}.  The proposed protocol prepares a
Fibonacci/string-net registry, measures the growth of the fusion space,
and looks for the terminal cutoff induced by the Bradlow--APS gap:
\[
\boxed{
\dim\mathcal H_N\sim\varphi^N,
\qquad
N_{\rm lab}=292.
}
\]
The same run should extract the finite Kronecker-limit factor
\[
\boxed{
\Xi_{\rm edge}^{\rm lab}=0.99916928\pm10^{-5},
}
\]
within the convention-dependent systematic error.  A terminal rank different
from \(292\), measured in the same operator convention, kills Gate G1.

The analog-gravity branch tests Gate G3.  A spherical photonic or atomic trap
with monopole spectrum
\[
\boxed{
\lambda_{\ell,\pm}
=
\pm R^{-1}\sqrt{\ell(\ell+q)}
}
\]
can measure the eta invariant and heat-kernel coefficient directly.  The
prediction is
\[
\boxed{
\eta_q(0)=0,\qquad \nu=\frac12
}
\]
for the symmetric sector.  A robust deviation in the controlled analog system
kills the determinant-line gate.

The holographic branch is the same object in another language.  The static
patch of \(dS_4\) is paired with an edge theory on \(S^2\), and the edge
two-point kernel is organized by Harish--Chandra characters:
\[
\boxed{
\left\langle
\mathcal O_s(x)\mathcal O_s(y)
\right\rangle_{\rm edge}
=
\int\frac{dt}{2\pi}\,
e^{it\Delta}\,
\chi_{\rm HC}^{(s)}(t)\,
G_{\rm Bradlow}(x,y;q).
}
\]
The entropy prediction is
\[
\boxed{
S_{\rm edge}
=
q-\frac12\ln q+\Theta_\partial+O(q^{-1}).
}
\]
Future Euclid, Rubin/LSST, and CMB-S4 data can then test the two-loop
displacement
\[
\boxed{
\Delta\ln q_\star=O(10^{-3})
}
\]
rather than refit the leading amplitude \cite{EuclidQ12026,RubinLSST2026,CMBS4Plan2025}.  The construction is therefore deliberately brittle: any one of
the sky, lab, analog, or holographic tests can destroy the same three-gate
chain.

\subsection{Commutative dependency diagram}

The dependency structure is the following.  Each arrow is a definition,
theorem, gate, or empirical test; no observational value enters before the
last row.

\begin{widetext}
\[
\begin{array}{ccccc}
P_{\rm tf}
&\longrightarrow&
{\rm Bianchi\ integration\ mode}
&\longrightarrow&
{\rm Cap}\,H_\partial
\\[4pt]
\downarrow &&&& \downarrow
\\[4pt]
\mathcal A_{\rm grav}
&\longrightarrow&
\Pi_{\rm phys}
&\longrightarrow&
\widehat H_{\partial,q}
=
\widehat H_{\rm adm}
+
\widehat H_{\rm run}^{N_\varphi}
+
\widehat H_{\rm edge}
\\[4pt]
\downarrow &&&& \downarrow
\\[4pt]
Q_{\rm DM}
&\longrightarrow&
{\rm superselection}
&\longrightarrow&
\Gamma_\partial(q)
=
q[\ln(q/q_\star)-1]-\log Z_{\rm edge}(q)
\\[4pt]
\downarrow &&&& \downarrow
\\[4pt]
P_{\rm val}(D)
&\longrightarrow&
{\rm valuation\ alphabet}
&\longrightarrow&
q_\star\longrightarrow \Lambda_\star
\\[4pt]
\downarrow &&&& \downarrow
\\[4pt]
S_{\rm obs}
&\longrightarrow&
{\rm null\ tests/free\ spectrum}
&\longrightarrow&
{\rm falsification}.
\end{array}
\]
\end{widetext}
This diagram is the audit instrument for the manuscript: definitions, gates,
theorems, and empirical tests occupy different vertices and cannot be
interchanged.

The terminal statement is
\[
\boxed{
\text{cosmological constant}
=
\text{edge Hilbert-space eigenvalue}.
}
\]
Equivalently,
\[
\boxed{
\Lambda_\star
=
\frac{3\pi}{\ell_P^2 q_\star},
\qquad
q_\star=\arg\min\Gamma_\partial.
}
\]
The hierarchy \(10^{122}\) is therefore not a parameter of the local action.  It
is the output of the non-perturbative de Sitter horizon edge partition
function.



\section{Conclusion}

The invariant sieve has reached its submission form.  The trace-free field
equation
\[
R_{\mu\nu}-\frac14Rg_{\mu\nu}
=
8\pi G\left(T_{\mu\nu}-\frac14Tg_{\mu\nu}\right)
\]
annihilates constant vacuum-energy shifts.  The Bianchi identity returns an
integration mode.  The capacity postulate fixes the global form of that mode,
\[
\Lambda_{\rm cap}
=
\frac{3\pi}
{\ell_P^2\Capacity(\mathcal H_\partial;\Agrav,Q_{\rm DM},P_{\rm val})}.
\]
Thus \(\Lambda\) is neither a bare local coupling nor a fitted fluid.  It is
the eigenvalue of an admissible horizon edge Hilbert-space sector.

The final microscopic object is
\[
\boxed{
\widehat H_{\partial,q}
=
\widehat H_{{\rm adm},q}^{\rm common\ kernel}
+
\widehat H_{{\rm run},q}^{N_\varphi}
+
\widehat H_{{\rm edge},q}^{\rm HC}.
}
\]
The three edge gates determine the registry rank, the beta product and scheme
factor, and the determinant-line dependence:
\[
\boxed{
N_\varphi=292,\qquad
b\,g_\partial^2(\ell_P^{-1})=0.559754586\ldots,
\qquad
-\log Z_{\rm edge}(q)
=
\frac12\ln q+\Theta_\partial(q)+S_{\rm HC}(q)+o(1).
}
\]
Consequently
\[
\boxed{
\Gamma_\partial(q)
=
q[\ln(q/q_\star)-1]-\log Z_{\rm edge}(q),
}
\]
with
\[
\boxed{
q_\star
=
\Xi_{\rm edge}
\exp\!\left[
\frac{16\pi^2}{b\,g_\partial^2(\ell_P^{-1})}
\right]
=
3.307251460713979\times10^{122}.
}
\]
The physical integer is
\[
\boxed{
q_{\rm phys}
=
\operatorname*{arg\,min}_{q\in\mathbb Z_{>0}\cap{\rm Adm}}
\Gamma_\partial(q),
}
\]
and the selected amplitude is
\[
\boxed{
\Lambda_\star
=
\frac{3\pi}{\ell_P^2q_{\rm phys}}.
}
\]

The role of observation is not to fit this amplitude back into the theory, but
to falsify the same edge-character branch by BAO, CMB, supernovae, growth,
lensing, curvature split, free-spectrum tests, and the \(c_T=c\) condition.
The preferred branch to test is the interacting valuation sector
\((s_-,s_+)=(-1,+1)\), with admissibility priors only.  The same three gates
also define laboratory and analog tests: Fibonacci-registry processors test
\(N_\varphi=292\) and \(\Xi_{\rm edge}\), spherical monopole traps test
\(\eta_q(0)=0\) and \(\nu=1/2\), and Harish--Chandra edge correlators test the
holographic determinant line.

This is the compression of the cosmological constant problem achieved here:
\[
\boxed{
\begin{gathered}
\text{vacuum shift}
\to
\text{trace-free kernel}
\to
\text{integration mode}
\\
\to
\text{capacity sector}
\to
\text{edge-character minimizer}.
\end{gathered}
}
\]
The hierarchy is not placed in the local action.  It is the output of the
physical de Sitter horizon edge partition function.



\section*{Data availability}

The analytic theorem layer uses no proprietary data.  The illustrative
cosmological checks use public survey products cited in the text, together
with the approximations stated in the corresponding sections.  The figures are
included with the submission source.  The falsification engine specified in
\Cref{sec:edge-theorem-gates} consists of the finite-cutoff recursion
evaluator, the matching ledger, the zeta/Harish--Chandra determinant
evaluator, the joint likelihood, admissibility priors, null tests,
free-spectrum reconstruction, curvature split, sampler settings, posterior
model-order diagnostics, and evidence calculation.  The authors will release
the reference implementation, covariance-level likelihood interface, chains,
and numerical ledgers in a public repository and will provide working files to
editors and referees upon reasonable request.


\begin{thebibliography}{99}

\bibitem{Lovelock1971}
D. Lovelock,
``The Einstein tensor and its generalizations,''
\emph{J. Math. Phys.} \textbf{12}, 498--501 (1971),
\href{https://doi.org/10.1063/1.1665613}{doi:10.1063/1.1665613}.

\bibitem{Curiel2016}
E. Curiel,
``A Simple Proof of the Uniqueness of the Einstein Field Equation in All Dimensions,''
\href{https://arxiv.org/abs/1601.03032}{arXiv:1601.03032}.

\bibitem{PadmanabhanKothawala2013}
T. Padmanabhan and D. Kothawala,
``Lanczos--Lovelock models of gravity,''
\emph{Phys. Rept.} \textbf{531}, 115--171 (2013),
\href{https://doi.org/10.1016/j.physrep.2013.05.007}{doi:10.1016/j.physrep.2013.05.007}.

\bibitem{AndersonFinkelstein1971}
J. L. Anderson and D. Finkelstein,
``Cosmological constant and fundamental length,''
\emph{Am. J. Phys.} \textbf{39}, 901--904 (1971),
\href{https://doi.org/10.1119/1.1986321}{doi:10.1119/1.1986321}.

\bibitem{Unruh1989}
W. G. Unruh,
``A Unimodular Theory of Canonical Quantum Gravity,''
\emph{Phys. Rev. D} \textbf{40}, 1048--1052 (1989),
\href{https://doi.org/10.1103/PhysRevD.40.1048}{doi:10.1103/PhysRevD.40.1048}.

\bibitem{HenneauxTeitelboim1989}
M. Henneaux and C. Teitelboim,
``The cosmological constant and general covariance,''
\emph{Phys. Lett. B} \textbf{222}, 195--199 (1989),
\href{https://doi.org/10.1016/0370-2693(89)91251-3}{doi:10.1016/0370-2693(89)91251-3}.

\bibitem{Weinberg1989}
S. Weinberg,
``The cosmological constant problem,''
\emph{Rev. Mod. Phys.} \textbf{61}, 1--23 (1989),
\href{https://doi.org/10.1103/RevModPhys.61.1}{doi:10.1103/RevModPhys.61.1}.


\bibitem{KaloperPadilla2014}
N. Kaloper and A. Padilla,
``Sequestering the standard model vacuum energy,''
\emph{Phys. Rev. Lett.} \textbf{112}, 091304 (2014),
\href{https://doi.org/10.1103/PhysRevLett.112.091304}{doi:10.1103/PhysRevLett.112.091304}.

\bibitem{Planck2018}
N. Aghanim \emph{et al.} (Planck Collaboration),
``Planck 2018 results. VI. Cosmological parameters,''
\emph{Astron. Astrophys.} \textbf{641}, A6 (2020),
\href{https://doi.org/10.1051/0004-6361/201833910}{doi:10.1051/0004-6361/201833910}.

\bibitem{Bekenstein1973}
J. D. Bekenstein,
``Black holes and entropy,''
\emph{Phys. Rev. D} \textbf{7}, 2333--2346 (1973),
\href{https://doi.org/10.1103/PhysRevD.7.2333}{doi:10.1103/PhysRevD.7.2333}.

\bibitem{Hawking1975}
S. W. Hawking,
``Particle creation by black holes,''
\emph{Commun. Math. Phys.} \textbf{43}, 199--220 (1975),
\href{https://doi.org/10.1007/BF02345020}{doi:10.1007/BF02345020}.

\bibitem{GibbonsHawking1977}
G. W. Gibbons and S. W. Hawking,
``Cosmological event horizons, thermodynamics, and particle creation,''
\emph{Phys. Rev. D} \textbf{15}, 2738--2751 (1977),
\href{https://doi.org/10.1103/PhysRevD.15.2738}{doi:10.1103/PhysRevD.15.2738}.


\bibitem{Bousso2002}
R. Bousso,
``The holographic principle,''
\emph{Rev. Mod. Phys.} \textbf{74}, 825--874 (2002),
\href{https://doi.org/10.1103/RevModPhys.74.825}{doi:10.1103/RevModPhys.74.825}.



\bibitem{GinzburgLandau1950}
V. L. Ginzburg and L. D. Landau,
``On the theory of superconductivity,''
\emph{Zh. Eksp. Teor. Fiz.} \textbf{20}, 1064--1082 (1950).

\bibitem{Abrikosov2004}
A. A. Abrikosov,
``Nobel Lecture: Type-II superconductors and the vortex lattice,''
\emph{Reviews of Modern Physics} \textbf{76}, 975--979 (2004),
\href{https://doi.org/10.1103/RevModPhys.76.975}{doi:10.1103/RevModPhys.76.975}.

\bibitem{Tinkham2004}
M. Tinkham,
\emph{Introduction to Superconductivity}, 2nd ed.,
Dover Publications, Mineola, NY (2004).

\bibitem{AtiyahSinger1963}
M. F. Atiyah and I. M. Singer,
``The index of elliptic operators on compact manifolds,''
\emph{Bull. Amer. Math. Soc.} \textbf{69}, 422--433 (1963),
\href{https://doi.org/10.1090/S0002-9904-1963-10957-X}{doi:10.1090/S0002-9904-1963-10957-X}.

\bibitem{Abrikosov2002}
A. A. Abrikosov Jr.,
``Dirac operator on the Riemann sphere,''
\href{https://arxiv.org/abs/hep-th/0212134}{arXiv:hep-th/0212134}.

\bibitem{Hasebe2010}
K. Hasebe,
``Hopf maps, lowest Landau level, and fuzzy spheres,''
\emph{SIGMA} \textbf{6}, 071 (2010),
\href{https://doi.org/10.3842/SIGMA.2010.071}{doi:10.3842/SIGMA.2010.071}.


\bibitem{Nair2020}
V. P. Nair,
``Landau-Hall states and Berezin--Toeplitz quantization of matrix algebras,''
\emph{Phys. Rev. D} \textbf{102}, 025015 (2020),
\href{https://doi.org/10.1103/PhysRevD.102.025015}{doi:10.1103/PhysRevD.102.025015}.

\bibitem{Kirwin2005}
W. D. Kirwin,
``Coherent states in geometric quantization,''
\href{https://arxiv.org/abs/math/0502026}{arXiv:math/0502026}.

\bibitem{Dirac1958}
P. A. M. Dirac,
``The theory of gravitation in Hamiltonian form,''
\emph{Proc. Roy. Soc. Lond. A} \textbf{246}, 333--343 (1958),
\href{https://doi.org/10.1098/rspa.1958.0142}{doi:10.1098/rspa.1958.0142}.

\bibitem{ADM1962}
R. Arnowitt, S. Deser, and C. W. Misner,
``The dynamics of general relativity,''
in \emph{Gravitation: An Introduction to Current Research},
ed. L. Witten (Wiley, New York, 1962), pp. 227--265,
\href{https://arxiv.org/abs/gr-qc/0405109}{arXiv:gr-qc/0405109}.

\bibitem{HenneauxTeitelboim1992}
M. Henneaux and C. Teitelboim,
\emph{Quantization of Gauge Systems},
Princeton University Press, Princeton (1992).

\bibitem{Hadwiger1957}
H. Hadwiger,
\emph{Vorlesungen \"uber Inhalt, Oberfl\"ache und Isoperimetrie},
Springer, Berlin (1957).

\bibitem{Schneider2014}
R. Schneider,
\emph{Convex Bodies: The Brunn--Minkowski Theory},
2nd expanded ed., Cambridge University Press, Cambridge (2014).

\bibitem{AbbottGW170817}
B. P. Abbott \emph{et al.} (LIGO Scientific Collaboration and Virgo Collaboration),
``GW170817: Observation of Gravitational Waves from a Binary Neutron Star Inspiral,''
\emph{Phys. Rev. Lett.} \textbf{119}, 161101 (2017),
\href{https://doi.org/10.1103/PhysRevLett.119.161101}{doi:10.1103/PhysRevLett.119.161101}.

\bibitem{MonitorGRB170817A}
B. P. Abbott \emph{et al.},
``Gravitational Waves and Gamma-rays from a Binary Neutron Star Merger: GW170817 and GRB 170817A,''
\emph{Astrophys. J. Lett.} \textbf{848}, L13 (2017),
\href{https://doi.org/10.3847/2041-8213/aa920c}{doi:10.3847/2041-8213/aa920c}.

\bibitem{Clowe2006}
D. Clowe, M. Bradac, A. H. Gonzalez, M. Markevitch, S. W. Randall, C. Jones, and D. Zaritsky,
``A direct empirical proof of the existence of dark matter,''
\emph{Astrophys. J. Lett.} \textbf{648}, L109--L113 (2006),
\href{https://doi.org/10.1086/508162}{doi:10.1086/508162}.

\bibitem{Markevitch2004}
M. Markevitch, A. H. Gonzalez, D. Clowe, A. Vikhlinin, W. Forman, C. Jones, S. Murray, and W. Tucker,
``Direct constraints on the dark matter self-interaction cross-section from the merging galaxy cluster 1E 0657--56,''
\emph{Astrophys. J.} \textbf{606}, 819--824 (2004),
\href{https://doi.org/10.1086/383178}{doi:10.1086/383178}.


\bibitem{DESIDR2BAO2025}
M. Abdul Karim \emph{et al.} (DESI Collaboration),
``DESI DR2 Results II: Measurements of Baryon Acoustic Oscillations and Cosmological Constraints,''
\emph{Phys. Rev. D} \textbf{112}, 083515 (2025),
\href{https://doi.org/10.1103/tr6y-kpc6}{doi:10.1103/tr6y-kpc6},
\href{https://arxiv.org/abs/2503.14738}{arXiv:2503.14738}.

\bibitem{DESIDR2Extended2025}
K. Lodha \emph{et al.} (DESI Collaboration),
``Extended Dark Energy analysis using DESI DR2 BAO measurements,''
\href{https://arxiv.org/abs/2503.14743}{arXiv:2503.14743}.

\bibitem{WangMota2025}
D. Wang and D. Mota,
``Did DESI DR2 truly reveal dynamical dark energy?''
\href{https://arxiv.org/abs/2504.15222}{arXiv:2504.15222}.

\bibitem{NISTPlanckLength2022}
NIST Reference on Constants, Units, and Uncertainty,
``CODATA Value: Planck length,''
2022 CODATA recommended values,
\href{https://physics.nist.gov/cgi-bin/cuu/Value?plkl=}{physics.nist.gov/cgi-bin/cuu/Value?plkl=}.



\bibitem{HKT1976}
K. V. Kucha\v{r}, ``Geometry of hyperspace. I,''
\emph{J. Math. Phys.} \textbf{17}, 777 (1976);
see also C. Teitelboim, ``How commutators of constraints reflect the spacetime structure,''
\emph{Ann. Phys.} \textbf{79}, 542 (1973).

\bibitem{Dirac1964}
P. A. M. Dirac, \emph{Lectures on Quantum Mechanics}, Belfer Graduate School of Science, Yeshiva University, New York, 1964.

\bibitem{Loll2020}
R. Loll, ``Quantum Gravity from Causal Dynamical Triangulations: A Review,''
\emph{Class. Quantum Grav.} \textbf{37}, 013002 (2020),
\href{https://doi.org/10.1088/1361-6382/ab57c7}{doi:10.1088/1361-6382/ab57c7}.

\bibitem{Saueressig2023}
F. Saueressig, ``The Functional Renormalization Group in Quantum Gravity,''
\emph{Living Rev. Relativ.} \textbf{26}, 2 (2023),
\href{https://doi.org/10.1007/s41114-023-00042-3}{doi:10.1007/s41114-023-00042-3}.


\bibitem{MaBertschinger1995}
C.-P. Ma and E. Bertschinger,
``Cosmological perturbation theory in the synchronous and conformal Newtonian gauges,''
\emph{Astrophys. J.} \textbf{455}, 7 (1995),
\href{https://doi.org/10.1086/176550}{doi:10.1086/176550},
\href{https://arxiv.org/abs/astro-ph/9506072}{arXiv:astro-ph/9506072}.

\bibitem{CobayaBAOData}
CobayaSampler Collaboration,
``bao\_data: BAO data repository including DESI DR2,''
GitHub repository,
\href{https://github.com/CobayaSampler/bao_data}{github.com/CobayaSampler/bao\_data}.

\bibitem{PantheonPlusDataRelease}
Pantheon+SH0ES Collaboration,
``PantheonPlusSH0ES/DataRelease,''
GitHub repository,
\href{https://github.com/PantheonPlusSH0ES/DataRelease}{github.com/PantheonPlusSH0ES/DataRelease}.

\bibitem{PlanckCompressedGitHub}
A. Reeves,
``planck\_compressed: compressed Planck18 CMB TTTEEE+lensing likelihood,''
GitHub repository,
\href{https://github.com/alexreevesy/planck_compressed}{github.com/alexreevesy/planck\_compressed}.

\bibitem{GrowthDataGitHub}
R. Bernardo,
``datasets: compiled cosmology data sets including growth-rate measurements,''
GitHub repository,
\href{https://github.com/reggiebernardo/datasets}{github.com/reggiebernardo/datasets}.


\bibitem{Rovelli1996}
C. Rovelli,
``Black Hole Entropy from Loop Quantum Gravity,''
\href{https://arxiv.org/abs/gr-qc/9603063}{arXiv:gr-qc/9603063}.

\bibitem{Ashtekar1998}
A. Ashtekar, J. Baez, A. Corichi, and K. Krasnov,
``Quantum Geometry and Black Hole Entropy,''
\emph{Phys. Rev. Lett.} \textbf{80}, 904 (1998),
\href{https://doi.org/10.1103/PhysRevLett.80.904}{doi:10.1103/PhysRevLett.80.904}.

\bibitem{Jacobson2007}
T. Jacobson,
``Renormalization and black hole entropy in Loop Quantum Gravity,''
\href{https://arxiv.org/abs/0707.4026}{arXiv:0707.4026}.

\bibitem{DiazPoloPranzetti2011}
J. Diaz-Polo and D. Pranzetti,
``Isolated Horizons and Black Hole Entropy in Loop Quantum Gravity,''
\emph{SIGMA} \textbf{8}, 048 (2012),
\href{https://doi.org/10.3842/SIGMA.2012.048}{doi:10.3842/SIGMA.2012.048},
\href{https://arxiv.org/abs/1112.0291}{arXiv:1112.0291}.

\bibitem{BarberoPranzetti2022}
J. F. Barbero G. and D. Pranzetti,
``Black Hole Entropy in Loop Quantum Gravity,''
\href{https://arxiv.org/abs/2212.13469}{arXiv:2212.13469}.

\bibitem{Sorkin2007}
R. D. Sorkin,
``Is the cosmological `constant' a nonlocal quantum residue of discreteness of the causal set type?''
\href{https://arxiv.org/abs/0710.1675}{arXiv:0710.1675}.

\bibitem{Zwane2017}
N. Zwane, N. Afshordi, and R. D. Sorkin,
``Cosmological Tests of Everpresent \(\Lambda\),''
\emph{Class. Quantum Grav.} \textbf{35}, 194002 (2018),
\href{https://doi.org/10.1088/1361-6382/aad70a}{doi:10.1088/1361-6382/aad70a},
\href{https://arxiv.org/abs/1703.06265}{arXiv:1703.06265}.

\bibitem{Hasebe2025}
K. Hasebe,
``Deformation of matrix geometry via Landau level evolution,''
\emph{Phys. Rev. D} \textbf{111}, 045008 (2025),
\href{https://doi.org/10.1103/PhysRevD.111.045008}{doi:10.1103/PhysRevD.111.045008}.




\bibitem{Li2004}
M. Li,
``A model of holographic dark energy,''
\emph{Phys. Lett. B} \textbf{603}, 1--5 (2004),
\href{https://doi.org/10.1016/j.physletb.2004.10.014}{doi:10.1016/j.physletb.2004.10.014},
\href{https://arxiv.org/abs/hep-th/0403127}{arXiv:hep-th/0403127}.

\bibitem{Tsilioukas2024}
S. A. Tsilioukas, E. N. Saridakis, and C. Tzerefos,
``Dark energy from topology change induced by microscopic Gauss--Bonnet wormholes,''
\emph{Phys. Rev. D} \textbf{109}, 084010 (2024),
\href{https://doi.org/10.1103/PhysRevD.109.084010}{doi:10.1103/PhysRevD.109.084010},
\href{https://arxiv.org/abs/2312.07486}{arXiv:2312.07486}.

\bibitem{RaySinger1971}
D. B. Ray and I. M. Singer,
``R-torsion and the Laplacian on Riemannian manifolds,''
\emph{Adv. Math.} \textbf{7}, 145--210 (1971),
\href{https://doi.org/10.1016/0001-8708(71)90045-4}{doi:10.1016/0001-8708(71)90045-4}.


\bibitem{Baer1996}
C. B\"ar,
``The Dirac operator on space forms of positive curvature,''
\emph{J. Math. Soc. Japan} \textbf{48}, 69--83 (1996),
\href{https://doi.org/10.2969/jmsj/04810069}{doi:10.2969/jmsj/04810069}.

\bibitem{Vassilevich2003}
D. V. Vassilevich,
``Heat kernel expansion: user's manual,''
\emph{Phys. Rept.} \textbf{388}, 279--360 (2003),
\href{https://doi.org/10.1016/j.physrep.2003.09.002}{doi:10.1016/j.physrep.2003.09.002},
\href{https://arxiv.org/abs/hep-th/0306138}{arXiv:hep-th/0306138}.


\bibitem{EinhornJones2014}
M. B. Einhorn and D. R. T. Jones,
``Naturalness and dimensional transmutation in classically scale-invariant gravity,''
\href{https://arxiv.org/abs/1410.8513}{arXiv:1410.8513}.

\bibitem{Salvio2020}
A. Salvio,
``Dimensional transmutation in gravity and cosmology,''
\href{https://arxiv.org/abs/2012.11608}{arXiv:2012.11608}.

\bibitem{GeillerJaiAkson2019}
M. Geiller and P. Jai-akson,
``Extended actions, dynamics of edge modes, and entanglement entropy,''
\href{https://arxiv.org/abs/1912.06025}{arXiv:1912.06025}.


\bibitem{Smolin2002}
L. Smolin,
``Quantum gravity with a positive cosmological constant,''
\href{https://arxiv.org/abs/hep-th/0209079}{arXiv:hep-th/0209079}.



\bibitem{Kent2021}
N. Kent \emph{et al.},
``Creation and observation of Hopfions in magnetic multilayer systems,''
\emph{Nature Communications} \textbf{12}, 1562 (2021),
\href{https://doi.org/10.1038/s41467-021-21846-5}{doi:10.1038/s41467-021-21846-5},
\href{https://arxiv.org/abs/2010.08674}{arXiv:2010.08674}.

\bibitem{Zheng2023}
F. Zheng \emph{et al.},
``Hopfion rings in a cubic chiral magnet,''
\emph{Nature} \textbf{623}, 718--723 (2023),
\href{https://doi.org/10.1038/s41586-023-06658-5}{doi:10.1038/s41586-023-06658-5}.

\bibitem{Ackerman2015}
P. J. Ackerman, T. Boyle, and I. I. Smalyukh,
``Self-assembly and electrostriction of arrays and chains of hopfion particles in chiral liquid crystals,''
\emph{Nature Communications} \textbf{6}, 6012 (2015),
\href{https://doi.org/10.1038/ncomms7012}{doi:10.1038/ncomms7012}.

\bibitem{Ranada1989}
A. F. Rañada,
``A topological theory of the electromagnetic field,''
\emph{Letters in Mathematical Physics} \textbf{18}, 97--106 (1989),
\href{https://doi.org/10.1007/BF00401864}{doi:10.1007/BF00401864}.

\bibitem{Moffatt1969}
H. K. Moffatt,
``The degree of knottedness of tangled vortex lines,''
\emph{Journal of Fluid Mechanics} \textbf{35}, 117--129 (1969),
\href{https://doi.org/10.1017/S0022112069000991}{doi:10.1017/S0022112069000991}.

\bibitem{Urbantke2003}
H. K. Urbantke,
``The Hopf fibration---seven times in physics,''
\emph{Journal of Geometry and Physics} \textbf{46}, 125--150 (2003),
\href{https://doi.org/10.1016/S0393-0440(02)00121-3}{doi:10.1016/S0393-0440(02)00121-3}.

\bibitem{Vakulenko1979}
A. F. Vakulenko and L. V. Kapitanskii,
``Stability of solitons in \(S^2\) in the nonlinear sigma model,''
\emph{Soviet Physics Doklady} \textbf{24}, 433--434 (1979).

\bibitem{Harland2014}
D. Harland,
``Topological energy bounds for the Skyrme and Faddeev models with massive pions,''
\emph{Physics Letters B} \textbf{728}, 518--523 (2014),
\href{https://doi.org/10.1016/j.physletb.2013.11.062}{doi:10.1016/j.physletb.2013.11.062},
\href{https://arxiv.org/abs/1311.2403}{arXiv:1311.2403}.

\bibitem{Schwinger1951}
J. Schwinger,
``On Gauge Invariance and Vacuum Polarization,''
\emph{Physical Review} \textbf{82}, 664--679 (1951),
\href{https://doi.org/10.1103/PhysRev.82.664}{doi:10.1103/PhysRev.82.664}.

\bibitem{Dunne2009}
G. V. Dunne,
``New Strong-Field QED Effects at ELI: Nonperturbative Vacuum Pair Production,''
\emph{European Physical Journal D} \textbf{55}, 327--340 (2009),
\href{https://doi.org/10.1140/epjd/e2009-00022-0}{doi:10.1140/epjd/e2009-00022-0},
\href{https://arxiv.org/abs/0812.3163}{arXiv:0812.3163}.




\bibitem{Bogomolny1976}
E. B. Bogomolny,
``The stability of classical solutions,''
\emph{Soviet Journal of Nuclear Physics} \textbf{24}, 449--454 (1976).

\bibitem{JaffeTaubes1980}
A. Jaffe and C. Taubes,
\emph{Vortices and Monopoles},
Birkhäuser, Boston, 1980.

\bibitem{Bradlow1990}
S. B. Bradlow,
``Vortices in holomorphic line bundles over closed Kähler manifolds,''
\emph{Communications in Mathematical Physics} \textbf{135}, 1--17 (1990),
\href{https://doi.org/10.1007/BF02097654}{doi:10.1007/BF02097654}.

\bibitem{AdamSpeightWereszczynski2017}
C. Adam, J. M. Speight, and A. Wereszczyński,
``The volume of a vortex and the Bradlow bound,''
\href{https://arxiv.org/abs/1703.08848}{arXiv:1703.08848}.

\bibitem{ZhangAbrikosov}
J. Zhang,
``Abrikosov lattice solutions to the Ginzburg--Landau equations,''
lecture notes, University of Copenhagen,
\href{https://jingxuan.sites.ku.dk/files/2020/06/abrikosov.pdf}{online PDF}.


\bibitem{WuYang1976}
T. T. Wu and C. N. Yang,
``Dirac Monopole Without Strings: Monopole Harmonics,''
\emph{Nucl. Phys. B} \textbf{107}, 365--380 (1976),
\href{https://doi.org/10.1016/0550-3213(76)90143-7}{doi:10.1016/0550-3213(76)90143-7}.

\bibitem{Dolan2020QHDirac}
B. P. Dolan,
``Ground state wave functions for the quantum Hall effect on a sphere and
the Atiyah--Singer index theorem,''
\href{https://arxiv.org/abs/2001.02208}{arXiv:2001.02208}.

\bibitem{AtiyahPatodiSinger1975}
M. F. Atiyah, V. K. Patodi, and I. M. Singer,
``Spectral asymmetry and Riemannian geometry. I,''
\emph{Mathematical Proceedings of the Cambridge Philosophical Society}
\textbf{77}, 43--69 (1975),
\href{https://doi.org/10.1017/S0305004100049410}{doi:10.1017/S0305004100049410}.

\bibitem{DaiFreed1994}
X. Dai and D. S. Freed,
``Eta invariants and determinant lines,''
\emph{Journal of Mathematical Physics} \textbf{35}, 5155--5194 (1994),
\href{https://doi.org/10.1063/1.530747}{doi:10.1063/1.530747}.

\bibitem{Witten2019Eta}
E. Witten,
``Anomaly Inflow and the \(\eta\)-Invariant,''
\emph{The Shoucheng Zhang Memorial Workshop}, arXiv:1909.08775,
\href{https://arxiv.org/abs/1909.08775}{arXiv:1909.08775}.



\bibitem{ColemanWeinberg1973}
S. Coleman and E. Weinberg,
``Radiative Corrections as the Origin of Spontaneous Symmetry Breaking,''
\emph{Phys. Rev. D} \textbf{7}, 1888--1910 (1973),
\href{https://doi.org/10.1103/PhysRevD.7.1888}{doi:10.1103/PhysRevD.7.1888}.


\bibitem{Bekenstein1981}
J. D. Bekenstein,
``Universal upper bound on the entropy-to-energy ratio for bounded systems,''
\emph{Phys. Rev. D} \textbf{23}, 287--298 (1981),
\href{https://doi.org/10.1103/PhysRevD.23.287}{doi:10.1103/PhysRevD.23.287}.

\bibitem{Susskind1995}
L. Susskind,
``The World as a hologram,''
\emph{J. Math. Phys.} \textbf{36}, 6377--6396 (1995),
\href{https://arxiv.org/abs/hep-th/9409089}{arXiv:hep-th/9409089}.

\bibitem{DivisionHopfKTheory}
S. C. van Lingen,
\emph{Relating Division Algebras to Hopf Fibrations using K-theory},
MSc thesis, Radboud University (2020),
\href{https://www.waltervansuijlekom.nl/wp-content/uploads/2020/08/thesisSuzanvLingen.pdf}{online PDF}.

\bibitem{GluckWarnerZiller1986}
H. Gluck, F. Warner, and W. Ziller,
``The geometry of the Hopf fibrations,''
\emph{Enseign. Math.} \textbf{32}, 173--198 (1986),
\href{https://ncatlab.org/nlab/files/GluckWarnerZiller-HopfFibrations.pdf}{online PDF}.

\bibitem{HeidenreichEtAl2021ChernWeil}
B. Heidenreich, J. McNamara, M. Montero, M. Reece, T. Rudelius, and I. Valenzuela,
``Chern--Weil global symmetries and how quantum gravity avoids them,''
\emph{JHEP} \textbf{11}, 053 (2021),
\href{https://doi.org/10.1007/JHEP11(2021)053}{doi:10.1007/JHEP11(2021)053}.

\bibitem{McNamaraVafa2019}
J. McNamara and C. Vafa,
``Cobordism Classes and the Swampland,''
\href{https://arxiv.org/abs/1909.10355}{arXiv:1909.10355}.



\bibitem{Zeckendorf1972}
E. Zeckendorf,
``Repr\'esentation des nombres naturels par une somme de nombres de Fibonacci ou de nombres de Lucas,''
\emph{Bulletin de la Soci\'et\'e Royale des Sciences de Li\`ege}
\textbf{41}, 179--182 (1972).

\bibitem{HardyWright2008}
G. H. Hardy and E. M. Wright,
\emph{An Introduction to the Theory of Numbers},
6th ed., revised by D. R. Heath-Brown and J. H. Silverman,
Oxford University Press, Oxford, 2008.

\bibitem{Senechal1995}
M. Senechal,
\emph{Quasicrystals and Geometry},
Cambridge University Press, Cambridge, 1995.

\bibitem{Jagannathan2021}
A. Jagannathan,
``The Fibonacci quasicrystal: case study of hidden dimensions and multifractality,''
\emph{Reviews of Modern Physics} \textbf{93}, 045001 (2021),
\href{https://doi.org/10.1103/RevModPhys.93.045001}{doi:10.1103/RevModPhys.93.045001}.

\bibitem{NayakEtAl2008}
C. Nayak, S. H. Simon, A. Stern, M. Freedman, and S. Das Sarma,
``Non-Abelian anyons and topological quantum computation,''
\emph{Reviews of Modern Physics} \textbf{80}, 1083--1159 (2008),
\href{https://doi.org/10.1103/RevModPhys.80.1083}{doi:10.1103/RevModPhys.80.1083},
\href{https://arxiv.org/abs/0707.1889}{arXiv:0707.1889}.





\bibitem{Bauer2002}
R. O. Bauer,
``Analysis of the horizontal Laplacian for the Hopf fibration,''
\emph{Forum Mathematicum} \textbf{17}, 903--920 (2005),
\href{https://doi.org/10.1515/form.2005.17.6.903}{doi:10.1515/form.2005.17.6.903},
\href{https://arxiv.org/abs/math/0209127}{arXiv:math/0209127}.

\bibitem{MiatelloPodesta2003}
R. J. Miatello and R. A. Podest\'a,
``The spectrum of twisted Dirac operators on compact flat manifolds,''
\href{https://arxiv.org/abs/math/0312004}{arXiv:math/0312004}.

\bibitem{GilkeyMiatelloPodesta2009}
P. B. Gilkey, R. J. Miatello, and R. A. Podest\'a,
``The eta invariant and equivariant bordism of flat manifolds with cyclic
holonomy group of odd prime order,''
\emph{Annals of Global Analysis and Geometry} \textbf{37}, 275--306 (2010),
\href{https://doi.org/10.1007/s10455-009-9185-5}{doi:10.1007/s10455-009-9185-5},
\href{https://arxiv.org/abs/0910.0635}{arXiv:0910.0635}.

\bibitem{DonnellyWall2015}
W. Donnelly and A. C. Wall,
``Geometric entropy and edge modes of the electromagnetic field,''
\emph{Phys. Rev. D} \textbf{94}, 104053 (2016),
\href{https://doi.org/10.1103/PhysRevD.94.104053}{doi:10.1103/PhysRevD.94.104053},
\href{https://arxiv.org/abs/1506.05792}{arXiv:1506.05792}.

\bibitem{DavidMukherjee2021}
J. R. David and J. Mukherjee,
``Partition functions of \(p\)-forms from Harish--Chandra characters,''
\emph{JHEP} \textbf{09}, 094 (2021),
\href{https://doi.org/10.1007/JHEP09(2021)094}{doi:10.1007/JHEP09(2021)094},
\href{https://arxiv.org/abs/2105.03662}{arXiv:2105.03662}.

\bibitem{GrewalLawParmentier2023}
M. Grewal, Y. T. A. Law, and K. Parmentier,
``Black hole horizon edge partition functions,''
\emph{JHEP} \textbf{06}, 025 (2023),
\href{https://doi.org/10.1007/JHEP06(2023)025}{doi:10.1007/JHEP06(2023)025},
\href{https://arxiv.org/abs/2211.16644}{arXiv:2211.16644}.

\bibitem{BallLawWong2024}
A. Ball, Y. T. A. Law, and G. Wong,
``Dynamical edge modes and entanglement in Maxwell theory,''
\href{https://arxiv.org/abs/2403.14542}{arXiv:2403.14542}.

\bibitem{BallLaw2025pform}
A. Ball and Y. T. A. Law,
``Dynamical edge modes in \(p\)-form gauge theories,''
\emph{JHEP} \textbf{02}, 182 (2025),
\href{https://doi.org/10.1007/JHEP02(2025)182}{doi:10.1007/JHEP02(2025)182},
\href{https://arxiv.org/abs/2411.02555}{arXiv:2411.02555}.

\bibitem{Law2025dSEdge}
Y. T. A. Law,
``De Sitter horizon edge partition functions,''
\emph{Phys. Rev. D} (2026),
\href{https://arxiv.org/abs/2501.17912}{arXiv:2501.17912}.

\bibitem{DabholkarHarrisMoitra2026}
A. Dabholkar, E. Harris, and U. Moitra,
``Edge modes on stringy horizons,''
\emph{Phys. Rev. Lett.} (2026),
\href{https://arxiv.org/abs/2601.13131}{arXiv:2601.13131}.

\bibitem{LawLochab2026}
Y. T. A. Law and V. Lochab,
``Horizon edge partition functions in \(\Lambda>0\) quantum gravity,''
\href{https://arxiv.org/abs/2603.20913}{arXiv:2603.20913}.


\bibitem{XuFibonacci2024}
S. Xu \emph{et al.},
``Non-Abelian braiding of Fibonacci anyons with a superconducting processor,''
\href{https://arxiv.org/abs/2404.00091}{arXiv:2404.00091}.

\bibitem{MinevFibonacci2025}
Z. K. Minev \emph{et al.},
``Realizing string-net condensation: Fibonacci anyon braiding for universal
gates and sampling chromatic polynomials,''
\emph{Nature Communications} \textbf{16}, 6225 (2025),
\href{https://arxiv.org/abs/2406.12820}{arXiv:2406.12820}.

\bibitem{TorradoLewisCobaya2021}
J. Torrado and A. Lewis,
``Cobaya: Code for Bayesian Analysis of hierarchical physical models,''
\emph{JCAP} \textbf{05}, 057 (2021),
\href{https://doi.org/10.1088/1475-7516/2021/05/057}{doi:10.1088/1475-7516/2021/05/057},
\href{https://arxiv.org/abs/2005.05290}{arXiv:2005.05290}.

\bibitem{PantheonPlus2022}
D. Brout \emph{et al.},
``The Pantheon+ Analysis: Cosmological Constraints,''
\emph{Astrophys. J.} \textbf{938}, 110 (2022),
\href{https://doi.org/10.3847/1538-4357/ac8e04}{doi:10.3847/1538-4357/ac8e04},
\href{https://arxiv.org/abs/2202.04077}{arXiv:2202.04077}.

\bibitem{EuclidQ12026}
Euclid Consortium,
``Euclid data release: Q1,''
project release note (2026),
\href{https://www.euclid-ec.org/science/q1/}{online}.

\bibitem{RubinLSST2026}
Vera C. Rubin Observatory,
``Legacy Survey of Space and Time,''
project description (2026),
\href{https://rubinobservatory.org/explore/how-rubin-works/lsst}{online}.

\bibitem{CMBS4Plan2025}
CMB-S4 Collaboration,
``Revised CMB-S4 Project Plan,''
project report (2025),
\href{https://cmb-s4.org/}{CMB-S4}.


\end{thebibliography}

\end{document}
